\pdfminorversion=7
\documentclass[parskip=half]{scrreprt}
\usepackage[T1]{fontenc}
\usepackage[utf8]{inputenc}
\usepackage[english]{babel}
\usepackage{setspace}
\usepackage{graphicx} % Required for inserting images
\usepackage[force]{feynmp-auto}
\usepackage[backend=biber, uniquename=false, 
    uniquelist=false,
    maxcitenames=3]{biblatex}
\addbibresource{literature.bib}
\usepackage{amsmath,amssymb}
\usepackage{slashed}
\usepackage{autobreak}
\usepackage{graphicx}
\usepackage{siunitx}
\usepackage{bm}
\usepackage[table]{xcolor}
\usepackage{pdfpages}
\usepackage{booktabs}
\usepackage[a-2b,pdf17]{pdfx}
\usepackage{csquotes}
\usepackage[normalem]{ulem}
\usepackage{multirow}
\usepackage{tikz}
\usetikzlibrary{decorations.markings, decorations.pathmorphing, positioning, arrows.meta}
\usepackage{pgfplots}

\pgfplotsset{compat=1.18}
\DeclareOldFontCommand{\bf}{\normalfont\bfseries}{\mathbf}
\DeclareOldFontCommand{\rm}{\normalfont\rmfamily}{\mathrm}
\def\nn{\nonumber\\ }
\def\gamfY#1#2#3{\gamma^{(Y)}_{\substack{#1 \\ #2 #3}}}
\def\rd{{\rm d}}
\def\abs#1{\left| #1 \right| }
\def\lsix{ \mathcal{L}^{(6)}}
\def\lchi{\Lambda_\chi}
\def\lqcd{ \Lambda_{\text{QCD}}}
\def\vev#1{\left\langle #1 \right\rangle}
\def\msbar{$\overline{\hbox{MS}}$}
\def\hyp{\mathsf{y}}
\def\tr{{\rm Tr}\,}
\def\gamHY{\gamma^{(Y)}_H}
\def\gamfY#1#2#3{\gamma^{(Y)}_{\substack{#1 \\ #2 #3}}}
\newcommand{\de}{\text{d}}
\newcommand{\im}{\text{i}}
\newcommand{\op}{\mathcal{O}}
\title{Muon decay: From Michel via LEFT and SMEFT to BSM}
\author{Carl Jakob Moritz}
\date{June 2026}

\begin{document}
%\doublespacing
\begin{fmffile}{main}
\fmfset{arrow_len}{3mm}
\newsavebox{\muondecaySM}
\newsavebox{\muondecayLEFT}


\section*{Abstract}
How would a heavy new field from beyond the Standard Model leave its imprint on a
process like muon decay that occurs well below the electroweak scale? In order to
bridge the energy regimes and investigate the traces of new physics within the
phenomenology of muon decay, we perturb the Michel parameters in the presence of
Low Energy Effective Field Theory (LEFT) operators. To infer their UV origin, we
collect the matching Standard Model Effective Field Theory (SMEFT) operators at
dimension six and seven. Finally, we open them up and ask: Which possible Beyond Standard
Model (BSM) fields could produce them? Focusing on two specific BSM fields -- a scalar doublet and a vector singlet -- we derive the constraints imposed by precision measurements of muon decay. This
approach synthesizes experimental data and effective field theory matching into a
comprehensive exploration of energy regimes and theoretical layers.
\section*{Zusammenfassung}
Wie würde sich die Erweiterung des Standardmodells um schwere neue Felder auf einen Prozess wie den Myonenzerfall auswirken, welcher in einem Energiebereich weit unterhalb der elektroschwachen Skala stattfindet? Um die Energiebereiche miteinander zu verbinden und die möglichen Spuren Neuer Physik in der Phänomenologie des Myonenzerfalls zu untersuchen, führen wir Störungen der Michelparameter durch Low Energy Effective Field Theory (LEFT) Operatoren ein. Um den Ursprung dieser effektiven Operatoren im UV-Bereich zu finden, tragen wir zunächst die dazugehörigen Operatoren der Standard Model Effective Field Theory (SMEFT) der Dimensionen sechs und sieben zusammen. Abschließend lösen wir diese Operatoren in mögliche UV-Freiheitsgrade auf und fragen: Welche Felder jenseits des Standardmodells könnten sie erzeugen? Mit besonderem Fokus auf zwei spezifische neue Felder -- ein skalares Dublett und ein Vektorsingulett -- leiten wir die Massebeschränkungen her, die sich aus Präzisionsmessungen des Myonzerfalls ergeben. Dieser Ansatz verbindet experimentelle Daten mit dem Matching effektiver Feldtheorien und ermöglicht so eine umfassende Untersuchung verschiedener Energiebereiche und theoretischer Beschreibungsebenen.
\section*{Acknowledgments}
I'd like to thank my parents Karin and Matthias who ignited my interest for physics and philosophy and supported my studies for the course of 22 (\enquote{twenty-two}) semesters in Leipzig, Hong Kong and Vienna. I hope this work is part of a movement of thought that truly spans the disciplines and ways of thinking -- like you two do, individually and together. Küsschen :* Moreover I want to thank Lucas, my much-loved partner in crime. I hope any of this makes sense to you. My third thanks goes to the 3 itself. Fourth is time: Thank you, Tyler, for spending so much with me and this project. Howdy!
%My third thanks goes to the 3 itself -- the Shem -- the spirit that is substance. Praised~be.\par 
%Fourth is time: Thank you Tyler for spending so much of it with me and this project. Thank you for showing me the Standard Model. Thank you for you optimism, trust and spirit. Howdy \& God bless America!
\tableofcontents
\chapter{Introduction}
\textit{The true is the whole. However, the whole is only the essence completing itself through its own development.}\par
\begin{raggedleft}
    G.W.F Hegel: Phenomenology of Spirit \cite[p.11]{hegel1977pheno}\par
\end{raggedleft} \vspace{0.8cm}
The muon was not welcomed quite enthusiastically by particle physicists when it made its first appearance in cloud chamber experiments in the 1930s \cite{neddermeyer1937note}: \enquote{Who ordered that?} was Isidor Rabi's famous reaction to the strange new particle \cite{bartusiak1987who}, encapturing the disorientation of a scientific community who firmly believed matter to consist of electrons, neutrons and protons. The muon fitted nowhere in this scheme: Sharing all its properties with the electron except having a mass 200 times higher, the muon seemed to be basically a heavy copy of the electron. The mysterious sibling raised more questions than it answered.

But soon the muon took its place at the very heart of particle research, acting as a messenger for a plethora of unknown physics. Being only short-lived itself and quickly decaying into an electron, it is ecactly the spectrum of latter that provides the main subject for examination. As the produced electron varies both in energy and momentum at each decay instance, it follows that two more particles are being produced which carry away the missing quantities \cite{steinberger1949range}. Those particles were soon identified as neutrinos and in the Brookhaven experiment \cite{danby1962neutrinos} shown to exhibit the same sibling-like relation as the muon and the electron -- thus introducing the fundamental idea of the Standard Model that all matter particles exist as generations. But the electron spectrum becomes even more insightful with the parametrization that Louis Michel introduced \cite{michel1950interaction}, which turned it into a systematic probe of the underlying interaction. This mechanism for muon decay shares many characteristic with beta decay. When both processes were experimentally proven to maximally violate parity, the long-held premise that nature was mirror-symmetric, it was clear that they represent instances of the same theory: The universal vector-axial interaction \cite{sudarshan1958chirality,feynman1958theory} of weak theory. So muon decay arguably has served for many decades as almost a kind of Rosetta stone for particle physics. But -- has its translative potential truly been exhausted yet?

The apparently local interaction governing muon decay was eventually understood as the low-energy limit of the mediation by a $W$ boson. Invisible at the energy scale of muon decay, this heavy field becomes a dynamic degree of freedom only in the high-energy regime such as present in particle colliders \cite{arnison1983w,banner1983w} -- but leaves an imprint on the low-energy regime nevertheless: Its mass determines the decay rate of the muon. Could there be other heavy particles veiled by a similar fashion? In order to bridge the energy regimes and systematically probe for new physics, we perturb the standard set of Michel parameters $\rho,\eta,\eta'',\delta,\xi,\xi',\xi'',\tfrac{\alpha'}{A},\tfrac{\beta'}{A}$ and the Fermi constant $G_F$ with Low Energy Effective Field Theory (LEFT) operators. They represent the adequate physical description for interactions below the electroweak scale and by matching them onto the Standard Model Effective Field Theory (SMEFT) we trace the perturbations upward the energy regime. At the scale of new physics we open up the SMEFT operators and ask: Which heavy new fields could produce them? This approach allows us to systematically collect all Beyond Standard Model (BSM) fields which at tree-level produce effective operators that could significantly perturb muon decay. We constraint them by translating the experimental values of the Michel parameters into boundaries for the masses of the heavy new fields. In this way we hold yet another careful reading of the muonic Rosetta stone and ask: Has the illustrious sibling already revealed each of its secrets?

The structure of the thesis is as follows: The second chapter revolves around the Michel parameters of muon decay, whereby we first briefly introduce the Standard Model and its effective description of muon decay. Subsequently the parameters themselves are introduced and how they shed light on the underlying interaction of muon decay. We end the chapter with a discussion of the observational values for the Michel parameters and how to appropriately account for their correlations. 

In Chapter three the LEFT is introduced and we show how the presence of LEFT operators perturbs the Michel parameters. We present our calculations for the case of SM-like neutrino output and the full four-fermion results. This allows us to present constraints for various LEFT operators of dimension six. Finally the Renormalization group equations are applied to evolve the LEFT operators up to the electroweak scale.

As we reach the electroweak scale in chapter four, we are ready to delve into the Standard Model Effective Field Theory (SMEFT). After a brief introduction we match our LEFT operators of dimension six to SMEFT operators of dimensions six and seven. By using the Renormalization Group Equations we translate the LEFT boundaries to SMEFT boundaries both at the EW scale and at the scale of new physics. 

At the energy scale of new physics in chapter five we use the complete tree-level dictionary for BSM fields to SMEFT operators up to dimension seven to collect the possible UV origins for our set of SMEFT operators. 

In chapter six we present our results: The constraints for a charged electroweak scalar doublet and a correlated neutral vector mediator which produces multiple LEFT operators that perturb muon decay are derived. We discuss the effects an extension of the SM by each of the heavy particles would have and compare our constraints to the literature.

The outlook in chapter seven gives us the opportunity to reflect on the directive of our investigation. Starting from the low-energy regime and the experimental data for muon decay, our approach is effectively bottom-up: We expand the theoretical space of muon decay and trace the perturbations up to the ultra-violet -- and subsequently back down to muon decay. This contradicts the implicit assumption of many studies, starting from a concrete UV model, that the \enquote{true} theory is the one valid at all energies -- that the \textit{true} is the \textit{whole}, so to say. But experimental data and theoretical descriptions stretch a far more complex space than such a top-down system. For the course of this thesis we consequently focus on the translations between energy regimes and operator bases -- the renormalization group equations and matching conditions -- rather than on the UV models themselves. 
\section{Notation}
We use the following notation:
\begin{table}[h]
    \centering
    \begin{tabular}{ll}
         $L^{(d)}_i$ &  Wilson coefficients of the LEFT\\
         $C^{(d)}_i$ &  Wilson coefficients of the SMEFT\\
         $\mathcal{O}_X$ & Operators of the LEFT or the SMEFT\\
         $l$ & Lepton doublet\\
         $v$ & VEV\\
         $S_i$ & UV scalar\\
         $F_i$ & UV fermion\\
         $V_i$ & UV vector \\
         $q$ & Muon momentum\\
         $m_\mu$ & Muon mass\\
         $p$ & Electron momentum\\
         $E_e$ & Electron energy\\
         $m_e$ & Electron mass\\
         $k_i$ & Neutrino momenta
    \end{tabular}
    \caption{The notation used in this thesis.}
    \label{tab:notation}
\end{table}

\chapter{Muon decay in the Standard Model}
We briefly introduce the Standard Model (SM) of particle physics -- its gauge symmetry, field contents and Lagrangian density. For muon decay, however, we do not need the full SM: At the energy scale of the muon mass, the $W$ boson is not a dynamical degree of freedom and can be intergrated out. The effective theory obtained this way resembles the historical four-Fermi theory and we use it to derive the decay width for muon decay.
\begin{figure}[ht]
    \centering
    \vspace*{1.5cm}
    \hspace*{-2cm}
    \begin{fmfgraph*}(80,60)
        \fmfleft{l1,l2}
        \fmfright{r1,r2,r3}
        \fmf{fermion,tension=2}{l2,o1}
        \fmf{fermion}{o1,r3}
        \fmf{photon,label=$W^-$}{o1,o2}
        \fmf{fermion}{r1,o2,r2}
        \fmflabel{$\mu^-$}{l2}
        \fmflabel{$\nu_\mu$}{r3}
        \fmflabel{$e^-$}{r2}
        \fmflabel{$\bar\nu_e$}{r1}
    \end{fmfgraph*}\par\vspace{1cm}
    \caption{Muon decay in the Standard Model}
    \label{fig:mdSM}
\end{figure}
\section{The Standard Model}
The fundamental structuring principle for the Standard Model is its local gauge symmetry group:
\begin{equation}
    SU(3)_C \times SU(2)_L \times U(1)_Y
\end{equation}
The first subgroup $SU(3)_C$ defines quantum chromodynamics, the theory of strong interactions governing quarks by means of 8 massless gluons. The other two subgroups $SU(2)_L \times U(1)_Y$ jointly represents the unified electroweak sector. The three generations of left-handed lepton fields $L$ and quark fields $Q$ transform as doublets under $SU(2)_L$, whereas the right-handed lepton $e$, up-type quark $u$ and down-type quark $d$ transform as singlets. The gauge bosons of $SU(2)_L$ thus couple exclusively to left-handed fermions, rendering electroweak theory chiral. The completing building block of the SM is the Higgs doublet $H$. After it acquires a vacuum expectation value (VEV) at an energy of $v\approx 246\text{ GeV}$ it no longer transforms as a $SU(2)_L$ doublet but as a singlet -- and the three remaining degrees of freedom are absorbed by three vector bosons obtaining a mass. The Higgs VEV thus defines the electroweak scale, below which the $SU(2)_L \times U(1)_Y$ symmetry group is broken and only the electromagnetic $U(1)_\text{em}$ symmetry applies. Above the electroweak scale the entire symmetry group $SU(2)_L \times U(1)_Y$ applies and the full Lagrangian density is constructed as a renormalizable and gauge invariant theory:
\begin{align}
\mathcal{L} _{\rm SM} &= -\frac14 G_{\mu \nu}^A G^{A\mu \nu}-\frac14 W_{\mu \nu}^I W^{I \mu \nu} -\frac14 B_{\mu \nu} B^{\mu \nu}
+ (D_\mu H^\dagger)(D^\mu H)+\sum_{\psi=Q,L,e,u,d} \overline \psi\, i \slashed{D} \, \psi\nn
&-\lambda \left(H^\dagger H -\frac12 v^2\right)^2- \biggl[ H^{\dagger j} \overline d\, Y_d\, Q_{j} + \widetilde H^{\dagger j} \overline u\, Y_u\, Q_{j} + H^{\dagger j} \overline e\, Y_e\,  L_{j} + \hbox{h.c.}\biggr]
\label{eq:LagSM}
\end{align}
The gauge covariant derivative is:
\begin{equation}
    D_\mu = \partial_\mu + i g_3 T^A A^A_\mu + \tfrac{1}{2}i g_2  \sigma^I W^I_\mu + i g_1 \hyp B_\mu
\end{equation}
Here, $T^A$ are the $SU(3)_C$ generators,  the Pauli matrices $\tfrac{1}{2}\sigma^I$ generate $SU(2)_L$, and $\hyp$ is the $U(1)_Y$ hypercharge generator.  The fermion fields have weak-eigenstate indices $r=1, ..., 3$ and the Yukawa couplings are $3\times 3$ matrices. $\widetilde H$ is defined by
\begin{align}
\widetilde H_j &= \epsilon_{jk} H^{\dagger\, k}
\end{align}
Here, an $SU(2)$ invariant tensor $\epsilon_{jk}$ is defined by $\epsilon_{12}=1$ and $\epsilon_{jk}=-\epsilon_{kj}$, $j,k=1,2$.\par
As we mentioned, the vacuum of the Higgs is invariant under the electromagnetic generator
$Q=T^3+\hyp$, but not under the remaining generators of
$SU(2)_L\times U(1)_Y$. The electroweak symmetry is therefore
spontaneously broken according to:
\begin{equation}
    SU(2)_L\times U(1)_Y
    \longrightarrow U(1)_{\mathrm{em}}
\end{equation}
In unitary gauge, the Higgs doublet can be written as:
\begin{equation}
    H(x)
    =
    \frac{1}{\sqrt{2}}
    \begin{pmatrix}
        0\\ v+h(x)
    \end{pmatrix}\label{eq:higgsVEV}
\end{equation}
It is useful at this point to introduce the gauge fields of the broken
phase. The charged fields are defined by:
\begin{equation}
    W_\mu^\pm
    =
    \frac{1}{\sqrt{2}}
    \left(
        W_\mu^1\mp iW_\mu^2
    \right)\label{eq:physicalW}
\end{equation}
The two neutral gauge fields mix according to:
\begin{align}
    Z_\mu
    &=
    \cos\theta_W\,W_\mu^3
    -
    \sin\theta_W\,B_\mu
    \\
    A_\mu
    &=
    \sin\theta_W\,W_\mu^3
    +
    \cos\theta_W\,B_\mu
\end{align}
The rotation into the physical fields is in terms of the weak mixing angle $\theta_W$:
\begin{equation}
    \tan\theta_W=\frac{g_1}{g_2}
    \qquad
    \sin\theta_W
    =
    \frac{g_1}{\sqrt{g_1^2+g_2^2}}
    \qquad
    \cos\theta_W
    =
    \frac{g_2}{\sqrt{g_1^2+g_2^2}}
\end{equation}
Substituting the Higgs VEV of Eq.~\ref{eq:higgsVEV} into the Higgs kinetic term in
gives:
\begin{align}
    (D_\mu H)^\dagger(D^\mu H)
    \supset{}&
    \frac{g_2^2v^2}{4}W_\mu^+W^{-\mu}
    +
    \frac{v^2}{8}
    \left(g_2W_\mu^3-g_1B_\mu\right)^2
    \\
    ={}&
    m_W^2W_\mu^+W^{-\mu}
    +
    \frac{1}{2}m_Z^2Z_\mu Z^\mu ,
\end{align}
So the mass of the $W$ boson is set by the Higgs VEV and experimentally measured as $80.3692(133)\text{ GeV}$ \cite{particle2024review} -- three order of magnitudes higher than the mass of the muon, which resides at $105.6583755(23)\text{ MeV}$ \cite{particle2024review} and which is why we are permitted to integrate the $W$ boson out of the description of muon decay, as we will now see.
\section{Weak interactions and leptonic currents}
We wish to find a description for muon decay in the SM, whereby we will use the following momenta convention:
\begin{equation}
    \mu^-(q)\to \bar{\nu}_e(k_1)+\nu_\mu(k_2)+e^-(p)
    \label{eq:md}
\end{equation}
From the fermion sector of the SM Lagrangian \ref{eq:LagSM} we can derive the interaction term of the lepton doublet with the $SU(2)_L$ gauge bosons:
\begin{equation}
    \mathcal{L}_\text{int}^{SU(2)_L}=\bar{L} i\gamma^\mu \tfrac{1}{2}i g_2  \sigma^I W^I_\mu L
\end{equation}
We expand the product of the Pauli matrices and the weak gauge fields and rotate the latter to their physical, charged fields according to Eq.~\ref{eq:physicalW}. We then find the following interaction terms (dropping flavor indices):
\begin{equation}
    \mathcal{L}_\text{int}^{SU(2)_L} = \frac{g_2}{2} \left( \bar{\nu}_L \gamma^\mu W_\mu^3 \nu_L - \bar{e}_L \gamma^\mu W_\mu^3 e_L + \sqrt{2} \bar{\nu}_L \gamma^\mu W_\mu^+ e_L + \sqrt{2} \bar{e}_L \gamma^\mu W_\mu^- \nu_L \right)
\end{equation}
We see that two interactions $\sqrt{2} \bar{e}_L \gamma^\mu W_\mu^- \nu_L$ allow us to describe muon decay as mediated by a $W^-$ boson (see fig. \ref{fig:mdSM}). The full invariant amplitude can the be derived by applying Feynman rules, notably inserting two vertex factors $i\tfrac{g}{\sqrt{2}}\gamma^\mu P_L$ and the $W^-$ propagator $-i\left(g_{\mu\nu}-\frac{k_\mu k_\nu}{m_W^2}\right)/\left(k^2-m_W^2\right)^{-1}$, and setting the spinor momenta according to \ref{eq:md}:
\begin{equation}
    \im\mathcal{M}=\Big[\bar u(k_1)\big(i\tfrac{g}{\sqrt{2}}\gamma^\mu P_L\big)u(q)\Big]
                    \frac{-i\Big(g_{\mu\nu}-\frac{k_\mu k_\nu}{m_W^2}\Big)}{k^2-m_W^2}
                    \Big[\bar u(p)\big(\im\tfrac{g}{\sqrt{2}}\gamma^\mu P_L\big)u(k_2)\Big]
    \label{eq:mdAmp}
\end{equation}
The momentum $k$ of the $W^-$ boson relates to the momenta of the interacting leptons by $k=q-k_1=k_2+p$ -- but since the momentum exchange is very small compared to the $W$ boson mass, we can expand the propagator:
\begin{equation}
    \frac{-\im\Big(g_{\mu\nu}-\frac{k_\mu k_\nu}{m_W^2}\Big)}{k^2-m_W^2} \approx
    \im\frac{g_{\mu\nu}}{m_W^2}+\mathcal{O}\left(\frac{k^2}{m_W^4}\right)
\end{equation}
Substituting this simplified propagator back into the amplitude yields:
\begin{equation}
    \mathcal{M}=-\frac{g_2^2}{2m_W^2}[\bar u(k_1)\gamma^\mu P_L u(q)][\bar u(p)\gamma_\mu P_L u(k_2)]
\end{equation}
This result resembles the theory that Fermi presented in 1934 for beta decay \cite{fermi1934tentativo} where he conceptualized the process by a four-fermion vector-vector interaction Hamiltonian. He already speculated about a heavy, undiscovered mechanism and gave an account of the decay as an approximated local interaction between the two weak currents. In order to connect his theory with the SM and frame it as a low energy limit of a UV complete Lagrangian we are adjusting Fermi's approach to modern terminology, namely by writing it as an interaction term of two charged currents in a Lagrangian:
\begin{equation}
    \mathcal{L}_\text{IR} = \frac{-4G_F}{\sqrt{2}} \left[ \bar{e} \gamma^\mu P_L \nu_e + \bar{\mu} \gamma^\mu P_L \nu_\mu \right] \left[ \bar{\nu}_e \gamma_\mu P_L e + \bar{\nu}_\mu \gamma_\mu P_L \mu \right]
    \label{eq:LagIR}
\end{equation}
The coupling constant $G_F$ of this interaction has a factor $\tfrac{4}{\sqrt{2}}$ to appropriately connect this direct coupling to the SM parameters, as we will see soon. The IR Lagrangian gives as a Feynman rule the four-fermion vertex $-i \frac{4G_F}{\sqrt{2}} [\gamma^\mu P_L]_{ab} [\gamma_\mu P_L]_{cd}$ where $abcd$ are Dirac indices of the respective spinor pairs. Contracting the vertex with the four particles of muon decay gives the amplitude:
\begin{equation}
    \mathcal{M}=-\frac{4G_F}{\sqrt{2}}[\bar u(k_1)\gamma^\mu P_L u(q)][\bar u(p)\gamma_\mu P_L u(k_2)]
    \label{eq:Amp4F}
\end{equation}
We see that the IR Lagrangian and the SM Lagrangian in the large $m_W$ limit produce the same amplitude for muon decay if:
\begin{equation}
    -\frac{4G_F}{\sqrt{2}}=-\frac{g_2^2}{2m_W^2}\implies G_F=\frac{\sqrt{2}g_2^2}{8m_W^2}
\end{equation}
In order to make the connection between the the two theories precise we will now give an account of how to generally integrate out heavy fields from the Lagrangian and how to apply the technique to the $W^-$ boson.
\begin{figure}[t]
    \centering
    \sbox{\muondecaySM}{
        \begin{fmfgraph*}(80,60)
            \fmfleft{l1,l2}
            \fmfright{r1,r2,r3}
            \fmf{fermion,tension=2}{l2,o1}
            \fmf{fermion}{o1,r3}
            \fmf{photon,label=$W^-$}{o1,o2}
            \fmf{fermion}{r1,o2,r2}
            \fmflabel{$\mu^-$}{l2}
            \fmflabel{$\nu_\mu$}{r3}
            \fmflabel{$e^-$}{r2}
            \fmflabel{$\bar\nu_e$}{r1}
        \end{fmfgraph*}
    }
    \sbox{\muondecayLEFT}{
        \begin{fmfgraph*}(80,60)
            \fmfleft{l1,l2}
            \fmfright{r1,r2,r3}
            \fmf{fermion,tension=2}{l2,o1}
            \fmf{phantom,tension=25}{o1,o2}
            \fmf{fermion}{o1,r3}
            \fmf{fermion}{r1,o2,r2}
            \fmflabel{$\mu^-$}{l2}
            \fmflabel{$\nu_\mu$}{r3}
            \fmflabel{$e^-$}{r2}
            \fmflabel{$\bar\nu_e$}{r1}
        \end{fmfgraph*}
    }
    \begin{tikzpicture}
        \node (mdSM) {\usebox{\muondecaySM}};
        \node[anchor=west] (mdLEFT) at (mdSM.east) [xshift=3.3cm] {\usebox{\muondecayLEFT}};
        \draw[->, thick, dotted, shorten <=25pt, shorten >=-5pt] (mdSM.south) to[bend right] node[midway, below] {$q^2\ll m_W^2$} (mdLEFT.south);
        \node[anchor=north] at (mdSM.south) [yshift=0.5cm] {$\frac{g^2}{q^2-m_W^2}$};
        \node[anchor=north] (mdLEFTform) at (mdLEFT.south) [yshift=1.2cm, xshift=-0.6cm] {$-\frac{g^2}{m_W^2}\propto G_F$};
    \end{tikzpicture}
    \caption{The four-fermion contact interaction}
    \label{fig:mdSM2FF}
\end{figure}
\section{Integrating out heavy fields}\label{sec:intoutHDOF}
In the path integral formulation, the dynamics of a system are captured by the generating functional, which integrates over all possible field configurations. For a theory containing both light fields, denoted collectively as $\phi$, and a heavy field $\Phi$ with mass $M \gg E$ (where $E$ is the typical energy scale of the physical processes under consideration), the effective action $S_{\text{eff}}[\phi]$ for the light degrees of freedom is formally defined by integrating out the heavy field \cite{Donoghue:1992dd}:
\begin{equation}
    e^{i S_{\text{eff}}[\phi]} = \int \mathcal{D}\Phi \, e^{i S[\phi, \Phi]}
\end{equation}
If the heavy field $\Phi$ appears strictly quadratically in the action -- for instance, interacting with a current $J(\phi)$ composed of light fields via an interaction term $J \Phi$ -- the path integral is exactly solvable. It reduces to a multi-dimensional Gaussian integral. By shifting the integration variable to complete the square, the heavy field can be analytically integrated out, resulting in a non-local effective action. At low energies ($p^2 \ll M^2$), the heavy propagator $1/(p^2 - M^2)$ collapses to a point-like contact interaction $-1/M^2$, generating the local operators of the effective field theory.

However, if the heavy field possesses self-interactions or couples non-linearly, the path integral cannot be evaluated exactly. Instead, one employs the saddle-point approximation \cite{henning2015usestandardmodeleffective}. In the low-energy limit, the path integral is overwhelmingly dominated by the field configuration that minimizes the action. We expand the heavy field around this classical stationary configuration $\Phi_c$:
\begin{equation}
    \Phi = \Phi_c + \eta
\end{equation}
Here, $\eta$ represents quantum fluctuations, and $\Phi_c$ is the solution to the Euler-Lagrange Equation of Motion (EOM), defined by:
\begin{equation}
    \left. \frac{\delta S[\phi, \Phi]}{\delta \Phi} \right|_{\Phi = \Phi_c} = 0
\end{equation}
Expanding the action around $\Phi_c$ yields:
\begin{equation}
    S[\phi, \Phi] = S[\phi, \Phi_c] + \frac{1}{2} \left. \frac{\delta^2 S}{\delta \Phi^2} \right|_{\Phi_c} \eta^2 + \mathcal{O}(\eta^3)
\end{equation}
Notably, the linear term in $\eta$ vanishes by definition of the EOM. The integration over the quadratic quantum fluctuations $\eta$ generates a functional determinant, which corresponds to one-loop quantum corrections.
If we restrict our matching procedure strictly to tree-level, these loop corrections are neglected. The effective action simply becomes the original action evaluated at the classical minimum:
\begin{equation}
    S_{\text{eff}}^{\text{tree}}[\phi] \approx S[\phi, \Phi_c]
\end{equation}
Therefore, replacing a heavy mediator with its classical equation of motion and substituting it back into the Lagrangian evaluates the path integral at the tree-level saddle point. We now perform this calculation for the $W^\pm$ in order to derive the IR Lagrangian appropriate for muon decay. To this end we first isolate the terms in the electroweak Lagranigan that contain the $W^\pm$ boson:
\begin{equation}
    \mathcal{L}_ \text{SM} \supset -\frac{1}{2} W_{\mu\nu}^+ W^{-\mu\nu} + m_W^2 W_\mu^+ W^{-\mu} - \frac{g}{\sqrt{2}} \left( J^\mu W_\mu^- + J^{\mu \dagger} W_\mu^+ \right)
    \label{eq:LagW}
\end{equation}
The current $J^\mu$ is that of Eq. \ref{eq:LagIR}, containing only the first and second generation of leptons:
\begin{equation}
    J^\mu=\bar{e} \gamma^\mu P_L \nu_e + \bar{\mu} \gamma^\mu P_L \nu_\mu 
\end{equation}
To find the EOM we apply the Euler-Lagrange equation:
\begin{equation}
    \frac{\partial \mathcal{L}}{\partial W_\mu^-} - \partial_\nu \frac{\partial \mathcal{L}}{\partial (\partial_\nu W_\mu^-)} = 0
\end{equation}
This gives:
\begin{equation}
    m_W^2 W_\mu^+ - \frac{g}{\sqrt{2}} J^\mu + \partial^\nu W_{\nu\mu}^+ = 0
\end{equation}
The derivative of $W_{\nu\mu}$ corresponds to the momentum $p$ and the term can therefore be neglected in the low energy limit, i.e. $\tfrac{p}{m_W}\approx\tfrac{m_\mu}{m_W}\approx 0$. We therefore have:
\begin{equation}
    W_\mu^+ \approx \frac{g}{\sqrt{2} m_W^2} J^\mu
\end{equation}
Substituting this back into the Lagrangian \ref{eq:LagW} gives the familiar charged current interaction of Eq. \ref{eq:LagIR}:
\begin{equation}
    \mathcal{L}_\text{IR} = - \frac{g^2}{2 m_W^2} J^{\mu \dagger} J_\mu
\end{equation}
So this Lagrangian is an adequate description at $E=m_\mu=105.6583755(23)\text{ MeV}$, which three orders of magnitudes below $m_W=80.3692(133)\text{ GeV}$ \cite{particle2024review}.\par 
Systematically constructing an whole \textit{Effective Field Theory (EFT)} from a UV complete theory by integrating out the heavy degree of freedom represents an example of a \textit{top-down} approach to EFTs: Starting from a full theory valid in the high energy limit, in this case the SM Lagrangian, one derives the appropriate effective theory for the low energy limit. By requiring the amplitudes in both theories to agree in a certain energy regime, one \textit{matches} the couplings of the effective theory, called \textit{Wilson coefficients}, to the UV complete theory. The four-Fermi theory, on the other hand, historically started from the opposite direction: Without knowledge of the SM and the underlying heavy mechanism it gave an adequate description for weak interactions starting \textit{from scratch}, namely directly from the known, interacting fields. Experimental evaluation then allowed for constraining the effective couplings. For the course of this work both methods are applied to some extent but the overall directive is a bottom-up approach: Starting from the phenomenology of muon decay well below the electroweak scale we investigate possible UV origins. In the next section we see how to to determine Fermi's constant from the total decay width - the kinematic space of the electron is, however, far more complex and in the following section we get acquainted with the Michel parameters that form its foundation. The $\rho$ parameter, for example, brings to light the V-A nature of the weak interaction, as we will briefly show.
\section{Effective Field Theories}
Effective Field Theory (EFT) is the idea that physics is separated by scales -- or more precisely it is a set of ideas that justify why this separation works \cite{brivio2019standard}. Notably the Standard Model (SM) has been validated up to LHC energies despite not being well-defined in the UV regime due to Landau poles in the electroweak sector. Renormalization works by regularizing the short-distance physics and absorbing the effects in low energy effective parameters -- see the detailed discussion in chapter 4 of \cite{brivio2019standard} for how dimensional regularization and the $\overline{MS}$ scheme correspond to a systematic removal of UV physics which is the key principle of an EFT. The Lagrangian is then constructed as combining the IR physics as a local, analytical operator expansion with the UV physics controlling the corresponding Wilson coefficients.
\begin{equation}
    \mathcal{L}_\text{EFT}\simeq\sum_iL_i^\text{UV}(\mu)O_i^\text{IR}(\mu)
\end{equation}
The first two \enquote{prime directives} of EFTs are thus in \cite{brivio2019standard} summarized as:
\begin{enumerate}
    \item determine a characteristic scale in observables
    \item construct the Lagrangian out of the degrees of freedom, i.e. the propagating on-shell states
\end{enumerate}
Starting from the SM two characteristic scales directly come to mind: The scale of electroweak symmetry breaking dictated by the Higgs VEV $v$ and the scale of heavy new physics $\Lambda$. Depending on how one assumes heavy new physics to relate to the Higgs mechanism two  EFTs can be set up by the scale $\Lambda$:
\begin{enumerate}
    \item if one assumes the Higgs to enter the EFT in the form of its scalar doublet $H$, the full SM field content above the EW scale is used to construct the Standard Model Effective Field Theory (SMEFT).
    \item the Higgs Effective Field Theory (HEFT) instead treats the physical Higgs and the Goldstone bosons separatedly. The framework can be described as a chiral Lagrangian with a light scalar h, organized by a loop/chiral expansion rather than by canonical dimension \cite{buchalla2014complete}.
\end{enumerate}
While the HEFT might be more general and appropriate for non-linear UV models, the SMEFT is in our case the more restrictive and diagnostic tool: As its operators adhere to the full $SU(3)_C\times SU(2)_L\times U(1)_Y$ gauge group one can systematically derive all UV representations that would generate specific operators -- rendering the SMEFT a powerful mode of investigation in the bottom-up spirit \cite{li2022bottom}.\par 
The scale $v$ on the other hand separates the IR regime from the full SM and allows expanding the physics above the EW scale by Wilson coefficients divided by powers of $v$ and the effects of UV physics by powers of $\Lambda$. We will first introduce the LEFT in Ch.~\ref{ch:left} and subsequently the SMEFT in Ch.~\ref{ch:smeft}.
\section{Total decay width}
We want to derive the total decay width for muon decay from the amplitude of four-Fermi theory:
\begin{equation}
    \mathcal{M}=-\frac{4G_F}{\sqrt{2}}[\bar u(k_1)\gamma^\mu P_L u(q)][\bar u(p)\gamma_\mu P_L u(k_2)]
    \label{eq:Amp4F2}
\end{equation}
If the spins of the outgoing particles are not measured, the expectation value of the squared amplitude is the following spin sum:
\begin{equation}
    \langle|\mathcal{M}|^2\rangle=\sum_{\text{spins}}|\mathcal{M}|^2
\end{equation}
Squaring the amplitude, we obtain:
\begin{equation}
\begin{aligned}
    |\mathcal{M}|^2=\left(\frac{g^2}{2m_W^2}\right)^2&[\overline{u}(k_1)\gamma^\mu P_L u(q)][\overline{u}(p)\gamma_\mu P_Lu(2)]\\ &\times [\overline{u}(k_1)\gamma^\mu P_L u(q)]^*[\overline{u}(p)\gamma_\mu P_Lu(2)]^*
    \label{eq:T2}
\end{aligned}
\end{equation}
A detailed derivation can be found in \cite{griffiths2020introduction} - by applying the completeness theorem for the spin sums and trace identities, the amplitude evaluates to:
\begin{equation}
    \langle|\mathcal{M}|^2\rangle=2\left(\frac{g}{M_W}\right)^4(q\cdot k_3)(k_2\cdot p)
    \label{eq:T2summedF}
\end{equation}
The decay rate is calculated by applying Fermi's Golden Rule:
\begin{equation}
    \frac{\de^3\Gamma}{\de^3p} = \frac{1}{(2\pi)^5}\frac{1}{4m_\mu E_e}
    \int\frac{\de^3k_1}{2E_1}\frac{\de^3k_2}{2E_2}
    \delta^{(4)}(q-k_1-k_2-p)\langle|\mathcal{M}|^2\rangle
    \label{eq:fermiGoldenRule}
\end{equation}
The neutrino phase space is evaluated using the following identity \cite[p.308]{scheck2012electroweak}:
\begin{equation}
    \int\frac{d^3p_2}{2E_2}\frac{d^3p_3}{2E_3}=\delta^{(4)}(Q-p_2-p_3)p_2^\alpha p_3^\beta=\frac{\pi}{24}(g^{\alpha\beta}Q^2+2Q^\alpha Q^\beta)
    \label{eq:neutrinoPhaseSpace}
\end{equation}
Integrating over all electron momenta gives in the $m_e\to 0$ limit the total decay width:
\begin{equation}
    \Gamma =\frac{G_F^2m_\mu^5}{192\pi^3}
    \label{eq:totalDW}
\end{equation}
Current experiments determine $G_F$ to the value \parencite{particle2024review}:
\begin{equation}
    G_{F,SM}=\qty{1.166385(6)e-5}{\per\giga\electronvolt\squared}
    \label{eq:GFexp}
\end{equation}
\section{The Michel parameters}\label{sec:michel}
In his work from 1950 \cite{michel1950interaction}, Michel set up the most general Lorentz-invariant Hamiltonian in order to analyze the kinematics of four-fermion interactions:
\begin{equation}
    \mathcal{H}_{1950} = \sum_{i=S,V,T,A,P} L_i (\bar{\psi}_e \Gamma_i \psi_\mu) (\bar{\psi}_{\nu} \Gamma_i \psi_{\nu}) + \text{h.c.}
\end{equation}
Because he did not consider the possibility that such an interaction might violate parity, his Hamiltonian only couples currents of the same kind -- scalar, vector, tensor, axial and pseudo -- to each other, resulting in a overall parity-conserving expression. There are thus 5 coupling constants $L_i$ that define his four-fermion theory. Critically, however, Michel realized that the exact weights of the five different interaction types in the overall expression only alter the energy spectrum by a single constant. This is because the neutrinos are not observed and integrated over. When evaluating the squared amplitude $|\mathcal{M}|^2$ one first obtains a product of two traces: One for the electron-muon line and one for the neutrino line, the latter containing the neutrino momenta $k_1,k_2$ - integrating over the neutrino phase space thus amounts to solving the following integral:
\begin{equation}
    I^{\alpha\beta} = \int \frac{d^3k_1}{2E_1} \frac{d^3k_2}{2E_2} \delta^4(Q - k_1 - k_2) k_1^\alpha k_2^\beta
\end{equation}
By pure Lorentz invariance, this can only depend on the total transferred momentum $Q=q-p$ and the only two rank-two tensors one can build out of $Q$ are the metric tensor $g^{\alpha\beta}$ and the product $Q^\alpha Q^\beta$:
\begin{equation}
    I^{\alpha\beta} = A(Q^2) g^{\alpha\beta} + B(Q^2) Q^\alpha Q^\beta
\end{equation}
Because the neutrinos are considered to be massless, dimensional analysis and Lorentz invariance fix $A(Q^2)\propto g^{\alpha\beta}$ and $B(Q^2)\propto 1$. If one contracts the neutrino tensor with the muon-electron trace, which is built of $q_\mu$ and $p_e$, scalars of the following kind can be formed:
\begin{align}
    Q^2 = (q_\mu - p_e)^2 \approx m_\mu^2 - 2m_\mu E_e\\
    Q \cdot q_\mu = (q_\mu - p_e) \cdot q_\mu \approx m_\mu^2 - m_\mu E_e\\
    Q \cdot p_e = (q_\mu - p_e) \cdot p_e \approx m_\mu E_e
\end{align}
Every single invariant is simply linear function of the elctron energy $E_e$ and because of the phase space factor $E_e^2\de E_e$ the differential decay rate $\de\Gamma/\de E_e$ is a polynomial in $E_e^2$ and $E_e^3$. Thus a single parameter suffices to describe the relative weight between the quadratic and cubic terms and thus the shape of the spectrum -- the original Michel parameter $\rho$. Expressed in terms of the dimensionless variable $x=E/W$ with W the maximum electron energy, the spectrum can be written as:
\begin{equation}
    P(x) dx \propto x^2 \left[ 3(1 - x) + 2\rho \left( \frac{4}{3}x - 1 \right) \right] dx
\end{equation}
Current global data gives for the $\rho$ parameter:
\begin{equation}
    \rho = 0.74979\pm 0.00026
\end{equation}
This implies the interaction to consist of vector and axial vector terms:
\begin{equation}
    \rho = \frac{3L_V^2 + 3L_A^2 + 6L_T^2}{L_S^2 + L_P^2 + 4L_V^2 + 4L_A^2 + 6L_T^2}
    \label{eq:rho}
\end{equation}
For $L_V=-L_A=1$ this gives $\rho=\tfrac{3}{4}$. However, from the experimental value of the $\rho$ parameter the exact ratio of the $V$ and $A$ term or whether axial and vector current mix to form a parity-violating term can not be deduced. A paradigm shift took place in 1956 when Lee and Yang, in order to explain the $\theta-\tau$ puzzle \parencite{lee1956question}, suggested that the weak force might violate parity  -- a breakthrough idea, which was experimentally tested a year later by Chien-Shiung Wu \parencite{wu1957experimental}: She confirmed that the electrons in Cobalt-60 decay are emitted preferentially in the direction opposite to the nuclear spin. In the phenomenology of muon decay parity violation generates an anisotropic part in the electron spectrum. If one measures the polarization of the decaying muon and the angle by which the electron is emitted relative to that polarization the probability is not distributed isotropically at all: The electron is most likely to be ejected in the exact opposite direction to the muon's polarization. And this clearly violates parity, since mirroring of the experiment would leave the polarization invariant but spatially reverse the electron momentum.\par 
In 1957, two groups (Bouchiat \& Michel \cite{bouchiat1957theory} and, independently, Kinoshita \& Sirlin \cite{kinoshita1957muon,kinoshita1957polarization}) expanded the formalism to match the new chiral reality and introduced additional
parameters for the anisotropic part of the spectrum. The parameter $\xi$ measures parity violation and $\delta$ is again a shape parameter for the anisotropic part of the spectrum. 
\begin{equation}
    \begin{aligned}
        \frac{\de ^2\Gamma_\text{anisotr.}}{\de x\de \hspace{-0.1em}\cos\theta}\propto\xi P_\mu\cos\theta \left( x^2(1-x)+2\delta\left(\frac{4}{3}x-1\right)\right)
    \end{aligned}
\end{equation}
As one can see, the peak of the differential decay width is located at $\theta=180^\circ$, the direction opposite to the muon polarization vector. The SM predicts $\xi = 1$, maximal parity violation, and $\delta = \frac{3}{4}$, which means the shape of the anisotropic spectrum is also V-A-like. Both SM predictions are supported by current data \parencite{particle2024review}:
\begin{align}
    P_\mu\xi &=1.0009 \pm 0.0012\\
    \xi P_\mu\delta/\rho&=1.00179 \pm 0.00114
\end{align}
Both groups also introduced a parameter $\eta$ to account for the finite electron mass shaping the low-energy end of the spectrum. It reflects the interference of scalar and tensor terms, which can only produce a non-zero term in the amplitude via mass insertion. As the Standard Model neither has tensor nor scalar interaction terms, it predicts $\eta$ to be zero, which current experimental data supports \parencite{particle2024review}:
\begin{equation}
    \eta =0.057 \pm 0.034
\end{equation}
The full spectrum of the electron, when the muon's polarization and the electron's momentum direction relative to it are measured (but not the spin of the electron), thus has the following form:
\begin{equation}
    \begin{aligned}
        \frac{\de ^2\Gamma}{\de x\de \hspace{-0.1em}\cos\theta}=&\frac{m_\mu}{4\pi^3}W^4G_F^2
        \sqrt{x^2-x_0^2}
        \left[2x(1-x)+\frac{4}{9}\rho(4x^2-3x-x_0^2)\right.\\
        &+2\eta x_0(1-x)-\frac{2}{3}\xi P_\mu\cos\theta\sqrt{x^2-x_0^2}\Big( [1-x]\\
        &\hspace{2.2cm}+\left.\frac{2}{3}\delta\left[(4x-3)+((1-x_0^2)^{-1}-1)\right]\Big)\right]
        \label{eq:MichelNoPe}
    \end{aligned}
\end{equation}

If the spin of the electron is also measured, additional Michel parameters can be defined: $\xi'$, $\xi''$, $\eta''$, $\alpha'/A$, and $\beta'/A$, whereby $A$ is an overall normalization factor and not independently determined. In order to account for the elctron spin dependence, we extend the parameterization of the differential decay rate as in \cite{particle2024review}:

\begin{equation}
\frac{d^2\Gamma}{dx\, d\cos\theta}=\frac{m_\mu }{4\pi^3} W^4G_F^2\sqrt{x^2-x_0^2}\left[F_{IS}(x)-P_\mu\cos\theta F_{AS}(x)+\hat \zeta\cdot \vec P_e(x,\theta)\right]\, ,
\end{equation}

The unit vector $\hat\zeta$ defines the direction along which the electron polarization, $\vec P_e$, is analyzed. From a theory perspective, we choose $\hat\zeta$ such that it projects the polarization onto a specific axis, such as longitudinal or transverse polarizations. The isotropic and anisotropic distributions correspond to those defined in Eq. \ref{eq:MichelNoPe} -- the factor of $2$ difference relative to Eq. \ref{eq:MichelNoPe} is because we no longer sum over the electron spin states.
\begin{eqnarray}
F_{IS}(x)&=&x(1-x)+\frac{2}{9}\rho(4x^2-3x-x_0^2)+\eta x_0(1-x)\, ,\\
F_{AS}(x)&=&\frac{\xi}{3}\sqrt{x^2-x_0^2}\left[(1-x)+\frac{2}{3}\delta\left(\left[4x-3\right]+\left[\sqrt{1-x_0^2}-1\right]\right)\right]\, .
\end{eqnarray}
%
The additional electron spin dependent terms follow from expanding the electron spin dependence as $\vec P_e(x,\theta)=P_{T_1}\hat x_1+P_{T_2}\hat x_2+P_L\hat x_3$. The direction $\hat x_3$ is chosen along the electron momentum, $\hat x_1$ lies in the decay plane and is transverse to the electron momentum, and $\hat x_2$ is orthogonal to both $\hat x_1$ and $\hat x_3$. The kinematics are outlined in further detail in appendix~\ref{ch:kinematics} and specifically our method for projecting out each Michel parameter is discussed in appendix~\ref{sec:michelbasis}. With this choice, the polarization components are given by:
%
\begin{align}
P_{T_1}=&\frac{P_\mu\sin\theta}{12}\left[-2\left(\xi''+12\left[\rho-\frac{3}{4}\right]\right)(1-x)x_0\right.\nonumber\\[5pt]
&\hspace*{4cm} -3\eta\left(x^2-x_0^2\right)+\eta''\left(-3x^2+4x-x_0^2\right)\Big]
\\[5pt]
P_{T_2}=&\frac{P_\mu\sin\theta}{3}\sqrt{x^2-x_0^2}\left[3\frac{\alpha'}{A}(1-x)+2\frac{\beta'}{A}\sqrt{1-x_0^2}\right]\\[5pt]
P_L=&-F_{IP}(x)+P_\mu\cos\theta F_{AP}(x)\\[5pt]
F_{IP}(x)=&\frac{\sqrt{x^2-x_0^2}}{54}\left[9\xi'\left(-2x+2+\sqrt{1-x_0^2}\right)\right.\nonumber\\[5pt]
&\hspace*{4.3cm}\left.+4\xi\left(\delta-\frac{3}{4}\right)\left(4x-4+\sqrt{1-x_0^2}\right)\right]\\[5pt]
F_{AP}(x)=&\frac{1}{6}\left[\xi''\left(2x^2-x-x_0^2\right)+4\left(\rho-\frac{3}{4}\right)\left(4x^2-3x-x_0^2\right)+2\eta''(1-x)x_0\right]
\end{align}
\clearpage
These different contributions can be qualitatively understood as follows:

\begin{itemize}
\item The longitudinal polarization, $P_L$, corresponds to alignment of the electron momentum with its polarization. The corresponding contributions in addition to the usual Michel parameters are $\xi'$, $\xi''$, and $\eta''$. These terms also include contributions from $\rho$ and $\delta$ as defined above.
\item The transverse polarization $P_{T_1}$ lies in the decay plane. It receives contributions from $\xi''$ and $\eta''$ along with contributions from the previously defined $\rho$ and $\eta$. It includes angular dependence in the form of $\sin\theta$. 
\item The transverse polarization $P_{T_2}$ is perpendicular to the decay plane. It is $T$-odd and therefore sensitive to time-reversal violating effects, we will see in the LEFT interpretation its only contributions (at leading order) come from operators that don't conserve lepton number. It also includes angular dependence of the form $\sin\theta$. The parameters $\alpha'/A$ and $\beta'/A$ are constrained by considering this kinematical region. We treat $\alpha'/A$ and $\beta'/A$ as independent variables, not ratios. Their labeling as ratios comes from a different, but equivalent treatment of the Michel parameters.
\end{itemize}

In the SM we have for these extended Michel parameters, $\eta''=0$, $\xi'=\xi''=1$, and $\alpha'/A=\beta'/A=0$. 

The Michel parameters lay the phenomenological foundation for the inquiry of this thesis: How would heavy fields from beyond the SM affect them? Expressing the Michel parameters in terms of SPVAT couplings as for $\rho$ in Eq. \ref{eq:rho} possesses inherent theoretical limitations. The SPVAT Lagrangian is constructed essentially as an ad hoc, bottom-up classification of the Lorentz-invariant four-fermion interactions. It does not systematically incorporate the underlying gauge structure of the SM, nor does it provide a framework for consistently evolving the couplings up to the high-energy matching scale. In order to systematically probe for new physics within muon decay, a more modern theoretical architecture is required: Low-Energy Effective Field Theory (LEFT).

\section{Experimental Michel-parameter likelihood}
\label{sec:michel-likelihood}

The experimental Michel-parameter inputs are not treated as a set of
independent one-dimensional Gaussian measurements. Instead, we divide
them into experimental likelihood blocks:
\begin{equation}
  \bm{O}_{\mathrm{fit}}
  =
  \underbrace{
    \bigl(\rho,\delta,P_\mu^\pi\xi\bigr)
  }_{\mathrm{TWIST}}
  \oplus
  \underbrace{
    \left(\eta,\eta'',\alpha'/A,\beta'/A\right)
  }_{\mathrm{Danneberg}}
  \oplus
  \underbrace{\xi'}_{\mathrm{Burkard}}
  \oplus
  \underbrace{\xi''}_{\mathrm{Prieels}}.
  \label{eq:michel-fit-vector}
\end{equation}
The different experimental blocks are assumed to be mutually
independent, whereas the correlations within each block are retained.

\paragraph{TWIST block}
For the TWIST measurement, we use the following native parameter basis \cite{Hillairet2012precision,bueno2011precise}:
\begin{equation}
  \bm{O}^{\mathrm{exp}}_{\mathrm{TWIST}}
  =
  \begin{pmatrix}
    \rho\\
    \delta\\
    P_\mu^\pi\xi
  \end{pmatrix}_{\!\mathrm{exp}}
  =
  \begin{pmatrix}
    0.74977\\
    0.75049\\
    1.00084
  \end{pmatrix}
  \label{eq:twist-input}
\end{equation}
with uncertainties
\begin{align}
  \rho
  &=
  0.74977
  \pm0.00012_{\mathrm{stat}}
  \pm0.00023_{\mathrm{syst}}
  \\
  \delta
  &=
  0.75049
  \pm0.00021_{\mathrm{stat}}
  \pm0.00027_{\mathrm{syst}}
  \\
  P_\mu^\pi\xi
  &=
  1.00084
  \pm0.00029_{\mathrm{stat}}
  {}^{+0.00165}_{-0.00063}{}_{\mathrm{syst}}
\end{align}
The reported correlation matrix, in the ordering
$\left(\rho,\delta,P_\mu^\pi\xi\right)$, is
\begin{equation}
  \bm{R}_{\mathrm{TWIST}}
  =
  \begin{pmatrix}
    1    & 0.19  & 0.21\\
    0.19 & 1     & -0.72\\
    0.21 & -0.72 & 1
  \end{pmatrix}
  \label{eq:twist-correlation}
\end{equation}
In a symmetrized Gaussian approximation, the covariance matrix is
constructed as:
\begin{equation}
  \bm{V}_{\mathrm{TWIST}}
  =
  \bm{D}_{\mathrm{TWIST}}\,
  \bm{R}_{\mathrm{TWIST}}\,
  \bm{D}_{\mathrm{TWIST}}
  \label{eq:twist-covariance}
\end{equation}
where
\begin{equation}
  \bm{D}_{\mathrm{TWIST}}
  =
  \operatorname{diag}
  \left(
    \sqrt{(0.00012)^2+(0.00023)^2},
    \sqrt{(0.00021)^2+(0.00027)^2},
    \sigma_{P_\mu^\pi\xi}
  \right).
\end{equation}
For a symmetric treatment of the uncertainty on $P_\mu^\pi\xi$, we symmetrize in quadrature:
\begin{equation}
  \sigma_{P_\mu^\pi\xi}
  \simeq
  \sqrt{
    (0.00029)^2+
    \left(\frac{0.00165+0.00063}{2}\right)^2
  }
  \label{eq:twist-xi-error}
\end{equation}

\paragraph{Danneberg block}
For the transverse-polarization measurement, we use the following general
four-parameter fit \cite{danneberg2005muon}:
\begin{equation}
  \bm{O}^{\mathrm{exp}}_{\mathrm{D}}
  =
  \begin{pmatrix}
    \eta\\
    \eta''\\
    \alpha'/A\\
    \beta'/A
  \end{pmatrix}_{\!\mathrm{exp}}
  =
  \begin{pmatrix}
     0.071\\
     0.105\\
    -0.0034\\
    -0.0005
  \end{pmatrix}
  \label{eq:danneberg-input}
\end{equation}
After combining statistical and systematic uncertainties in quadrature,
the standard-deviation matrix is:
\begin{equation}
  \bm{D}_{\mathrm{D}}
  =
  \operatorname{diag}
  \left(
    0.03734,\,
    0.05235,\,
    0.02186,\,
    0.00800
  \right)
  \label{eq:danneberg-errors}
\end{equation}
The corresponding correlation matrix, in the ordering
$\left(\eta,\eta'',\alpha'/A,\beta'/A\right)$, is:
\begin{equation}
  \bm{R}_{\mathrm{D}}
  =
  \begin{pmatrix}
    1     & 0.946 & 0      & 0\\
    0.946 & 1     & 0      & 0\\
    0     & 0     & 1      & -0.893\\
    0     & 0     & -0.893 & 1
  \end{pmatrix}
  \label{eq:danneberg-correlation}
\end{equation}
The covariance matrix is therefore:
\begin{equation}
  \bm{V}_{\mathrm{D}}
  =
  \bm{D}_{\mathrm{D}}\,
  \bm{R}_{\mathrm{D}}\,
  \bm{D}_{\mathrm{D}}
\end{equation}
Or numerically:
\begin{equation}
  \bm{V}_{\mathrm{D}}
  \simeq
  \begin{pmatrix}
    1.394\times10^{-3} & 1.849\times10^{-3} & 0 & 0\\
    1.849\times10^{-3} & 2.741\times10^{-3} & 0 & 0\\
    0 & 0 & 4.779\times10^{-4} & -1.562\times10^{-4}\\
    0 & 0 & -1.562\times10^{-4} & 6.400\times10^{-5}
  \end{pmatrix}
  \label{eq:danneberg-covariance}
\end{equation}
Consequently, $\eta$ and $\eta''$ must not be treated as statistically
independent and neither $\alpha'/A$ and $\beta'/A$ -- indeed both pairs are separately strongly correlated. 

\paragraph{Single-observable blocks}
The remaining polarization observables are included as independent
one-dimensional blocks \cite{burkard1985muon,prieels2014measurement}:
\begin{align}
  \xi'  &= 0.998\pm0.045,\\
  \xi'' &= 0.981
           \pm0.045_{\mathrm{stat}}
           \pm0.003_{\mathrm{syst}}
         \simeq 0.981\pm0.0451
\end{align}
No covariance between these measurements and the TWIST or Danneberg
datasets is available.
\paragraph{Combined likelihood}
The total experimental covariance is approximated by:
\begin{equation}
  \bm{V}_{\mathrm{exp}}
  \simeq
  \bm{V}_{\mathrm{TWIST}}
  \oplus
  \bm{V}_{\mathrm{D}}
  \oplus
  \sigma_{\xi'}^2
  \oplus
  \sigma_{\xi''}^2
  \label{eq:combined-michel-covariance}
\end{equation}
For a vector of LEFT Wilson coefficients $\bm{L}$, the corresponding
contribution to the likelihood is:
\begin{align}
  \chi^2_{\mathrm{Michel}}(\bm{L})
  =
  &\sum_b
  \left[
    \bm{O}^{\mathrm{th}}_b(\bm{L})
    -
    \bm{O}^{\mathrm{exp}}_b
  \right]^T
  \bm{V}_b^{-1}
  \left[
    \bm{O}^{\mathrm{th}}_b(\bm{L})
    -
    \bm{O}^{\mathrm{exp}}_b
  \right]
  \nonumber\\
  &b\in
  \left\{
    \mathrm{TWIST},
    \mathrm{D},
    \xi',
    \xi''
  \right\}
  \label{eq:michel-chi2}
\end{align}
The quantities involving $P_\mu^\pi$ are measured as polarization products. Setting $P_\mu^\pi=1$ is therefore an additional
phenomenological assumption rather than part of the experimental
measurement. Furthermore, the Michel shape and polarization parameters do not constrain the overall normalization of the muon-decay amplitude. If normalization-changing LEFT coefficients are included, the muon
lifetime or an equivalent normalization constraint must also be used. The full set of experimental-level Michel parameters entering the likelihood is given in Tab.~\ref{tab:michel-likelihood-inputs}.
\begin{table}[tbp]
    \centering
    \begin{tabular}{llll}
        \toprule
        Observable
        & Experimental value
        & SM value
        & Source
        \\
        \midrule
        $\rho$
        & $0.74977\pm0.00026$
        & $3/4$
        & TWIST \cite{Hillairet2012precision}
        \\
        $\delta$
        & $0.75049\pm0.00034$
        & $3/4$
        & TWIST \cite{Hillairet2012precision}
        \\
        $P_\mu^\pi\xi$
        & $1.00084^{+0.00168}_{-0.00069}$
        & $1$
        & TWIST \cite{bueno2011precise}
        \\
        $\eta$
        & $0.071\pm0.0373$
        & $0$
        & Danneberg et al.\ \parencite{danneberg2005muon}
        \\
        $\eta''$
        & $0.105\pm0.0523$
        & $0$
        & Danneberg et al.\ \parencite{danneberg2005muon}
        \\
        $\alpha'/A$
        & $-0.0034\pm0.0219$
        & $0$
        & Danneberg et al.\ \parencite{danneberg2005muon}
        \\
        $\beta'/A$
        & $-0.0005\pm0.0080$
        & $0$
        & Danneberg et al.\ \parencite{danneberg2005muon}
        \\
        $\xi'$
        & $0.998\pm0.045$
        & $1$
        & Burkard et al.\ \parencite{burkard1985muon}
        \\
        $\xi''$
        & $0.981\pm0.0451$
        & $1$
        & Prieels et al.\ \parencite{prieels2014measurement}
        \\
        \bottomrule
    \end{tabular}
    \caption{%
        Experiment-level Michel-parameter inputs entering the likelihood.
        Statistical and systematic uncertainties are combined in quadrature,
        except that the asymmetric uncertainty on $P_\mu^\pi\xi$ is retained.
        These values should be distinguished from the PDG averages quoted
        elsewhere in the text.
    }
    \label{tab:michel-likelihood-inputs}
\end{table}
\chapter{Low Energy Effective Field Theory (LEFT)}\label{ch:left}
The LEFT is constructed by integrating out all heavy degrees of freedom at the electroweak scale (the $W$, $Z$ and Higgs boson as well as the top quark) while strictly preserving the unbroken low-energy gauge symmetry of the Standard Model, $SU(3)_C \times U(1)_{em}$. This provides several critical advantages over the historical SPVAT approach:
\begin{itemize}
    \item \textbf{Gauge Invariance:} The operators in LEFT are explicitly constructed to be invariant under QED and QCD, ensuring theoretical consistency at low energies.
    \item \textbf{Systematic Expansion:} LEFT allows for a rigorous power-counting expansion in terms of operator dimensionality (suppressed by powers of the new physics scale $\Lambda$), guaranteeing that all possible interactions up to a given mass dimension are accounted for.
    \item \textbf{Renormalization Group (RG) Evolution:} A formal EFT framework allows the Wilson coefficients evaluated at the high-energy matching scale ($\mu \sim m_W$) to be systematically evolved down to the kinematic scale of muon decay ($\mu \sim m_\mu$) via the renormalization group equations, properly accounting for operator mixing.
\end{itemize}

By mapping the specific SPVAT couplings to the rigorously defined Wilson coefficients of the LEFT operator basis, we can modernize the historical Michel parameters. The primary objective of the subsequent chapters is to execute this exact mapping: deriving the analytical deviations of the Michel parameters driven by the insertion of dimension-six LEFT operators, thereby establishing a precise, gauge-invariant dictionary between high-energy BSM theory and low-energy muon decay observables.



At the energy scale of muon decay, determined by the muon mass $m_\mu \approx 106\text{ MeV}$, the heavy fields of the SM are not dynamical degrees of freedom - the $W^\pm$, $Z$ and $H$ boson as well as the $t$ quark all have masses clustered around the electroweak scale $v\approx 246\text{ GeV}$, three orders of magnitude above the muon mass. For energies well below this scale, an effective theory can be derived by integrating out the heavy fields and adding higher-dimensional effective operators built from the remaining field content. 

\section{The LEFT basis}
\paragraph{LEFT symmetries}
Since the particles responsible for electroweak symmetry breaking are integrated out, only the unbroken gauge symmetry remains:
\begin{equation}
    SU(3)_C \times U(1)_\text{em}
\end{equation}
This is conceptually different from the broken phase of the SM:  While left-handed fermions only exist as doublets in the SM, regardless of the energy scale, the LEFT treats left-handed neutrinos and left-handed charged leptons as completely independent, unrelated fields. As a consequence, left- and right-handed particles can mix more freely in the LEFT and new mathematical structures like scalar and tensor interactions arise.
\paragraph{LEFT particle content}
As degrees of freedom remain the massless bosons and all fermions except the top quark.
\begin{itemize}
    \item \textbf{Gauge bosons} Photon $\gamma$ and gluons $g$
    \item \textbf{Charged fermions} Charged leptons $e,\mu,\tau$ and the five lightest quarks $u,d,c,s,b$. Since the $SU(2)_L$ group is missing, left- and right-handed charged fermions have the same gauge charges and as they are moreover coupled together in the Lagrangian by explicit mass terms after symmetry breaking they are written as single Dirac spinors, e.g. $e=e_L+e_R$.
    \item \textbf{Neutral fermions} For the neutrinos, however, standard LEFT only contains left-handed fields $\nu_e,\nu_\mu,\nu_\tau$.
\end{itemize}
\paragraph{LEFT Lagrangian}
The renormalizable parts of QCD and QED plus an infinite tower of higher-dimensional operators built from the light fields form the LEFT Lagrangian:
\begin{equation}
    \mathcal{L}_{\text{LEFT}} = \mathcal{L}_{\text{QED} + \text{QCD}} + \sum_{d \ge 5} \sum_i L_i^{(d)} \mathcal{O}_i^{(d)}
\end{equation}
\begin{itemize}
    \item $\mathcal{O}_i^{(d)}$ are the local, gauge-invariant operators of mass dimension $d$
    \item $L_i^{(d)}$ are the Wilson coefficients
\end{itemize}
The number of operators at each mass dimension is given in table \ref{tab:LEFTnum} \parencite{jenkins2018lowOM}. 
\begin{table}[t]
    \centering
    \begin{tabular}{cccc} \toprule
             d&quantum numbers&$n_g=1$&$n_g=3$ \\\midrule
             3&$(\Delta L=2$) + h.c.&2&12\\
             5&$\Delta L=\Delta B=0$&10&70\\
             5&$(\Delta L=2)$ + h.c.&0&6\\
             6&$\Delta B=\Delta L=0$&80&3631\\
             6&$(\Delta L=2)$ + h.c.&22&1200\\
             6&$(\Delta L=4)$ + h.c.&0&12\\
             6&$(\Delta B=\Delta L=1)$ + h.c.&12&576\\
             6&$(\Delta B=-\Delta L=1)$ + h.c.&4&456\\\bottomrule
    \end{tabular}
    \caption{Number of LEFT operators with dimension 3,5 and 6 ($n_g=$ generations)}
    \label{tab:LEFTnum}
\end{table}
For muon decay, the following operators are relevant:\par
Dimension-five dipole operators:
\begin{align}
    \op_{eA}&=\bar e_{L,p}\sigma^{\mu\nu}e_{R,r}F_{\mu\nu}\\
    \op_{\nu A}&=(\nu^T_{L,p}C\sigma^{\mu\nu}e_{R,r})F_{\mu\nu}
\end{align}
Dimension-six operators:
\begin{align}
    \op_{\nu e}^{V,LL}&=(\bar\nu_{L,p}\gamma^\mu\nu_{L,r})(\bar e_{L,s}\gamma_\mu e_{L,t})\\
    \op_{\nu e}^{V,LR}&=(\bar\nu_{L,p}\gamma^\mu\nu_{L,r})(\bar e_{R,s}\gamma_\mu e_{R,t})\\
    \op_{\nu e}^{S,LL}&=(\nu_{L,p}^TC\nu_{L,r})(\bar e_{R,s} e_{L,t})\\
    \op_{\nu e}^{S,LR}&=(\nu_{L,p}^TC\nu_{L,r})(\bar e_{L,s} e_{R,t})\\
    \op_{\nu e}^{T,LL}&=(\nu_{L,p}^TC\sigma^{\mu\nu}\nu_{L,r})(\bar e_{R,s}\sigma_{\mu\nu} e_{L,t})
\end{align}
\begin{itemize}
    \item[$\bullet$] Flavor-changing coupling of leptons to photon highly supressed
    \item[$\bullet$] $C$ is charge conjugation matrix
    \item[$\bullet$] $p,r,s,t$ flavor indices
    \item[$\bullet$] Only $\op_{\nu e}^{V,LL}$ and $\op_{\nu e}^{V,LR}$ preserve lepton number
\end{itemize}
\section{Michel parameters in the LEFT}
\subsection{SM-like neutrino output}

We begin by deriving the Michel parameters for $\mu^-\to e^-\bar\nu_e\nu_\mu$, which corresponds to the SM process. We follow the discussion from \ref{sec:michel} and further elaborated in \ref{sec:kinPS}. We consider strictly the SM-like process as a simple first look at our results. In the next subsection we consider all four-fermion operators simultaneously which allows for different neutrino final states.

Starting with the contributions insensitive to the electron polarization we find:
%
\begin{align}
\frac{G_F^2}{G_{F,{\rm SM}}^2}=&\;1-\frac{v^2}{2}\left[L_{\nu e,2112}^{V,LL}+L_{\nu e,1221}^{V,LL}\right]+\frac{v^4}{4}\left[L_{\nu e,2112}^{V,LL}L_{\nu e,1221}^{V,LL}+L_{\nu e,2112}^{V,LR}L_{\nu e,1221}^{V,LR}\right]\label{eq:SMlikeGF}\\[5pt]
\rho=&\;\frac{3}{4}\\[5pt]
\eta=&\;\frac{v^2}{4}\left[L_{\nu e,2112}^{V,LR}+L_{\nu e,1221}^{V,LR}\right]+\frac{v^4}{8}\left[L_{\nu e,2112}^{V,LL}L_{\nu e,1221}^{V,LR}+L_{\nu e,2112}^{V,LL}L_{\nu e,1221}^{V,LR}\right]\\[5pt]
\xi=&\;1-\frac{v^4}{2}L_{\nu e,2112}^{V,LR}L_{\nu e,1221}^{V,LR}\\[5pt]
\delta=&\frac{3}{4}
\end{align}
Since the operators $\op_{\nu e}^{V,LL}$ and $\op_{\nu e}^{V,LR}$ are hermitian, their flavor structures $1221$ and $2112$ are equivalent and for their coefficients it holds that ${L_{\nu e,2112}^{V,LL}}^*=L_{\nu e,1221}^{V,LL}$ and ${L_{\nu e,2112}^{V,LR}}^*=L_{\nu e,1221}^{V,LR}$ so we can simplify the first order corrections for $G_F/G_{F,\rm SM}$ to the real part of one operator coefficient.\par 
In many analyses $G_F$ is viewed as an input parameter. However, we use the result of \cite{grunwald2025beyond} where the authors apply the $\{\hat m_W,\hat m_Z,\hat \alpha\}$ input parameter scheme to conduct a theory determination of $G_F$:
\begin{equation}
G_{F,{\rm SM}}=\frac{\pi\hat \alpha\hat m_Z^2}{\sqrt{2}\hat m_W^2(\hat m_Z^2-\hat m_W^2)}(1+\delta_R) =1.1658(10)\cdot 10^{-5}\, 
\end{equation}
For this they use the values $\delta_R=0.03686(21)$, $\hat\alpha=7.2973525693(11) \cdot10^{-3}$, $m_W=80.369(13)$, and $m_Z=91.1880(20)$ from \cite{particle2024review}. Comparing this SM prediction with the leading contribution in \ref{eq:SMlikeGF} and the experimental value in Eq.~\ref{eq:GFexp}, we find
%
\begin{equation}
-0.004\le v^2\text{Re}\left[L_{\nu e,2112}^{V,LL}\right]\le 0.002
\end{equation}
%
By taking $|L|\sim 1/\Lambda^2$ (and assuming that $|L|\approx\text{Re}[L]$) we obtain a 95\% confidence level bound on the scale of new physics:
%
\begin{equation}
\Lambda\ge 21v\sim {\rm\ 5.1\ TeV}
\end{equation}
If we allow for contributions that are quadratic in $L_{\nu e,2112}^{V,LR}$ and  $L_{\nu e,2112}^{V,LL}$ we can define a $\chi^2$ over $G_F, \eta$ and $\xi$. Assuming $P_\mu=1$, we find the following 95\% confidence level region for the two contributing Wilson coefficients:
\begin{align}
-0.036\le\quad v^2\, L_{\nu e,2112}^{V,LR}\quad\le0.064\, ,\\
-0.0042\le\quad  v^2\, L_{\nu  e,2112}^{V,LL}\quad\le 0.0025
\end{align}

Extending our analysis to include electron polarization dependence we find corrections to the following additional Michel parameters:
%
\begin{align}
\eta''=&-\frac{v^2}{2}c^{V,LR}_{\nu e,2112}-\frac{v^4}{4}L_{\nu e,2113}^{V,LL}L_{\nu e,2113}^{V,LR}=-\eta\\[5pt]
\xi'=&\;\;1-\frac{v^4}{2}L_{\nu e,2112}^{V,LR}L_{\nu e,1221}^{V,LR}=\xi\\[5pt]
\frac{\beta'}{A}=&-\frac{v^2}{4}{\rm Im}[L_{\nu e,2112}^{V,LR}]-\frac{v^4}{8}{\rm Im}[(L_{\nu e,2112}^{V,LL})^*L_{\nu e,2112}^{V,LR}]\nonumber\\[5pt]
&-\frac{v^4}{8}{\rm Im}[(L_{\nu e,2112}^{V,LL})L_{\nu e,2112}^{V,LR}]
\end{align}
%

The parameters $\xi''$ and $\alpha'/A$ remain unchanged -- in the LEFT we will see they only receive corrections from lepton number violating operators. As $\eta''=-\eta$ and $\xi'=\xi$ and these two additional parameters are more poorly measured than $\eta$ and $\xi$, we neglect to extend our fit to include them here. However, the leading interference for $\beta'/A$ together with the one for $\eta$ allow for determining the real and imaginary component of $L^{V,LR}_{\nu e,2112}$:
\begin{equation}
    \eta\approx\frac{1}{2}\text{Re}[L^{V,LR}_{\nu e,2112}] \quad\quad \frac{\beta'}{A} \approx  \frac{1}{4} \text{Im}[L^{V,LR}_{\nu e,2112}]
    \label{eq:etabetaprime}
\end{equation}
This operator can be transformed to to a scalar operator by a Fierz transformation:
\begin{align}
    \mathcal{O}_{\nu e,prst}^{V,LR}=&\;(\bar{\nu}_{p}\gamma^\mu P_L\nu_{r})(\bar e_{s}\gamma_\mu P_Re_{t})\nonumber\\[5pt]
    =&\;2(\bar e_{Rs}\nu_{Lr})(\bar{\nu}_{Lp}e_{Rt})\nonumber\\[8pt]
    &\text{by Fierz: } (\gamma^\mu P_L)_{ab}(\gamma_\mu P_R)_{cd}=-2(P_R)_{ad}(P_L)_{cb} \label{eq:fierzingV,LR2SRR}
\end{align}
This resembles for $prst=2112$ the Michel structure of $g_{RR}^S$ in the notation by Fetscher, Gerber \& Johnson \cite{fetscher1986muon} and thus our determination in \ref{eq:etabetaprime} reproduces the findings in~\cite{falkowski2016model} where the authors state that at leading order (and for SM like neutrino output) only the Michel parameters $\eta$ and $\beta'/A$ receive a correction stemming from the scalar interaction term $g_{RR}^S$. Since $L^{V,LR}_{\nu e,2112}$ and $g_{RR}^S$ produce the same operator structure we can also constrain $L^{V,LR}_{\nu e,2112}$ directly from the PDG results for the SPVAT matrix elements~\cite{particle2024review}:
\begin{equation}
    L^{V,LR}_{\nu e,2112}\simeq |g_{RR}^S|<0.035
\end{equation}
Before we are discussing the evolution of this constrain up to the EW scale we discuss the case of full neutrino output.
\subsection{Full four-fermion results}
Starting with neglecting measurements sensitive to electron spin, we find for $G_F$:
%
\begin{align}
G_F^2/G_{F,{\rm SM}}^2=&\;1-\frac{v^2}{2}\left[L_{\nu e,2112}^{V,LL}+L_{\nu e,1221}^{V,LL}\right]\nonumber\\[5pt]
&+\frac{v^4}{4}\sum_{a,b=1}^3\left[L_{\nu e,ab12}^{V,LL}L_{\nu e,ba21}^{V,LL}+L_{\nu e,ab12}^{V,LR}L_{\nu e,ba21}^{V,LR}\right]\nonumber\\[5pt]
&+\frac{v^4}{16}\sum_{a<b}\left[|L_{\nu e,ab12}^{S,LL}+L_{\nu e,ba12}^{S,LL}|^2+|L_{\nu e,ab12}^{S,LR}+L_{\nu e,ba12}^{S,LR}|^2+(12\to 21)\right]\nonumber\\[5pt]
&+2\times\frac{v^4}{16} \sum_{a=1}^{3}\left[|L_{\nu e,aa12}^{S,LL}|^2+|L_{\nu e,aa12}^{S,LR}|^2+(12\to 21)\right]\nonumber\\[5pt]
&+3 v^4\sum_{a<b}\left[|L_{\nu e,ab12}^{T,LL}-L_{\nu e,ba12}^{T,LL}|^2+(12\to 21)\right]
\end{align}
%
Here we have used ``$(12\to 21)$'' to mean swapping only the last two flavor indices from ``12'' to ``21.'' For sums with $a<b$ it is understood $b$ is summed from 1 to 3. The factor of two for like neutrinos arises from the counting factor (1/2) for identical final-state particles and symmetrizing the amplitude giving $|c+c|^2$ to $4|c|^2$ for $a=b$. The factor of 3 for the tensor can be understood as coming from the six independent components of $\sigma^{\mu\nu}$ and the trace over dirac matrices. The difference in sign structure between scalar and tensor terms is simply due to the symmetry properties of the two operators under interchange of the neutrino flavor indices.

For the other electron spin independent Michel parameters we find:
%
\begin{align}
\rho=&\;\;\frac{3}{4}-\frac{3v^4}{64}\sum_{a<b}\left[|L_{\nu e,ab12}^{S,LL}+L_{\nu e,ba12}^{S,LL}|^2+|L_{\nu e,ab12}^{S,LR}+L_{\nu e,ba12}^{S,LR}|^2+(12\to 21)\right]\nonumber\\[5pt]
&\;\;\phantom{\frac{3}{4}}-2\times \frac{3v^4}{64}\sum_{a=1}^3\left[|L_{\nu e,aa12}^{S,LL}|^2+|L_{\nu e,aa12}^{S,LR}|^2+(12\to 21)\right]\nonumber\\[5pt]
&\;\;\phantom{\frac{3}{4}}+\frac{3v^4}{4}\sum_{a<b}\left[|L_{\nu e,ab12}^{T,LL}-L_{\nu e,ba12}^{T,LL}|^2+(12\to 21)\right]\\[5pt]
\eta=&\;\;0+\frac{v^2}{4}\left[L_{\nu e,2112}^{V,LR}+L_{\nu e,1221}^{V,LR}\right]+\frac{v^4}{8}\left[L_{\nu e,1221}^{V,LL}L_{\nu e,2112}^{V,LR}+L_{\nu e,2112}^{V,LL}L_{\nu e,1221}^{V,LR}\right]\nonumber\\[5pt]
&\;\;\phantom{0}+\frac{v^4}{8}\left[L_{\nu e,2112}^{V,LL}L_{\nu e,2112}^{V,LR}+L_{\nu e,1221}^{V,LL}L_{\nu e,1221}^{V,LR}\right]\nonumber\\[5pt]
&\;\;\phantom{0}-\frac{v^4}{8}\sum_{a,b=1}^3\left[L_{\nu e,ba12}^{V,LL}L_{\nu e,ab12}^{V,LR}+L_{\nu e,ab12}^{V,LL}L_{\nu e,ba21}^{V,LR}\right]\nonumber\\[5pt]
&\;\;\phantom{0}+\frac{v^4}{8}\sum_{a<b}{\rm Re}\left[(L_{\nu e,ab12}^{S,LL}+L_{\nu e,ba12}^{S,LL})^*(L_{\nu e,ab12}^{S,LR}+L_{\nu e,ba12}^{S,LR})+(12\to 21)\right]\nonumber \\[5pt]
&\;\;\phantom{0}+\frac{v^4}{4}\sum_{a=1}^3{\rm Re}\left[(L_{\nu e,aa12}^{S,LL})^*(L_{\nu e,aa12}^{S,LR})+(12\to 21)\right]\\[5pt]
\xi=&\;\;1-\frac{v^4}{2}\sum_{a,b=1}^3L_{\nu e,ab12}^{V,LR}L_{\nu e,ba21}^{V,LR}\nonumber\\[5pt]
&\;\;\phantom{1}+\frac{v^4}{8}\sum_{a<b}\left[|L_{\nu e,ab12}^{S,LR}+L_{\nu e,ba21}^{S,LR}|^2-2|L_{\nu e,ab12}^{S,LL}+L_{\nu e,ba12}^{S,LL}|^2\right]\nonumber\\[5pt]
&\;\;\phantom{1}+\frac{v^4}{8}\sum_{a<b}\left[-2|L_{\nu e,ab21}^{S,LR}+L_{\nu e,ba21}^{S,LR}|^2+|L_{\nu e,ab21}^{S,LL}+L_{\nu e,ba21}^{S,LL}|^2\right]\nonumber\\[5pt]
&\;\;\phantom{1}+\frac{v^4}{4}\sum_{a=1}^3\left[|L_{\nu e,aa21}^{S,LL}|^2-2|L_{\nu e,aa21}^{S,LR}|^2-2|L_{\nu e,aa12}^{S,LL}|^2+|L_{\nu e,aa12}^{S,LR}|^2\right]\nonumber\\[5pt]
&\;\;\phantom{1}+\frac{v^4}{8}\sum_{a<b}\left[32|L_{\nu e,ab12}^{T,LL}-L_{\nu e,ba12}^{T,LL}|^2-80|L_{\nu e,ab21}^{T,LL}-L_{\nu e,ba21}^{T,LL}|^2\right]
\end{align}
%
\begin{align}
\delta=&\;\;\frac{3}{4}+\frac{9v^4}{64}\sum_{a<b}\left[|L_{\nu e,ab12}^{S,LL}+L_{\nu e,ba12}^{S,LL}|^2-|L_{\nu e,ab12}^{S,LR}+L_{\nu e,ba12}^{S,LR}|^2 -(12\to 21)\right]\nonumber\\[5pt]
&\;\;\phantom{\frac{3}{4}}+\frac{9v^4}{32}\sum_{a=1}^3\left[|L_{\nu e,aa12}^{S,LL}|^2-|L_{\nu e,aa12}^{S,LR}|^2-(12\to 21)\right]\nonumber\\[5pt]
&\;\;\phantom{\frac{3}{4}}-\frac{9v^4}{4}\sum_{a<b}\left[|L_{\nu e,ab12}^{T,LL}-L_{\nu e,ba12}^{T,LL}|^2-(12\to 21)\right]
\end{align}
%
The extra quadratic term in $\eta$ that does not appear as part of a sum arises from the cross term when expanding the ratio defining eta (see \ref{eq:solveforMichel}). It originates from the product of the linear term in the numerator ($\sim L_{\nu e,2112}^{V,LR}$) with the linear correction to $1/G_F^2$ ($\sim L_{\nu e,2112}^{V,LL}$). This contribution only exists for the SM flavor assignment. Notice that for $\xi$ there is a difference between the ``$12$'' and ``$21$'' contributions. From \ref{eq:LLvio} we can see that exchanging these labels swaps the chirality assigned to the electron and muon. The parameter $\xi$, which corresponds to the anisotropic part of the rate, is sensitive to precisely this difference in chiral assignments. For the same reason, $\delta$ also swaps signs under the flavor exchange $(12\to 21)$.


Measurements sensitive to the electron's polarization then access the additional Michel parameters:
%
\begin{align}
\xi'=&\;\;1-\frac{v^4}{2}\sum_{a,b=1}^3L_{\nu e,ab12}^{V,LR}L_{\nu e,ba21}^{V,LR}\nonumber\\[5pt]
&\;\;\phantom{1}-\frac{v^4}{8}\sum_{a<b}\left[|L_{\nu e,ab12}^{S,LL}+L_{\nu e,ba12}^{S,LL}|^2+|L_{\nu e,ab21}^{S,LR}+L_{\nu e,ba21}^{S,LR}|^2\right]\nonumber\\[5pt]
&\;\;\phantom{1}-\frac{v^4}{4}\sum_{a=1}^3\left[|L_{\nu e,aa12}^{S,LL}|^2+|L_{\nu e,aa21}^{S,LR}|^2\right]-6v^4\sum_{a<b}|L_{\nu e,ab12}^{T,LL}-L_{\nu e,ba12}^{T,LL}|^2\\[5pt]
\xi''=&\;\;1+\frac{v^4}{8}\sum_{a<b}\left[|L_{\nu e,ab12}^{S,LL}+L_{\nu e,ba12}^{S,LL}|^2+|L_{\nu e,ab12}^{S,LR}+L_{\nu e,ba12}^{S,LR}|^2+(12\to 21)\right]\nonumber\\[5pt]
&\;\;\phantom{1}+\frac{v^4}{4}\sum_{a=1}^3\left[|L_{\nu e,aa12}^{S,LL}|^2+|L_{\nu e,aa12}^{S,LR}|^2+(12\to 21)\right]\nonumber\\[5pt]
&\;\;\phantom{1}-10v^4\sum_{a<b}\left[|L_{\nu e,ab12}^{T,LL}-L_{\nu e,ba12}^{T,LL}|^2+(12\to 21)\right]
\end{align}

\begin{align}
\eta''=&\;\;0-\frac{v^2}{4}\left[L_{\nu e,2112}^{V,LR}+L_{\nu e,1221}^{V,LR}\right]-\frac{v^4}{8}[L_{\nu e,1221}^{V,LL}L_{\nu e,2112}^{V,LR}+L_{\nu e,2112}^{V,LL}L_{\nu e,1221}^{V,LR}]\nonumber\\[5pt]
&\;\;\phantom{0}-\frac{v^4}{8}[L_{\nu e,2112}^{V,LL}L_{\nu e,2112}^{V,LR}+L_{\nu e,1221}^{V,LL}L_{\nu e,1221}^{V,LR}]\nonumber\\[5pt]
&\;\;\phantom{0}+\frac{v^4}{8}\sum_{a,b=1}^3[L_{\nu e,ba21}^{V,LL}L_{\nu e,ab12}^{V,LR}+L_{\nu e,ab12}^{V,LL}L_{\nu e,ba21}^{V,LR}]\nonumber\\[5pt]
&\;\;\phantom{0}+\frac{3v^4}{8}\sum_{a<b}{\rm Re}\left[(L_{\nu e,ab12}^{S,LL}+L_{\nu e,ba12}^{S,LL})^*(L_{\nu e,ab12}^{S,LR}+L_{\nu e,ba12}^{S,LR})+(12\to 21)\right]\nonumber\\[5pt]
&\;\;\phantom{0}+\frac{3v^4}{4}\sum_{a=1}^3{\rm Re}\left[(L_{\nu e,aa12}^{S,LL})^*L_{\nu e,aa12}^{S,LR}+(12\to 21)\right]\\[5pt]
\frac{\alpha'}{A}=&\;\;0-\frac{v^4}{8}\sum_{a<b}{\rm Im}[(L_{\nu e,ab12}^{S,LL}+L_{\nu e,ba12}^{S,LL})^*(L_{\nu e,ab12}^{S,LR}+L_{\nu e,ba12}^{S,LR})+(12\to 21)]\nonumber\\[5pt]
&\;\;\phantom{0}-\frac{v^4}{4}\sum_{a=1}^3{\rm Im}[(L_{\nu e,aa12}^{S,LL})^*L_{\nu e,aa12}^{S,LR}+(12\to 21)]\\[5pt]
%%
\frac{\beta'}{A}=&\;\;0-\frac{v^2}{4}{\rm Im}[L_{\nu e,2112}^{V,LR}]-\frac{v^4}{8}{\rm Im}[(L_{\nu e,2112}^{V,LL})^*L_{\nu e,2112}^{V,LR}]-\frac{v^4}{8}{\rm Im}[(L_{\nu e,2112}^{V,LL})L_{\nu e,2112}^{V,LR}]\nonumber\\[5pt]
&\;\;\phantom{0}+\frac{v^4}{8}\sum_{a,b=1}^3{\rm Im}[(L_{\nu e,ab12}^{V,LL})^*L_{\nu e,ab12}^{V,LR}]
\end{align}
\section{Constraints on the LEFT Wilson coefficients}\label{sec:left-constraints}
Having expressed the Michel parameters as functions of the LEFT Wilson coefficients, we can now invert the relation: given the experimental likelihood of Sec.~\ref{sec:michel-likelihood} we constrain the coefficients themselves. Throughout we work with the dimensionless combination $v^2L^{(6)}_{\nu e}\sim(v/\Lambda)^2$, so that a bound on $v^2L$ translates, for an order-one UV coupling, into a lower bound on the new-physics scale via $\Lambda\gtrsim v/\sqrt{|v^2L|}$. The results quoted here are obtained at the scale of muon decay, $\mu\sim m_\mu$; their evolution to the electroweak scale is the subject of the following section.

\paragraph{Interference versus quadratic sensitivity}
The structure of the full four-fermion results dictates two qualitatively different regimes. The lepton-number-conserving vector operators $\mathcal{O}^{V,LL}_{\nu e}$ and $\mathcal{O}^{V,LR}_{\nu e}$ with the Standard-Model flavour assignment ($ab=21$, i.e.\ the final state $\nu_\mu\bar\nu_e$) interfere with the SM amplitude and therefore enter the observables \emph{linearly}, at order $v^2L$. Every lepton-number-violating operator, as well as the vector operators for any non-SM neutrino flavour, cannot interfere with the SM final state and enters only \emph{quadratically}, as $|v^2L|^2$. Since the neutrinos are not observed, the observables depend only on incoherent sums over the neutrino flavours; a bound on an individual flavour coefficient is therefore only meaningful under the assumption of single-operator dominance, and within a given operator class the bound is largely independent of the specific neutrino-flavour label.

The two regimes produce reach that differs by orders of magnitude. A linearly-sensitive coefficient is bounded as $|v^2L|\lesssim\sigma_{\rm exp}$, while a quadratic one only as $|v^2L|\lesssim\sqrt{\sigma_{\rm exp}}$, so that the mass reach scales as $\Lambda\propto\sigma_{\rm exp}^{-1/4}$ rather than $\sigma_{\rm exp}^{-1/2}$. This is the origin of the multi-TeV bounds on the interfering coefficients and the $\mathcal{O}(1\text{ TeV})$ bounds on the lepton-number-violating ones.

\paragraph{Single-operator bounds}
Switching on one coefficient at a time and profiling the block-covariance $\chi^2$ of Eq.~\eqref{eq:michel-chi2} (supplemented by the muon-lifetime/$G_F$ normalization constraint, required once normalization-changing coefficients are included) yields the $95\%$ confidence-level bounds of Table~\ref{tab:leftbounds}. The coefficients listed are representatives spanning the distinct operator classes: the lepton-number-conserving vectors in the SM flavour slot, a vector in a non-SM flavour slot, and the scalar and tensor structures in both the diagonal ($a=b$) and off-diagonal ($a\neq b$) neutrino configurations.

\begin{table}[tbp]
    \centering
    \begin{tabular}{lccc}\toprule
        Coefficient & $95\%$ CL interval on $v^2L$ & $\Lambda$ [TeV] & sensitivity \\\midrule
        $L^{V,LL}_{\nu e,2112}$ (Re) & $[-0.0044,\,+0.0024]$ & $3.7$ & linear ($G_F$) \\
        $L^{V,LR}_{\nu e,2112}$ (Re) & $[-0.021,\,+0.018]$ & $1.7$ & linear ($\eta,\eta''$) \\
        $L^{V,LR}_{\nu e,2112}$ (Im) & $[-0.017,\,+0.023]$ & $1.6$ & linear ($\beta'/A$) \\
        $L^{V,LL}_{\nu e,3312}$ (Re) & $[-0.13,\,+0.13]$ & $0.7$ & quadratic \\
        $L^{S,LL}_{\nu e,1112}$ (Re) & $[+0.036,\,+0.069]$ & $0.9$ & quadratic \\
        $L^{S,LR}_{\nu e,1112}$ (Re) & $[-0.025,\,+0.025]$ & $1.6$ & quadratic \\
        $L^{S,LR}_{\nu e,1212}$ (Re) & $[-0.036,\,+0.036]$ & $1.3$ & quadratic \\
        $L^{T,LL}_{\nu e,1212}$ (Re) & $[-0.008,\,+0.008]$ & $2.8$ & quadratic \\\bottomrule
    \end{tabular}
    \caption{Single-operator $95\%$ CL bounds on selected dimensionless LEFT Wilson coefficients $v^2L^{(6)}_{\nu e}$ at the muon scale, from the Michel-parameter likelihood of Sec.~\ref{sec:michel-likelihood}. The implied scale assumes $|L|\sim1/\Lambda^2$ with an order-one coupling. Flavour indices follow the convention $prst$ with $p,r$ the neutrino and $s,t$ the charged-lepton generations.}
    \label{tab:leftbounds}
\end{table}
The interfering coefficient $L^{V,LL}_{\nu e,2112}$ is the most tightly constrained, reproducing the $\Lambda\gtrsim5$~TeV normalization bound; its imaginary vector partner $L^{V,LR}_{\nu e,2112}$ is bounded through the transverse-polarization parameter $\beta'/A$, giving access to a CP-violating phase. The tensor operator is comparatively well constrained even though it enters quadratically, because it feeds the precisely-measured shape parameters $\rho$, $\delta$ and $\xi$ with large numerical coefficients.

It is worth noting that the scalar coefficient $L^{S,LL}_{\nu e,1112}$ has a $95\%$ interval that excludes zero. This is not evidence for new physics: a diagonal scalar lowers $\rho$ and raises $\delta$, and the current central values sit slightly below and above their SM predictions respectively, so the fit mildly prefers a non-zero value. With $\chi^2_{\rm SM}\approx21$ for the ten observables -- driven largely by the internal tension of the TWIST block, where $\delta$ and $P_\mu^\pi\xi$ are both pulled upward despite their $-0.72$ correlation -- this preference reflects the present experimental inputs rather than a significant signal.

\paragraph{Global fit and flat directions}
Relaxing single-operator dominance and profiling one coefficient with the others floating weakens the bounds through cancellations: For instance the $\Lambda\gtrsim3.7$~TeV bound on $\text{Re}\,L^{V,LL}_{\nu e,2112}$ degrades to $\gtrsim2.6$~TeV once $\text{Re}\,L^{V,LR}_{\nu e,2112}$ is allowed to compensate. A singular-value decomposition of the (covariance-whitened) observable response at the SM point shows that only the interfering combinations receive linear sensitivity, while all lepton-number-violating directions are flat at linear order and bounded purely quadratically. Consequently the Michel parameters alone constrain only a limited number of coefficient combinations; the remaining directions require complementary observables. These findings are subject to the assumed validity of the effective expansion, $|v^2L|\lesssim1$: some apparent flat directions in the global fit only open up for coefficient values at which the EFT description has broken down.

\section{RGEs of the LEFT}\label{sec:RGELEFT}
Before we can map the LEFT to the SMEFT we have to consider the RGEs for the LEFT up to the EW scale to check whether heavy mixing with other operators occurs or the Wilson coefficients dramatically change -- we will see, however, that this is not the case. 
From \parencite[p.29]{jenkins2018lowAD} we find for the anomalous dimensions of the LNC operators $\mathcal{O}^{V,LL}_{\nu e}$ and $\mathcal{O}^{V,LR}_{\nu e}$:
\begin{align}
    \dot{L}^{V,LL}_{\substack{\nu e \\ prst}} &= \frac{4}{3} e^2 \mathsf{q}_e \delta_{st} \Bigg[ N_c \mathsf{q}_u \left( L^{V,LL}_{\substack{\nu u \\ prww}} + L^{V,LR}_{\substack{\nu u \\ prww}} \right) + N_c \mathsf{q}_d \left( L^{V,LL}_{\substack{\nu d \\ prww}} + L^{V,LR}_{\substack{\nu d \\ prww}} \right)\nonumber \\[5pt] 
    &\quad + \mathsf{q}_e \left( L^{V,LL}_{\substack{\nu e \\ prww}} + L^{V,LR}_{\substack{\nu e \\ prww}} \right) \Bigg]\nonumber \\[5pt] 
    &\quad + 96 e^2 \mathsf{q}_e^2 L_{\substack{\nu \gamma \\ wr}} L^*_{\substack{\nu \gamma \\ wp}} \delta_{st}
\end{align}
\begin{align} 
    \dot{L}^{V,LR}_{\substack{\nu e \\ prst}} =&\;\; \frac{4}{3} e^2 \mathsf{q}_e \delta_{st} \Bigg[ N_c \mathsf{q}_d \left( L^{V,LL}_{\substack{\nu d \\ prww}} + L^{V,LR}_{\substack{\nu d \\ prww}} \right) + N_c \mathsf{q}_u \left( L^{V,LL}_{\substack{\nu u \\ prww}} + L^{V,LR}_{\substack{\nu u \\ prww}} \right)\nonumber \\[5pt] 
    &\;\;+ \mathsf{q}_e \left( L^{V,LL}_{\substack{\nu e \\ prww}} + L^{V,LR}_{\substack{\nu e \\ prww}} \right) \Bigg]\nonumber \\[5pt] 
    &\;\;- 96 e^2 \mathsf{q}_e^2 L_{\substack{\nu \gamma \\ wr}} L^*_{\substack{\nu \gamma \\ wp}} \delta_{st}
\end{align}
Since $st=12\leftrightarrow 21$, the Kronecker delta $\delta_{st}$ in the closed fermions loop term and the dipole operator term make the anomalous dimensions for the muon decay LNC operators vanish:
\begin{align}
    \dot{L}^{V,LL}_{\nu e, ab12} = \dot{L}^{V,LL}_{\nu e, ab21}=0 \\[5pt]
    \dot{L}^{V,LR}_{\nu e, ab12} = \dot{L}^{V,LR}_{\nu e, ab21}=0
\end{align}
For the scalar and tensor LEFT operators, however, the anomalous dimensions do not evaluate to zero:
\begin{align}
    \dot{L}^{S,LL}_{\substack{\nu e \\ prst}} &= -6 e^2 \mathsf{q}_e^2 L^{S,LL}_{\substack{\nu e \\ prst}}\\[5pt]
    \dot{L}^{T,LL}_{\substack{\nu e \\ prst}} &= 2 e^2 \mathsf{q}_e^2 L^{T,LL}_{\substack{\nu e \\ prst}}\\[5pt]
    \dot{L}^{S,LR}_{\substack{\nu e \\ prst}} &= -6 e^2 \mathsf{q}_e^2 L^{S,LR}_{\substack{\nu e \\ prst}}
\end{align}
Solving for the three operators yields:
\begin{align}
    L^{S,LL}_{\nu e}(v) &= L^{S,LL}_{\nu e}(m_\mu)\left(\frac{v}{m_\mu}\right)^{-\frac{3\alpha}{2\pi}}\approx0.973\cdot L^{S,LL}_{\nu e}(m_\mu)\\
    L^{T,LL}_{\nu e}(v) &= L^{T,LL}_{\nu e}(m_\mu)\left(\frac{v}{m_\mu}\right)^{\frac{\alpha}{2\pi}}\approx 1.009\cdot L^{T,LL}_{\nu e}(m_\mu)\\
    L^{S,LR}_{\nu e}(v) &= L^{S,LR}_{\nu e}(m_\mu)\left(\frac{v}{m_\mu}\right)^{-\frac{3\alpha}{2\pi}}\approx 0.973\cdot L^{S,LR}_{\nu e}(m_\mu)
\end{align}
So we see that running of the couplings from the energy scale of muon decay ($\approx m_\mu$) to the EW scale ($\approx v$) only gives minor corrections -- but for completeness we are including them in our calculations. In the next chapter, we will see how the five LEFT operators under consideration match onto the SMEFT and how the corresponding SMEFT coefficients evolve under the RGEs.

\chapter{Standard Model Effective Field Theory (SMEFT)}\label{ch:smeft}
For the energy regime above the electroweak scale, the full Standard Model Effective Field Theory (SMEFT) is the appropriate effective description. Since we are exploring the imprints of heavy new physics on the phenomenology of muon decay we first have to identify the set of relevant SMEFT operators -- those operators which result in their low-energy limit in the LEFT operators identified in the previous chapter. For this purpose we are using the tree-level matching results of \cite{jenkins2018lowOM} and \parencite{liao2020extending}. However, while the matching is performed at the electroweak scale, the investigation of potential heavy new fields possibly shielded by these effective operators naturally takes place at the energy scale of new physics -- roughly at the order of a fex $\text{TeV}$, as we know from Tab.~\ref{tab:leftbounds}. So before we can \enquote{open up} the SMEFT operators we first have to \texttt{run} them up to that energy scale. To accomplish this task we are using the python packages \texttt{wilson} \cite{aebischer2018python} and the \texttt{D7RGESolver} \parencite{liao2025rge}. This will allow us to present the full set of SMEFT operators at the new physics scale which in their low-energy limit affect muon decay.\par
The content of this chapter is ordered in the following way: First, we briefly summarize the structure of SMEFT and its operators at dimensions five, six and seven. Subsequently, we cite the tree-level matching conditions for the relevant SMEFT operators at dimension six and seven to the LEFT operators which perturb muon decay. Finally, the identified SMEFT operators are run up to the energy scale of new physics where, in the next chapter, we can open them up to explore the various heavy new fields possibly veiled by them.
\section{The SMEFT Lagrangian}
The SMEFT extends the SM Lagrangian by embedding it as the leading term of an expansion in terms of powers of a heavy scale $\Lambda$: At each canonical mass dimension $D$ (in our convention the $\Lambda$s are absorbed by the Wilson coefficients, following \cite{jenkins2013rge1}) we add all local Lorentz invariant operators built out of the SM field content:
\begin{equation}
    \mathcal{L}=\mathcal{L}_\text{SM}+\sum_{D\ge5}\mathcal{L}^D
\end{equation}
Unlike the HEFT, the SMEFT assumes that electroweak symmetry is linearly realized and that the physical Higgs belongs to the $SU(2)_L$ doublet $H$. The higher dimensional interactions are therefore constructed from complete SM muliplets and respect the full SM gauge group $SU(3)_C\times SU(2)_L\times U(1)_Y$.\par 
At dimension five, there is only one operator (class) we can add -- the Weinberg neutrino mass operator \cite{weinberg1967model}, which combines two Higgs doublets with two lepton doublets:
\begin{equation}
    \mathcal{L}^5=C_5^{rs}\epsilon ^{ik}\epsilon^{kl}(l^T_{ir}Cl_{ks})H_jH_l+\text{h.c.}
\end{equation}
Here, $r,s$ are flavor indices, $i,j,k,l$ weak-eigenstate indices. For three generations, $C_5$ has 6 complex entries. After spontaneous symmetry breaking, this operator produces a Majorana mass term for the neutrino: $M_\nu=-v^2C_5$.\par 
For dimension six, Buchmüller and Wyler counted 80 operators in 1986 \cite{buchmuller1986effective}. However, many of these operators turn out to be actually equivalent and refer to the same operator structure because they can be transformed into each other. Fierz transformations, for example -- one of which we used in Eq.~\ref{eq:fierzingVLR2SRR} to \enquote{fierz} $L_{\nu e}^{V,LR}$ into $g^S_{RR}$ -- account for many subtle equivalences of scalar, vector and tensor-like structures. Moreover, the SM equations of motion and field redefinitions also allow for transforming many operator structure into each other. An updated classification was therefore published in 2010 which is now mostly referred to as the Warsaw basis \cite{grzadkowski2010dimension} and which counts 63 independent operators, 59 of which conserve baryon number. Of those 59 operators, 42 are self-hermitian and 17 not -- which is why Jenkins, Manohar \& Stoffer count $42+2(17)=76$ baryon-number conserving operators in \cite{jenkins2018lowOM} and $4+4=8$ baryon-number violating operators. At dimension seven, there are 18 independent operators \cite{liao2020extending} -- originally Lehman~\cite{lehman2014extending} presented 20 operators, but 2 pairs of $\psi^4D$ operators were shown to refer to the same structure. In Tab.~\ref{tab:SMEFTnumOP} we collect the number of operators at dimension five, six and seven, following the counting in \cite{jenkins2018lowOM} and \cite{liao2020extending}. We cite the full bases for dimensions six and seven in the appendix in Ch.~\ref{sec:base}.
\begin{table}[t]
    \centering
    \begin{tabular}{llll} \toprule
             D&$(\Delta B,\Delta L)$&$n_g=1$&$n_g=3$\\\midrule
             5&$(0,+2)$ + h.c.&$1+1$&6 + 6\\
             6&$(0,0)$&76&2499\\
             6&$(1,1)$ + h.c.&4 + 4&273 + 273\\
             7&$(0,+2)$ + h.c.&12 + 12&474 + 474\\
             7&$(+1,-1)$ + h.c.&6 + 6&297 + 297\\\bottomrule
    \end{tabular}
    \caption{Number of SMEFT operators at dimensions five, six and seven from \cite{jenkins2018lowOM} and \cite{liao2020extending} ($n_g=$ number of generations)}
    \label{tab:SMEFTnumOP}
\end{table}
\section{SMEFT in the broken phase}
In the spontaneously broken SMEFT effective Higgs operators alter the VEV, the couplings and the weak mixing angle -- for this discussion we are following the section of the same title in \cite{jenkins2018lowOM}. Essential for muon decay, the gauge currents themselves receive operator-dependent corrections. For example $\mathcal{W}^+_\mu$ couples to $(\bar\nu_{Lp}\gamma^\mu e_{Lr})$ via a modified coupling $[W_l]_{pr}$:
\begin{equation}
    \Delta\mathcal{L}_{Wl\nu}=-\frac{\bar g_2}{\sqrt{2}}\mathcal{W}_\mu^+\bar\nu_{Lp}\gamma^\mu[W_l]_{pr}e_{Lr}+\text{h.c.}
\end{equation}
The modified coupling $[W_l]_{pr}$ as well as $[Z_\psi]_{pr}$ which couples $\mathcal{Z}_\mu$ to $(\bar\psi_p\gamma^\mu\psi_r)$ whereby $\psi=\nu_L,e_L,e_R$ in the spontaneously broken SMEFT receive corrections from $\mathcal{O}^{(3)}_{Hl}$ as well as from the modified VEV:
\begin{align}
    [W_l]_{pr} &= \left[\delta_{pr} + v_T^2  C^{(3),pr}_{Hl} \right]\\[5pt]
    [Z_{\nu_L}]_{pr} &= \left[\delta_{pr}\left(\frac 12\right) - \frac12 v_T^2
    C^{(1),pr}_{Hl} + \frac12 v_T^2  C^{(3),pr}_{Hl} \right]\\[5pt]
    [Z_{e_L}]_{pr} &= \left[\delta_{pr}\left(-\frac 12+\bar s^2 \right) - \frac12 v_T^2  C^{(1),pr}_{Hl} - \frac12 v_T^2  C^{(3),pr}_{Hl} \right]\\[5pt]
    [Z_{e_R}]_{pr} &= \left[\delta_{pr}\bar s^2  - \frac12 v_T^2  C^{pr}_{He}\right]
\end{align}
The VEV and the gauge coupling change in the following way, which also affects the weak mixing angle:
\begin{align}
    v_T = \left( 1+ \frac{3 L_H v^2}{8 \lambda} \right) v\qquad \bar g_1= g_1 \left(1 + L_{HB} \, v_T^2 \right)\qquad \bar g_2= g_2 \left(1 + L_{HW} \, v_T^2 \right)\label{eq:vtgaugecorr}\\[5pt]
    \bar s =\sin \theta =  \frac{\bar g_1}{\sqrt{ \bar g_1^2 + \bar g_2^2}} \left[ 1 + \frac{\epsilon}{2}\frac{\bar g_2}{\bar g_1} \left( \frac{\bar g_2^2 - \bar g_1^2}{\bar g_1^2 + \bar g_2^2} \right)   \right]\hspace{2.2cm}
\end{align}
And the effective couplings $\bar e$ and $\bar g_Z$ as well:
\begin{align}
\bar e &= \bar g_2 \, \sin \bar \theta - \frac{1}{2} \, \cos \bar \theta  \, \bar g_2 \, v_T^2 \, C_{HWB}, \nonumber\\[5pt]
\bar g_Z &= \frac{\bar e}{\sin \bar\theta \cos \bar\theta}  \left[1 +  \frac{\bar g_1^2+\bar g_2^2}{2 \bar g_1\bar g_2} v_T^2C_{HWB}\right]
\end{align}
Equipped with these corrections we can now take a look on the matching conditions of the SMEFT onto the LEFT.
\section{Matching conditions}\label{sec:matchingLEFTSMEFT}
From the LEFT paper by Jenkins, Manohar and Stoffer \cite{jenkins2018lowOM} on operators and matching at dimension six we collect the following matching conditions for the LNC operators $\mathcal{O}^{V,LL}_{\nu e}$ and $\mathcal{O}^{V,LR}_{\nu e}$:
\begin{align}
    L^{V,LL}_{\nu e,prst}&= C_{ll}^{prst}+C_{ll}^{stpr}-\tfrac{\bar g_2^2}{2M_W^2}[W_l]_{pt}[W_l]^*_{rs}-\tfrac{\bar g_Z^2}{M_Z^2}[Z_{\nu_L}]_{pr}[Z_{e_L}]^*_{st}\label{eq:LVLLRGE}\\[5pt]
    L^{V,LR}_{\nu e,prst}&= C_{le}^{prst} -\tfrac{\bar g_Z^2}{M_Z^2}[Z_{\nu_L}]_{pr}[Z_{e_R}]^*_{st}\label{eq:LVLRRGE}
\end{align}
For the $prst=ab12$ flavor combination under consideration we are therefore obtaining, while neglecting all terms quadratic in couplings, the following matching conditions:
\begin{align}
    L^{V,LL}_{\nu e,ab12}=\, &C_{ll}^{ab12}+C_{ll}^{12ab}-\tfrac{\bar g_2^2}{2M_W^2}\left(\delta_{a2}\delta_{b1}+\delta_{a2}C^{(3),b1}_{Hl}+\delta_{b1}C^{(3),a2}_{Hl}\right)\label{eq:LVLLmatching}\\[5pt]
    &+\tfrac{\bar g_Z^2}{M_Z^2}\left(\frac{1}{4}\delta_{ab}v_T^2\left[C_{Hl}^{(1),12*}+C_{Hl}^{(3),12*}\right]\right)\nonumber\\[5pt]
    L^{V,LR}_{\nu e,ab12}=\, &C_{le}^{ab12}+\tfrac{\bar g_Z^2}{M_Z^2}\left(\frac{1}{4}\delta_{ab}v_T^2C_{He}^{12*}\right)
\end{align}
We have therefore 5 SMEFT Operators at dimension six which the LEFT matches onto (full basis in \ref{tab:dim6SMEFTbasis}):
\begin{align}
    \op_{ll}^{prst} & =(\bar{l}_p \gamma^\mu l_r)(\bar{l}_s \gamma_\mu l_t) \label{Oll}\\[5pt]
    \op_{le}^{prst}&=(\bar{l}_p \gamma^\mu l_r)(\bar{e}_s \gamma_\mu e_t)\\[5pt]
    \op_{Hl}^{(1),pr}&=(H^\dagger i\overleftrightarrow{D}_\mu H)(\bar{l}_p \gamma^\mu l_r)\\[5pt]
    \op_{Hl}^{(3),pr} &=(H^\dagger i\overleftrightarrow{D}^I_\mu H)(\bar{l}_p \tau^I \gamma^\mu l_r)\\[5pt]
    \op_{He}^{pr}&=(H^\dagger i\overleftrightarrow{D}_\mu H)(\bar{e}_p \gamma^\mu e_r)
\end{align}
We see from the last term of Eq. \ref{eq:LVLLmatching} that the SMEFT operator $Q_{Hl}^{(1)}$ only produces LEFT operators with two identical neutrino generations (due to the $\delta_{ab}$) -- this is because the operator stems from a lepton number changing neutral current mediated by the $Z$ boson which couples the muon directly to the electron and decays itself into two neutrinos of the same generation.\par  
For the scalar and tensor operators, since they violate lepton number, the matching at dimension six is zero but at dimension seven we obtain from \parencite{liao2020extending}:
\begin{align}
L^{S,LL}_{\nu e,prst} &= -\frac{\sqrt{2}v}{8} \left( 2C_{\bar{e}lllH}^{stpr} + C_{\bar{e}lllH}^{sptr} + p \leftrightarrow r \right) \label{eq:SLLtoD7}\\
L^{S,LR}_{\nu e,prst} &= -\frac{\sqrt{2}v}{2} \left( C_{leHD}^{pt}\delta^{rs} + C_{leHD}^{rt}\delta^{ps} \right) \label{eq:SLRtoD7} \\
L^{T,LL}_{\nu e, prst} &= +\frac{\sqrt{2}v}{32} \left( C_{\bar{e}lllH}^{sptr} - C_{\bar{e}lllH}^{srtp} \right) \label{eq:TLLtoD7}
\end{align}
Notice that in this case the $v$ is the bare SM VEV -- corrections in terms of $v_T$ would start at order $1/\Lambda^5$ and are therefore neglected in this tree-level result. As we see in Eq.~\ref{eq:SLLtoD7} to Eq.~\ref{eq:TLLtoD7} there are two relevant operators at dimension 7 (full basis in \ref{tab:dim7SMEFTbl02} and \ref{tab:dim7SMEFTbl1-1}):
\begin{align}
    \mathcal{O}_{leHD}^{pr} &= \epsilon_{ij}\epsilon_{mn} (\overline{l_p^{iC}} \gamma_\mu e_r) H^j H^m iD^\mu H^n \label{eq:leHD}\\
    \mathcal{O}_{\bar{e}lllH}^{prst} &= \epsilon_{ij}\epsilon_{mn} (\overline{e_p} l_r^i) (\overline{l_s^{jC}} l_t^m) H^n \label{eq:elllH}
\end{align}
The specific matching for the seven operators that we derived $95\%$ CL intervals for (Tab.~\ref{tab:leftbounds}) is:
\begin{align}
        L^{V,LL}_{\nu e,2112} &=-\frac{2}{v_T^2} + C_{ll}^{2112} + C_{ll}^{1221}-2C^{(3),22}_{Hl}-2C^{(3),11}_{Hl}\label{eq:specMatchVLL2112}\\
        L^{V,LR}_{\nu e,2112} &= C_{le}^{2112}\label{eq:specMatchVLR2112}\\
        L^{V,LL}_{\nu e,3312} &= C_{ll}^{3312}+C_{ll}^{1233}+\frac{1}{4}v_T^2\tfrac{\bar g_Z^2}{M_Z^2}\left(C_{Hl}^{(1),12*}+C_{Hl}^{(3),12*}\right)\\
        L^{S,LL}_{\nu e,1112} &= -\frac{\sqrt{2}v}{4}\left(2C_{\bar elllH}^{1211}+C_{\bar e lllH}^{1121}\right)\\
        L^{S,LR}_{\nu e,1112} &=-\sqrt{2}vC_{LeDH}^{12}\\
        L^{S,LR}_{\nu e,1212} &=\frac{\sqrt{2}v}{2}C_{LeDH}^{22}\\
        L^{T,LL}_{\nu e,1212} &=\frac{\sqrt{2}v}{32}\left(C_{\bar elllH}^{1122}-C_{\bar elllH}^{1221}\right)\label{eq:specMatchTLL1221}
\end{align}
Let us analyze the very first matching conditions in Eq. \ref{eq:specMatchVLL2112} for $ L^{V,LL}_{\nu e,2112}$: The first term $-\tfrac{2}{v_T^2}$ amounts to the pure SM $W$-boson exchange -- however, in the spontaneously broken SMEFT the Higgs potential is modified by the operator $\mathcal{O}_H = (H^\dagger H)^3$, see Eq. \ref{eq:vtgaugecorr}. Since we are working in the  $\{\hat{m}_W, \hat{m}_Z, \hat{\alpha}\}$ input scheme, $v_T$ is not determined by the Higgs potential but by the physical mass of the $W$ boson:
\begin{equation}
    m_W^2 = \frac{1}{4} \bar{g}_2^2 v_T^2
\end{equation}
We can not really distinguish whether our LEFT deviation is due to an SM gauge shift, a vertex correction or the four-fermion operator. We  neglect the corrections to $\bar g$ that stem from the SMEFT in the broken phase and refer to the technical report of the LHC EFT working group in~\cite{brivio2022truncation}: Generally products of two dimension-six coefficients, including electroweak input-parameter shifts multiplying an already dimension-six contribution, are consequently of $\mathcal O(\Lambda^{-4})$ and can be neglected. 
For operators that do not interfere with the Standard-Model amplitude, the squared EFT amplitude is retained as the leading non-vanishing contribution. Such squared terms are separately gauge invariant and basis independent, but constitute only a partial contribution at the
corresponding higher order in the EFT expansion. Moreover, SM gauge shifts and vertex corrections are heavily restricted by Electroweak Precision Observables (EWPO) and LEP data -- we must unfortunately refrain from a detailed discussion at this point and only note that exclusively constraining the four-fermion operator is legitimate.\par
Since furthermore $C_{ll}^{2112}=C_{ll}^{1221*}$ we can simplify the matching to:
\begin{equation}
    \text{Re}[L^{V,LL}_{\nu e,2112}]\approx 2\text{Re}[C_{ll}^{2112}]
\end{equation}
The matching for $L^{V,LR}_{\nu e,2112} = C_{le}^{2112}$ \ref{eq:specMatchVLR2112} is much cleaner -- and moreover we have constraints for both its real and imaginary part.\par 
The matching of \(L_{\nu e,1112}^{S,LL}\) requires some
care. Here, the single-coefficient assumption at the LEFT level means
that only \(L_{\nu e,1112}^{S,LL}\) is varied -- it does not imply that the matching selects a single independent SMEFT flavour coefficient. The first index of \(\mathcal O_{\bar elllH}^{prst}\) labels the
charged-lepton singlet, whereas its last three indices label the three lepton doublets and obey non-trivial permutation relations. We therefore rewrite the raw flavour components in the irreducible flavour basis used by \texttt{D7RGESolver} \parencite{liao2025rge}. For the flavour multiset \((r,s,t)=(1,1,2)\), the totally antisymmetric component vanishes, leaving one symmetric and one mixed-symmetry coefficient. We define:
\begin{equation}
     C_S\equiv C_{\bar elllH}^{(S),1112}
 \qquad
 C_M\equiv C_{\bar elllH}^{(M),1112}
\end{equation}

The projector relations then imply:
\begin{equation}
     C_{\bar elllH}^{1211}=C_S-C_M
 \qquad
 C_{\bar elllH}^{1121}=C_S
 \qquad
 C_{\bar elllH}^{1112}=C_S+C_M
\end{equation}

It follows that:
\begin{align}
 L_{\nu e,1112}^{S,LL}
 &=-\frac{\sqrt2v}{4}
   \left(2C_{\bar elllH}^{1211}
         +C_{\bar elllH}^{1121}\right)
 \nonumber\\
 &=-\frac{\sqrt2v}{4}\left(3C_S-2C_M\right).
 \label{eq:SLL-SM-decomposition}
\end{align}
The weak-isospin contractions by the Levi-Civita tensors should not be
confused with this flavour decomposition: The latter follows from
permutations of the three lepton-doublet flavour indices.
\section{RGEs for the SMEFT}
In order to consistently constraint UV fields we have to bridge the energy scale upwards to them. From the electroweak scale at about $246\text{ GeV}$ to a few TeV, as estimated in Tab.~\ref{tab:leftbounds}, it is only about one order of magnitude, so we don't expect large corrections -- but on a conceptual level the running and mixing of operators is a key element of exploring energy scales and translating between them: UV fields might produce effective operators that look quite different at a low-energy process like muon decay.\par 
For the dimension-six operators, we are using the \texttt{wilson}~2.5.2 package \parencite{aebischer2018python} and a continuously fixed Warsaw weak basis, for the dim-7 operators the \texttt{D7RGESolver} from \parencite{liao2025rge} in the down basis.
The full RGEs are given in section \ref{sec:apprge} in the appendix. In Tab.~\ref{tab:smeft-rge-primary} we give the complete complete one-coefficient translation of the muon decay LEFT bounds into SMEFT bounds and subsequently into UV bounds together with the total relative change of the components between $\mu_\text{EW}$ and $\Lambda_\text{est}$. In Tab.~\ref{tab:smeft-rge-mixing-d7} we collect the mixing of the of the dimension-seven coefficient into the strongest generated coefficient and the Weinberg operator in order to check whether those values are consistent with current neutrino mass bounds. In Tab.~\ref{tab:smeft-rge-mixing-d6} we collect the mixing of the dimension-six coefficients into several of the coefficients they produce.\par 
Equation~\ref{eq:SLL-SM-decomposition} shows that the LEFT measurement
constrains one linear combination of two independent SMEFT flavour
coefficients:
\begin{equation}
     v^2L_{\nu e,1112}^{S,LL}
    =-\frac{\sqrt2}{4}
   \left(3z_{S,1112}-2z_{M,1112}\right)
    \qquad z_{S,M}=v^3C_{S,M}
\end{equation}
If both coefficients are allowed simultaneously, only this combination
is bounded. The two rows in Tab.~\ref{tab:smeft-rge-primary} instead
represent the mutually exclusive one-coefficient hypotheses
\(z_M=0\) and \(z_S=0\) -- they must not be interpreted as a simultaneous
fit. \par 
The naive contributions in Tab.~\ref{tab:induced-neutrino-masses}
range from $\mathcal{O}(10^2\,\mathrm{eV})$ to
$\mathcal{O}(10^4\,\mathrm{eV})$ and are consequently many orders of
magnitude above the sub-eV neutrino-mass scale indicated by oscillation
data and constrained by direct and cosmological measurements
\parencite{particle2024review}. The strongest direct measurement comes from the KATRIN experiment from the kinematics of tritium $\beta$-decay. From the combined analysis of its first five measurement campaigns, KATRIN obtains \cite{katrin2025direct}:
\begin{equation}
    m_\beta^2
    \equiv \sum_{i=1}^{3}\lvert U_{ei}\rvert^2m_i^2
    =-0.14^{+0.13}_{-0.15}\,\mathrm{eV}^2
    \quad\implies\quad
    m_\beta<0.45\,\mathrm{eV}
    \quad (90\%~\mathrm{CL})
\end{equation}

If one imposed the genuinely
single-operator boundary condition $C_5(\Lambda)=0$ at the UV scale,
the downward RGE evolution would therefore generate neutrino masses
that are incompatible with observation when the dimension-seven
coefficients saturate their muon-decay bounds. In such a scenario,
neutrino-mass constraints would be considerably stronger than those
obtained from the Michel parameters, unless the mixing is suppressed by
a symmetry or cancelled by an independent contribution to the Weinberg
operator. The values in the table should nevertheless not be interpreted
as direct low-energy predictions of the present calculation: The RGE
evolution was performed upwards with $C_5(\mu_{\mathrm{EW}})=0$, so that
the generated $C_5(\Lambda)$ instead quantifies the UV coefficient
required to compensate the mixing and recover the chosen electroweak
boundary condition \parencite{liao2025rge}.
% Auto-generated transcription of smeft_rge_results.md.
% The tables use only packages already loaded in main.tex:
% booktabs, graphicx, amsmath, amssymb.
%
% Definitions:
%   x_LEFT = v^2 L,  z_d = v^(d-4) C_d,  v = mu_EW = 246.22 GeV.
% All intervals are single-coefficient 95% CL intervals.

\begin{table}[htbp]
    \vspace*{1cm}
    \centering
    \resizebox{\textwidth}{!}{%
    \renewcommand{\arraystretch}{1.25}
    \begin{tabular}{lllllllll}
        \toprule
        LEFT
        & SMEFT
        & Matching
        & \(x_{\mathrm{LEFT}}(m_\mu)\)
        & \(x_{\mathrm{LEFT}}(\mu_{\mathrm{EW}})\)
        & \(z_d(\mu_{\mathrm{EW}})\)
        & \(\Lambda_{\mathrm{est}}\,[\mathrm{TeV}]\)
        & \(z_d(\Lambda_{\mathrm{est}})\)
        & \(\Delta_{\mathrm{RGE}}\)
        \\
        \midrule
        \multicolumn{9}{c}{dimension-six SMEFT coefficients}\\
        \midrule
        \(\operatorname{Re}L_{\nu e,2112}^{V,LL}\)
        & \(\operatorname{Re}C_{ll}^{2112}\)
        & \(x_{\mathrm{LEFT}}=2z_6\)
        & \([-0.0044,+0.0024]\)
        & \([-0.0044,+0.0024]\)
        & \([-0.0022,+0.0012]\)
        & \(5.277\)
        & \([-0.0022,+0.0012]\)
        & \(-0.41\%\)
        \\
        \(\operatorname{Re}L_{\nu e,2112}^{V,LL}\)
        & \(\operatorname{Re}C_{Hl}^{(3),11}\)
        & \(x_{\mathrm{LEFT}}=-2z_6\)
        & \([-0.0044,+0.0024]\)
        & \([-0.0044,+0.0024]\)
        & \([-0.0012,+0.0022]\)
        & \(5.277\)
        & \([-0.0012,+0.0023]\)
        & \(+4.84\%\)
        \\
        \(\operatorname{Re}L_{\nu e,2112}^{V,LL}\)
        & \(\operatorname{Re}C_{Hl}^{(3),22}\)
        & \(x_{\mathrm{LEFT}}=-2z_6\)
        & \([-0.0044,+0.0024]\)
        & \([-0.0044,+0.0024]\)
        & \([-0.0012,+0.0022]\)
        & \(5.277\)
        & \([-0.0012,+0.0023]\)
        & \(+4.84\%\)
        \\
        \(\operatorname{Re}L_{\nu e,2112}^{V,LR}\)
        & \(\operatorname{Re}C_{le}^{2112}\)
        & \(x_{\mathrm{LEFT}}=z_6\)
        & \([-0.0213,+0.0177]\)
        & \([-0.0213,+0.0177]\)
        & \([-0.0213,+0.0177]\)
        & \(1.687\)
        & \([-0.0211,+0.0176]\)
        & \(-0.95\%\)
        \\
        \(\operatorname{Im}L_{\nu e,2112}^{V,LR}\)
        & \(\operatorname{Im}C_{le}^{2112}\)
        & \(x_{\mathrm{LEFT}}=z_6\)
        & \([-0.0165,+0.0226]\)
        & \([-0.0165,+0.0226]\)
        & \([-0.0165,+0.0226]\)
        & \(1.638\)
        & \([-0.0164,+0.0224]\)
        & \(-0.94\%\)
        \\
        \(\operatorname{Re}L_{\nu e,3312}^{V,LL}\)
        & \(\operatorname{Re}C_{ll}^{3312}\)
        & \(x_{\mathrm{LEFT}}=2z_6\)
        & \([-0.1320,+0.1320]\)
        & \([-0.1320,+0.1320]\)
        & \([-0.0660,+0.0660]\)
        & \(0.958\)
        & \([-0.0656,+0.0656]\)
        & \(-0.63\%\)
        \\
        \(\operatorname{Re}L_{\nu e,3312}^{V,LL}\)
        & \(\operatorname{Re}C_{Hl}^{(1),12}\)
        & \(x_{\mathrm{LEFT}}=\operatorname{Re}z_6^\ast\)
        & \([-0.1320,+0.1320]\)
        & \([-0.1320,+0.1320]\)
        & \([-0.1320,+0.1320]\)
        & \(0.678\)
        & \([-0.1363,+0.1363]\)
        & \(+3.21\%\)
        \\
        \(\operatorname{Re}L_{\nu e,3312}^{V,LL}\)
        & \(\operatorname{Re}C_{Hl}^{(3),12}\)
        & \(x_{\mathrm{LEFT}}=\operatorname{Re}z_6^\ast\)
        & \([-0.1320,+0.1320]\)
        & \([-0.1320,+0.1320]\)
        & \([-0.1320,+0.1320]\)
        & \(0.678\)
        & \([-0.1342,+0.1342]\)
        & \(+1.63\%\)
        \\
        \midrule
        \multicolumn{9}{c}{dimension-seven SMEFT coefficients}\\
        \midrule
        \(\operatorname{Re}L_{\nu e,1112}^{S,LL}\)
        & \(\operatorname{Re}C_{\bar e LLLH}^{(S),1112}\)
        & \(x_{\mathrm{LEFT}}=-\frac{3\sqrt2}{8}z_7\)
        & \([+0.0364,+0.0689]\)
        & \([+0.0355,+0.0671]\)
        & \([-0.1265,-0.0669]\)
        & \(0.490\)
        & \([-0.1282,-0.0678]\)
        & \(+1.36\%\)
        \\
        \(\operatorname{Re}L_{\nu e,1112}^{S,LL}\)
        & \(\operatorname{Re}C_{\bar e LLLH}^{(M),1112}\)
        & \(x_{\mathrm{LEFT}}=\frac{\sqrt2}{4}z_7\)
        & \([+0.0364,+0.0689]\)
        & \([+0.0355,+0.0671]\)
        & \([+0.1003,+0.1898]\)
        & \(0.428\)
        & \([+0.1005,+0.1902]\)
        & \(+0.22\%\)
        \\
        \(\operatorname{Re}L_{\nu e,1112}^{S,LR}\)
        & \(\operatorname{Re}C_{LeDH}^{12}\)
        & \(x_{\mathrm{LEFT}}=-\sqrt2\,z_7\)
        & \([-0.0251,+0.0251]\)
        & \([-0.0244,+0.0244]\)
        & \([-0.0173,+0.0173]\)
        & \(0.953\)
        & \([-0.0184,+0.0184]\)
        & \(+6.40\%\)
        \\
        \(\operatorname{Re}L_{\nu e,1212}^{S,LR}\)
        & \(\operatorname{Re}C_{LeDH}^{22}\)
        & \(x_{\mathrm{LEFT}}=\frac{\sqrt2}{2}z_7\)
        & \([-0.0355,+0.0355]\)
        & \([-0.0345,+0.0345]\)
        & \([-0.0488,+0.0488]\)
        & \(0.674\)
        & \([-0.0511,+0.0511]\)
        & \(+4.81\%\)
        \\
        \(\operatorname{Re}L_{\nu e,1212}^{T,LL}\)
        & \(\operatorname{Re}C_{\bar e LLLH}^{(M),1122}\)
        & \(x_{\mathrm{LEFT}}=\frac{\sqrt2}{32}z_7\)
        & \([-0.0076,+0.0076]\)
        & \([-0.0077,+0.0077]\)
        & \([-0.1743,+0.1743]\)
        & \(0.441\)
        & \([-0.1747,+0.1747]\)
        & \(+0.23\%\)
        \\
        \bottomrule
    \end{tabular}%
    }
    \caption{%
        Complete one-coefficient translation of the muon decay LEFT bounds
        into SMEFT constraints and their subsequent evolution to the estimated
        UV scale. We define \(x_{\mathrm{LEFT}}=v^2L\) and
        \(z_d=v^{d-4}C_d\), with
        \(v=\mu_{\mathrm{EW}}=246.22\,\mathrm{GeV}\). The scale estimate assumes
        \(\lvert C_d\rvert=\Lambda^{4-d}\), so that
        \(\Lambda_{\mathrm{est}}
        =v/\max\lvert z_d\rvert^{1/(d-4)}\).
        The column \(\Delta_{\mathrm{RGE}}\) gives the relative change of the
        indicated real or imaginary component between
        \(\mu_{\mathrm{EW}}\) and \(\Lambda_{\mathrm{est}}\).
    }
    \label{tab:smeft-rge-primary}
\end{table}



\begin{table}[htbp]
    \centering
    \renewcommand{\arraystretch}{1.25}
    \begin{tabular}{lllll}
        \toprule
        EW-scale seed
        & \(\Lambda_{\mathrm{est}}\)
        & Change
        & D7 Cf.
        & Induced Weinberg Cf.
        \\
        \midrule
        \(\operatorname{Re}C_{\bar eLLLH}^{(S),1112}\)
        & \(0.490\,\mathrm{TeV}\)
        & \(+1.36\%\)
        & \(C_{\bar Q uLLH}^{3321}\;(3.55\times10^{-6}\%)\)
        & \(C_{LH5}^{12}\;(1.172\times10^{-8})\)
        \\
        \(\operatorname{Re}C_{\bar eLLLH}^{(M),1112}\)
        & \(0.428\,\mathrm{TeV}\)
        & \(+0.22\%\)
        & \(C_{\bar Q uLLH}^{3312}\;(3.76\times10^{-6}\%)\)
        & \(C_{LH5}^{12}\;(3.101\times10^{-9})\)
        \\
        \(\operatorname{Re}C_{LeDH}^{12}\)
        & \(0.953\,\mathrm{TeV}\)
        & \(+6.40\%\)
        & \(C_{\bar d LueH}^{3132}\;(0.0703\%)\)
        & \(C_{LH5}^{12}\;(1.610\times10^{-6})\)
        \\
        \(\operatorname{Re}C_{LeDH}^{22}\)
        & \(0.674\,\mathrm{TeV}\)
        & \(+4.81\%\)
        & \(C_{\bar d LueH}^{3232}\;(0.0527\%)\)
        & \(C_{LH5}^{22}\;(2.363\times10^{-6})\)
        \\
        \(\operatorname{Re}C_{\bar eLLLH}^{(M),1122}\)
        & \(0.441\,\mathrm{TeV}\)
        & \(+0.23\%\)
        & \(C_{\bar Q uLLH}^{3322}\;(1.99\times10^{-6}\%)\)
        & \(C_{LH5}^{22}\;(6.527\times10^{-9})\)
        \\
        \bottomrule
    \end{tabular}%
    \caption{%
        Complete dimension-seven mixing generated by
        \texttt{D7RGESolver} using integrated one-loop running in the down
        basis. The forth column displays the strongest induced dimension-seven operator and the percentages refers to the EWSB-normalized
        induced-to-input ratios. The final column separately displays the
        largest generated Weinberg coefficient as
        \(v\lvert C_5\rvert/(v^3\lvert C_{7,\mathrm{input}}\rvert)\);
        it is a dimensionless ratio rather than a percentage.
    }
    \label{tab:smeft-rge-mixing-d7}
\end{table}

\begin{table}[htbp]
    \centering
    \small
    \renewcommand{\arraystretch}{1.25}
    \begin{tabular}{lllllll}
        \toprule
        EW seed
        & \(\Lambda_{\mathrm{est}}\)
        & Notable?
        & Induced &$\tfrac{C_i(\Lambda_\text{est})}{C_\text{input}(\mu_\text{EW})}$ &Induced&$\tfrac{C_i(\Lambda_\text{est})}{C_\text{input}(\mu_\text{EW})}$
        \\
        \midrule
        \(\operatorname{Re}C_{ll}^{2112}\)
        & \(5.277\,\mathrm{TeV}\)
        & yes
        & \(C_{ll}^{1122}\)&\(4.33\%\)&
          \(C_{ll}^{2332}\)&\(0.536\%\)\\&&&
          \(C_{ll}^{1331}\)&\(0.536\%\)&
          \(C_{ll}^{2222}\)&\(0.368\%\)\\&&&
          \(C_{ll}^{1111}\)&\(0.368\%\)&
          \(C_{Hl}^{(3),22}\)&\(0.266\%\)\\&&&
          \(C_{Hl}^{(3),11}\)&\(0.266\%\)&
          \(C_{lq}^{(3),2211}\)&\(0.266\%\)
        \\\midrule
        \(\operatorname{Re}C_{Hl}^{(3),11}\)
        & \(5.277\,\mathrm{TeV}\)
        & yes
        & \(C_{H\Box}\)&\(1.75\%\)&
          \(C_{lq}^{(3),1133}\)&\(1.18\%\)\\&&&
          \(C_{uH}^{33}\)&\(0.953\%\)&
          \(C_{Hl}^{(3),22}\)&\(0.569\%\)\\&&&
          \(C_{Hl}^{(3),33}\)&\(0.569\%\)&
          \(C_{Hq}^{(3),11}\)&\(0.569\%\)\\&&&
          \(C_{Hq}^{(3),22}\)&\(0.569\%\)&
          \(C_{Hq}^{(3),33}\)&\(0.567\%\)
        \\\midrule
        \(\operatorname{Re}C_{Hl}^{(3),22}\)
        & \(5.277\,\mathrm{TeV}\)
        & yes
        & \(C_{H\Box}\)&\(1.75\%\)&
          \(C_{lq}^{(3),2233}\)&\(1.18\%\)\\&&&
          \(C_{uH}^{33}\)&\(0.953\%\)&
          \(C_{Hl}^{(3),11}\)&\(0.569\%\)\\&&&
          \(C_{Hl}^{(3),33}\)&\(0.569\%\)&
          \(C_{Hq}^{(3),11}\)&\(0.569\%\)\\&&&
          \(C_{Hq}^{(3),22}\)&\(0.569\%\)&
          \(C_{Hq}^{(3),33}\)&\(0.567\%\)
        \\\midrule
        \(\operatorname{Re}C_{le}^{2112}\)
        & \(1.687\,\mathrm{TeV}\)
        & no
        & \(C_{lequ}^{(1),1133}\)&\(0.00257\%\)&&
        \\\midrule
        \(\operatorname{Im}C_{le}^{2112}\)
        & \(1.638\,\mathrm{TeV}\)
        & no
        & \(C_{lequ}^{(1),1133}\)&\(0.00253\%\)&&
        \\\midrule
        \(\operatorname{Re}C_{ll}^{3312}\)
        & \(0.958\,\mathrm{TeV}\)
        & yes
        & \(C_{ll}^{1332}\)&\(2.12\%\)&
          \(C_{le}^{1222}\)&\(0.150\%\)\\&&&
          \(C_{le}^{1211}\)&\(0.150\%\)&
          \(C_{le}^{1233}\)&\(0.150\%\)\\&&&
          \(C_{lu}^{1211}\)&\(0.101\%\)&
          \(C_{lu}^{1222}\)&\(0.101\%\)\\&&&
          \(C_{lu}^{1233}\)&\(0.101\%\)&
          \(C_{ll}^{1222}\)&\(0.0774\%\)
        \\\midrule
        \(\operatorname{Re}C_{Hl}^{(1),12}\)
        & \(0.678\,\mathrm{TeV}\)
        & yes
        & \(C_{lu}^{1233}\)&\(1.03\%\)&
          \(C_{lq}^{(1),1233}\)&\(0.539\%\)\\&&&
          \(C_{le}^{1211}\)&\(0.0558\%\)&
          \(C_{le}^{1222}\)&\(0.0558\%\)\\&&&
          \(C_{le}^{1233}\)&\(0.0556\%\)&
          \(C_{lu}^{1211}\)&\(0.0373\%\)\\&&&
          \(C_{lu}^{1222}\)&\(0.0373\%\)&
          \(C_{ll}^{1222}\)&\(0.0281\%\)
        \\\midrule
        \(\operatorname{Re}C_{Hl}^{(3),12}\)
        & \(0.678\,\mathrm{TeV}\)
        & no
        & \(C_{lq}^{(3),1233}\)&\(0.431\%\)&
          \(C_{ll}^{1332}\)&\(0.176\%\)\\&&&
          \(C_{ll}^{1112}\)&\(0.0892\%\)&
          \(C_{ll}^{1222}\)&\(0.0892\%\)\\&&&
          \(C_{lq}^{(3),1211}\)&\(0.0878\%\)&
          \(C_{lq}^{(3),1222}\)&\(0.0870\%\)\\&&&
          \(C_{ll}^{1233}\)&\(0.0868\%\)&
          \(C_{lq}^{(3),1232}\)&\(0.0217\%\)
        \\
        \bottomrule
    \end{tabular}
        \caption{%
        Complete dimension-six operator mixing generated between the EW scale 
        \(\mu_{\mathrm{EW}}=246.22\,\mathrm{GeV}\) and the
        one-coefficient scale estimate \(\Lambda_{\mathrm{est}}\).
        Every percentage is the magnitude of the induced coefficient divided
        by the input coefficient. The numerical evolution uses
        \texttt{wilson}~2.5.2 in a continuously fixed Warsaw weak basis.
        ``Notable'' denotes mixing of at least \(1\%\); coefficients down to
        \(0.001\%\) are listed.
    }
    \label{tab:smeft-rge-mixing-d6}
\end{table}

\begin{table}[htbp]
    \centering
    \renewcommand{\arraystretch}{1.25}
    \begin{tabular}{llll}
        \toprule
        D7 Cf.
        & Weinberg
        & Neutrino mass-matrix entry
        & Naïve induced mass scale
        \\
        \midrule
        $C_{\bar e LLLH}^{(S),1112}$
        & $C_5^{12}$
        & $\lvert(M_\nu)_{e\mu}\rvert$
        & $193\text{--}365\,\mathrm{eV}$
        \\
        $C_{\bar e LLLH}^{(M),1112}$
        & $C_5^{12}$
        & $\lvert(M_\nu)_{e\mu}\rvert$
        & $76.6\text{--}145\,\mathrm{eV}$
        \\
        $C_{LeHD}^{12}$
        & $C_5^{12}$
        & $\lvert(M_\nu)_{e\mu}\rvert$
        & $\lesssim 6.86\,\mathrm{keV}$
        \\
        $C_{LeHD}^{22}$
        & $C_5^{22}$
        & $\lvert(M_\nu)_{\mu\mu}\rvert$
        & $\lesssim 28.4\,\mathrm{keV}$
        \\
        $C_{\bar e LLLH}^{(M),1122}$
        & $C_5^{22}$
        & $\lvert(M_\nu)_{\mu\mu}\rvert$
        & $\lesssim 280\,\mathrm{eV}$
        \\
        \bottomrule
    \end{tabular}
    \caption{%
        Naïve Majorana-neutrino mass scales associated with the
        Weinberg coefficients induced by the dimension-seven SMEFT
        operators. The values are obtained from
        $\lvert(M_\nu)_{pr}\rvert=v^2\lvert C_5^{pr}\rvert$
        using the endpoints of the corresponding $95\%$ confidence
        intervals from muon decay. The first two rows are ranges because the corresponding
        Michel-parameter intervals do not contain zero. For the remaining
        rows, whose intervals do contain zero, only the maximal mass scale
        at the endpoint of the interval is quoted. The coefficients
        $C_5^{12}$ and $C_5^{22}$ describe entries of the Majorana mass
        matrix in the charged-lepton flavor basis and should not be
        identified with individual mass eigenvalues $m_1$, $m_2$, or $m_3$.
    }
    \label{tab:induced-neutrino-masses}
\end{table}



\chapter{Beyond Standard Model (BSM) physics}
We now want to investigate the UV fields that could possibly leave their imprint on muon decay. The direction of our inquiry, as we are coming from the low-energy regime, makes the necessary line of reasoning quite intricate. Inferring the effective low-energy limit of a given UV extension of the SM is relatively straight-forward: We construct an amplitude with a heavy degree of freedom and integrate it out, as we did in Sec.~\ref{sec:intoutHDOF}. As we match the resulting expression to the corresponding effective contact interaction we find its Wilson coefficient. And as we carry out this calculation for all constructable amplitudes we derive the complete set of Wilson coefficients and thus the full low-energy EFT. However, finding all UV origins for a given effective operator by such a top-down approach requires extensive research -- which comes down to constructing all low-energy EFTs for all possible UV models (of a certain type). So for the directive of this work a more effective bottom-up strategy would clearly be more adequate.\par 
Luckily the bottom-up approach has been made effective by the work on local on-shell amplitudes of recent years. Specifically Li, Ni, Xiao \& Yu showed in \cite{li2022bottom} the correspondence of on-shell contact amplitudes and effective operators and De Angelis provided in \cite{de2022amplitude} an algorithm for finding all kinematically independent massive amplitudes from UV states. This allows Li, Ren \& Yu in \cite{li2024complete} to present a complete UV dictionary for SMEFT operators up to dimension seven. Precisely, their dictionary allows us to look up all BSM particles of spin $0,\tfrac{1}{2},1$ that lead to renormalizable UV interactions which would then match at tree-level to a SMEFT operator of our interest -- which for our endevaour means a SMEFT operator that could disturb muon decay.\par 
We order the steps of the thought process in the following way: First we look up all UV fields that produce the SMEFT operators we derived in Sec.~\ref{sec:matchingLEFTSMEFT}. In the following section we analyze the Lagrangian for a specific UV field -- $S_4$ -- so we are well prepared to derive our results in the next chapter.

\section{The UV dictionary}
\begin{table}[tbp]
    \centering
    %\renewcommand{\arraystretch}{1.5}
    \small
    \begin{tabular}{lccccccccc}
        \toprule
        Notation \cite{li2024complete}& $S_2$ & $S_4$ & $S_6$ & $F_1$ & $F_2$ & $F_3$&$F_4$\\
        Notation \cite{de2018effective} & $\mathcal{S}_1$ &$\varphi$ & $\Xi_1$ & $N$ & $E^c$ & $\Delta^c_1$& $\Delta^c_3$\\
        Irrep & $({\bf 1,1})_1$ & $({\bf 1,2})_{\frac{1}{2}}$ & $({\bf 1,3})_1$ & $({\bf 1,1})_0$ & $({\bf 1,1})_1$ & $({\bf 1,2})_{\frac{3}{2}}$&$({\bf 1,2})_{\frac{1}{2}}$\\
        \midrule
        Notation \cite{li2024complete}& $F_5$ & $F_6$ & $V_1$ & $V_2$ & $V_3$ & $V_4$&  \\
        Notation \cite{de2018effective} & $\Sigma$ & $\Sigma_1^c$ & $\mathcal{B}$ & $\mathcal{B}_1$ & $\mathcal{L}_3^\dagger$ & $\mathcal{W}$& \\
        Irrep & $({\bf 1,3})_0$& $({\bf 1,3})_1$& $({\bf 1,1})_0$ & $({\bf 1,1})_1$ & $({\bf 1,2})_{\frac{3}{2}}$ & $({\bf 1,3})_0$&\\
        \bottomrule
    \end{tabular}
    \caption{BSM scalars, fermions and vectors that produce the operators $\mathcal{O}_{ll}, \mathcal{O}_{le}, \mathcal{O}^{(1)}_{Hl}, \mathcal{O}^{(3)}_{Hl}, \mathcal{O}_{He},\mathcal{O}_{leHD}, \mathcal{O}_{\bar elllH}$ at tree-level. Fermions have a superscript ${}^c$ if they are conjugated with each other. Majorana fermions $F_1$ and $F_5$ have parity left, i.e. $F = F_L = P_L F$. Vectors have a superscript ${}^\dagger$ if they are conjugated with each other.}
    \label{tab:rBSMsfv}
\end{table}
We use the dictionary from \cite{li2024complete} to determine the possible UV origins of our SMEFT operators -- the complete list of new UV states that at tree-level produce SMEFT operators of dimension five, six and seven is cited in Tab.~\ref{tab:BSMsfv} in the appendix. We see that the dimension six operators are produced by adding a single new UV field to the Lagrangian. The four-lepton operators $\op_{ll}$ and $\op_{le}$ naturally need a bosonic mediator -- a vertex of three fermions would not produce a Lorentz scalar and have half-integer canonical mass dimension after all. The operators $\op^{(1,3)}_{Hl}$ on the other hand allows for both a vector and a fermion channel -- either connecting the currents $(H^\dagger DH)$ and $(ll)$ with each other or $H^\dagger l$ with $Hl$. Our two dimension seven operators $\mathcal{O}_{leHD}$ and $\mathcal{O}_{\bar elllH}$ are produced by either adding $F_1$ or $F_5$ or combinations of fields. The fermions $F_1\sim({\bf 1,1})_0$ and $F_5\sim({\bf 1,3})_0$ both have zero hypercharge and transform as real representations of $SU(2)_L$ -- so they are self-conjugate and their Majorana mass exactly provides the $\Delta L=2$ for producing a dimension seven SMEFT operator that violates lepton number by two. For the remaining completions, no single non-Majorana field suffices -- two heavy propagators are needed in order to attach all the external fields and violate lepton number. 
\begin{align}
    \mathcal{O}_{ll}&\to S_2,S_6, V_1, V_4\label{eq:bsmOll}\\
    \mathcal{O}_{le}&\to S_4, V_1, V_3\\
    \mathcal{O}^{(1)}_{Hl}&\to F_1,F_2,F_5,F_6,V_1\\
    \mathcal{O}^{(3)}_{Hl}&\to F_1,F_2,F_5,F_6,V_4\\
    \mathcal{O}_{He}&\to F_3,F_4,V_1\\
    \mathcal{O}_{leHD}&\to F_1,F_5,(F_3,S_6),(F_3,V_2),(V_2,V_3)\label{eq:bsmOleHD}\\
    \mathcal{O}_{\bar elllH}&\to F_1,F_5,(S_2,S_4),(S_2,F_4),(S_4,S_6),(S_6,F_4)\label{eq:bsmOelllH}
\end{align}
The representations of all the UV fields above are given in Tab.~\ref{tab:rBSMsfv}. The notation in the first line of the table is that of Li, Ren \& Yu in \cite{li2024complete}, the notation in the second line the more common names as cited in \cite{de2018effective}. We would now like to analyze the (possible) UV origins of the dimension six coefficients and obtain mass limits. As we see in Eq.~\ref{eq:bsmOll} the coefficient $C_{ll}$ could be produced by scalar mediator $S_2$ or $S_6$ or a vector mediator $V_1$ or~$V_4$.\par 
Let us take a look at the precise matching conditions from the dictionary:
\begin{align}
\begin{autobreak}
 {C_{ll}} =
-\frac{\mathcal{D}_{{llS_{2}}}^{{f_1f_2p_1}} \mathcal{D}_{{llS_{2}}}^{{f_4f_3p_1*}}}{2 {M_{S_{2}}^2}}
+\frac{\mathcal{D}_{{llS_{6}}}^{{f_1f_2p_1}} \mathcal{D}_{{llS_{6}}}^{{f_4f_3p_1*}}}{4 {M_{S_{6}}^2}}
-\frac{\mathcal{D}_{{ll^\dagger V_{1}}}^{{f_1f_3p_1}} \mathcal{D}_{{ll^\dagger V_{1}}}^{{f_2f_4p_1}}}{2 {M_{V_{1}}^2}}
+\frac{\mathcal{D}_{{ll^\dagger V_{4}}}^{{f_1f_3p_1}} \mathcal{D}_{{ll^\dagger V_{4}}}^{{f_2f_4p_1}}}{4 {M_{V_{4}}^2}}
-\frac{\mathcal{D}_{{ll^\dagger V_{4}}}^{{f_1f_4p_1}} \mathcal{D}_{{ll^\dagger V_{4}}}^{{f_2f_3p_1}}}{2 {M_{V_{4}}^2}}
+\frac{\mathcal{D}_{{llS_{2}}}^{{f_2f_1p_1}} \mathcal{D}_{{llS_{2}}}^{{f_4f_3p_1*}}}{2 {M_{S_{2}}^2}}
+\frac{\mathcal{D}_{{llS_{6}}}^{{f_2f_1p_1}} \mathcal{D}_{{llS_{6}}}^{{f_4f_3p_1*}}}{4 {M_{S_{6}}^2}} 
\end{autobreak}
\end{align}
The indices $(f_1,f_2,f_3,f_4)$ are flavor indices which in Eq.~\ref{Oll} are labeled $(s,r,t,p)$. We see that the matching conditions have the following general structure:
\begin{equation}
    C_i^{(6)}(M_X)=\kappa\frac{\lambda_a\lambda_b^*}{M_X^2}
\end{equation}
Here we use $i=ll,le,Hl(1),Hl(3),He$ (the full matching conditions are cited in Eq.~\ref{eq:ll-UV} to Eq.~\ref{eq:elllH-UV} in the appendix) and $a,b$ are flavor indices and $\kappa$ is a numerical or group theoretic factor from the matching procedure. If we denote by $R$ the correction from the (LEFT or SMEFT) RGE and remember that our constrain $\epsilon_\text{allowed}$ relates to the LEFT coefficient by a factor of VEV squared $\epsilon_\text{allowed}\equiv v^2L_\text{allowed}(\mu_\mu)$ we can derive the general formula for the mass bounds on our heavy mediators:
\begin{equation}
    M_X>v\cdot\sqrt{\frac{|R\kappa\lambda_a\lambda_b^*|}{\epsilon_\text{allowed}}}
\end{equation}
For the dimension seven operators we take a look at the matching condition for $C_{leHD}$ -- the operator $\op_{\bar e lllH}$ has quite an opulent matching and we only cite it in the appendix in Eq.~\ref{eq:elllH-UV}. 
\begin{align}
\begin{autobreak}
 {C_{leHD}} =
-\frac{{y}_{{e}}^{{p_1f_5}} \mathcal{D}_{{F_{1}Hl}}^{{p_2f_1}} \mathcal{D}_{{F_{1}Hl}}^{{p_2p_1}}}{2 {M_{F_{1}}^3}}
-\frac{{y}_{{e}}^{{p_1f_5}} \mathcal{D}_{{llS_{6}}}^{{f_1p_1p_2}} \mathcal{C}_{{HHS_{6}^\dagger }}^{{p_2}}}{{M_{S_{6}}^4}}
-\frac{{y}_{{e}}^{{p_1f_5}} \mathcal{D}_{{F_{5}Hl}}^{{p_2f_1}} \mathcal{D}_{{F_{5}Hl}}^{{p_2p_1}}}{4 {M_{F_{5}}^3}}
+\frac{{y}_{{e}}^{{p_1f_5}} \mathcal{C}_{{HHS_{6}^\dagger }}^{{p_2}} \mathcal{D}_{{llS_{6}}}^{{p_1f_1p_2}}}{{M_{S_{6}}^4}}
+\frac{2 i \mathcal{D}_{{eF_{1}V_{2}}}^{{f_5p_1p_2}} \mathcal{D}_{{F_{1}Hl}}^{{p_1f_1}} \mathcal{D}_{{HHV_{2}^\dagger D}}^{{p_2}}}{{M_{F_{1}}} {M_{V_{2}}^2}}
-\frac{\mathcal{D}_{{e^\dagger F_{3R}^\dagger H^\dagger }}^{{f_5p_1*}} \mathcal{D}_{{F_{1}F_{3R}^\dagger H}}^{{p_2p_1}} \mathcal{D}_{{F_{1}Hl}}^{{p_2f_1}}}{{M_{F_{1}}} {M_{F_{3}}^2}}
-\frac{\mathcal{D}_{{e^\dagger F_{3R}^\dagger H^\dagger }}^{{f_5p_1*}} \mathcal{D}_{{F_{3L}F_{5}H^\dagger }}^{{p_1p_2*}} \mathcal{D}_{{F_{5}Hl}}^{{p_2f_1}}}{{M_{F_{3}}} {M_{F_{5}}^2}}
+\frac{2 i \mathcal{D}_{{e^\dagger F_{3R}^\dagger H^\dagger }}^{{f_5p_1*}} \mathcal{D}_{{F_{3L}l^\dagger V_{2}^\dagger }}^{{p_1f_1p_2*}} \mathcal{D}_{{HHV_{2}^\dagger D}}^{{p_2}}}{{M_{F_{3}}} {M_{V_{2}}^2}}
-\frac{\mathcal{D}_{{e^\dagger F_{3R}^\dagger H^\dagger }}^{{f_5p_1*}} \mathcal{D}_{{F_{3R}^\dagger F_{5}H}}^{{p_1p_2}} \mathcal{D}_{{F_{5}Hl}}^{{p_2f_1}}}{2 {M_{F_{3}}^2} {M_{F_{5}}}}
-\frac{\mathcal{D}_{{e^\dagger F_{3R}^\dagger H^\dagger }}^{{f_5p_1*}} \mathcal{D}_{{F_{3R}^\dagger lS_{6}}}^{{p_1f_1p_2}} \mathcal{C}_{{HHS_{6}^\dagger }}^{{p_2}}}{{M_{F_{3}}^2} {M_{S_{6}}^2}}
+\frac{\mathcal{D}_{{e^\dagger F_{5}S_{6}^\dagger }}^{{f_5p_1p_2*}} \mathcal{D}_{{F_{5}Hl}}^{{p_1f_1}} \mathcal{C}_{{HHS_{6}^\dagger }}^{{p_2}}}{{M_{F_{5}}^2} {M_{S_{6}}^2}}
-\frac{2 i \mathcal{D}_{{e^\dagger l^\dagger V_{3}^\dagger }}^{{f_5f_1p_1*}} \mathcal{C}_{{HHS_{6}^\dagger }}^{{p_2}} \mathcal{D}_{{HS_{6}V_{3}^\dagger D}}^{{p_2p_1}}}{{M_{S_{6}}^2} {M_{V_{3}}^2}}
-\frac{4 i \mathcal{D}_{{e^\dagger l^\dagger V_{3}^\dagger }}^{{f_5f_1p_1*}} \mathcal{D}_{{HHV_{2}^\dagger D}}^{{p_2}} \mathcal{C}_{{HV_{2}V_{3}^\dagger }}^{{p_2p_1}}}{{M_{V_{2}}^2} {M_{V_{3}}^2}} 
\end{autobreak}\label{eq:leHD-UV}
\end{align}
We see that for a SM extension of a single fermion $F_1$ or $F_5$ the mass is bound by a term of the following form:
\begin{equation}
    M_X=v\cdot\left(\frac{|R\kappa\lambda_\text{eff}|}{\epsilon_\text{allowed}}\right)^\frac{1}{3}
\end{equation}
For a two-field completion the experiment only constraints:
\begin{equation}
    M_X^2M_Y>\frac{|R\kappa\lambda_1\lambda_2\lambda_3|v^3}{\epsilon_\text{allowed}}
\end{equation}
Only if we assume $M_X\approx M_Y$ the cube root formula applies.
\section{Correlated mediators}
Especially interesting are mediators that produce multiple SMEFT operators that disturb muon decay -- as for example the BSM vector field $V_1$, which matches onto $\op_{ll},\op_{le}$ and $\op_{Hl}^{(1)}$ all-together. Schematically $V_1$ couples to three currents:
\begin{equation}
    \mathcal{L}_\text{BSM}\supset V_{1\mu}(\lambda_HJ^\mu_H+\lambda_{l,pr}J_l^{\mu,pr}+\lambda_{e,pr}J_e^{\mu,pr})
\end{equation}
So integrating out $V_1$ generates simultaneously:
\begin{equation}
    C_{ll}\sim\frac{\lambda_l\lambda_l}{M_{V_1}^2}\qquad C_{le}\sim\frac{\lambda_l\lambda_e}{M_{V_1}^2}\qquad 
    C_{Hl}^{(1)}\sim\frac{\lambda_H\lambda_l}{M_{V_1}^2}\label{eq:V}
\end{equation}
These should be constrained together by the LEFT observables -- which we will do in the following chapter. 


\chapter{Results}
We begin by discussing the constraints for UV fields that would produce the operator $\mathcal{O}_{le}$. From the LEFT bounds in Tab.~\ref{tab:leftbounds} and from the matching condition in Eq.~\ref{eq:specMatchVLR2112} we know that the real and imaginary part of its coefficient lie in this $95\%$-confidence region:
\begin{equation}
    \text{Re}[C_{le}^{2112}]\text{ in }[-0.021,+0.018]\text{ and Im}[C_{le}^{2112}]\text{ in }[-0.017,+0.023]
\end{equation}
The UV matching of this operator was cited in Eq.~\ref{eq:UVMatchCle} and for $C_{le}^{2112}$ amounts to:
\begin{equation}
     C_{le}^{2112} =
     \frac{\mathcal{D}_{{e^\dagger eV_{1}}}^{{12p_1}} \mathcal{D}_{{ll^\dagger V_{1}}}^{{12p_1}}}{{M_{V_{1}}^2}}
    -\frac{\mathcal{D}_{{e^\dagger lS_{4}^\dagger }}^{{11p_1}} \mathcal{D}_{{e^\dagger lS_{4}^\dagger }}^{{22p_1*}}}{2 {M_{S_{4}}^2}}
    +\frac{\mathcal{D}_{{e^\dagger l^\dagger V_{3}^\dagger }}^{{12p_1}} \mathcal{D}_{{e^\dagger l^\dagger V_{3}^\dagger }}^{{21p_1*}}}{{M_{V_{3}}^2}}
\end{equation}
If we assume single-mediator dominance, a single heavy copy accordingly, and a unit real coupling product, the appropriate limits for the masses are:
\begin{align}
    M_{V_1},M_{V_3}&>\frac{v}{\sqrt{0.018}}\simeq1.84\text{ TeV}\\
    M_{S_4}&>\frac{v}{\sqrt{2(0.021)}}\simeq1.2\text{ TeV}
\end{align}
For $S_4$ we used the lower end of the $95\%$ CL region and for $V_1,V_3$ the upper end due to the sign in their matching condition. The SMEFT RGEs only change the Wilson coefficient by about $1\%$ as we see from Tab.~\ref{tab:smeft-rge-primary}.
\section{Heavy Scalar \texorpdfstring{$\mathbf{S_4}$}{S4}}
It's time to \textit{change gears} and reverse the direction: We take a look at which SMEFT operators $S_4$ produces and if they would show up in muon decay. From the dictionary we can look up the top-down translation of $S_4$:
\begin{equation}
    S_4\to \op_H,\op_{eH},\op_{uH},\op_{dH},\op_{le},\op^{(1)}_{qu},\op^{(8)}_{qu},\op^{(1)}_{qd},\op^{(8)}_{qd},
    \op_{leqd},\op_{quqd}^{(1)},\op^{(1)}_{lequ}
\end{equation}
But we do can not really make any educated assumptions on the hypothetical coupling of $S_4$ to the quark sector or the Higgs. Let us assume those other couplings are zero -- can we infer something nevertheless? \par
The Lagrangian term of $S_4$ relevant for its matching conditions is this:
\begin{align}
\begin{autobreak}
{\Delta\mathcal{L}_{S_4}} =
- 2 \mathcal{D}_{{d^\dagger QS_{4}^\dagger }}^{{rsp}} [(\bar{d}_{{r}})^{{a}}({q}_{{s}})_{{ai}}]({S}_{{4p}}^{\dagger })^{{i}} 
- \mathcal{D}_{{e^\dagger lS_{4}^\dagger }}^{{rsp}} [\bar{e}_{{r}}({l}_{{s}})_{{i}}]({S}_{{4p}}^{\dagger })^{{i}} 
- 2 \mathcal{D}_{{QS_{4}u^\dagger }}^{{rps}}\epsilon ^{{ji}} [(\bar{u}_{{s}})^{{a}}({q}_{{r}})_{{aj}}]({S}_{{4p}})_{{i}}  
- \mathcal{D}_{{HH^\dagger H^\dagger S_{4}(1)}}^{{p}}\delta _{{k}}^{{i}} \delta _{{l}}^{{j}}{H}_{{j}}{H}^{{\dagger k}}{H}^{{\dagger l}}({S}_{{4p}})_{{i}} 
- \mathcal{D}_{{HH^\dagger H^\dagger S_{4}(2)}}^{{p}}\delta _{{l}}^{{i}} \delta _{{k}}^{{j}}{H}_{{j}}{H}^{{\dagger k}}{H}^{{\dagger l}}({S}_{{4p}})_{{i}}
\end{autobreak}\label{eq:UVLagS4}
\end{align}
So the matching condition of $C_{le}$ to $S_4$ in $prst$ convention is the following:
\begin{equation}
 {C_{le,prst}} = -\frac{\mathcal{D}_{{e^\dagger lS_{4}^\dagger }}^{{sr}} \mathcal{D}_{{e^\dagger lS_{4}^\dagger }}^{{tp*}}}{2 {M_{S_{4}}^2}} \label{eq:UVMatchCle}
\end{equation}
We get for $C_{le,2112}$ the following matching:
\begin{equation}
{C_{le,2112}} = -\frac{\mathcal{D}_{{e^\dagger lS_{4}^\dagger }}^{{11a}} \mathcal{D}_{{e^\dagger lS_{4}^\dagger }}^{{22a*}}}{2 {M_{S_{4}}^2}}=-\frac{\lambda_e\lambda_\mu^*}{2M_{S_4}^2}
\end{equation}
We thus can infer:
\begin{equation}
{|C_{le,1111}}| = -\frac{\lambda_e^2}{2M_{S_4}^2}\qquad{|C_{le,2222}|} = -\frac{\lambda_\mu^2}{2M_{S_4}^2}\qquad{|C_{le,1221}|} = -\frac{\lambda_e^*\lambda_\mu}{2M_{S_4}^2}
\end{equation}
Together with our assumptions $\lambda_e\simeq1$ and $\lambda_\mu\simeq1$ this translates to:
\begin{equation}
    |C_{le}^{1111}| > 1.2\text{ TeV}\qquad{|C_{le}^{2222}|} > 1.2\text{ TeV} \qquad{|C_{le}^{2112}|} > 1.2\text{ TeV}
\end{equation}
For $\op_{le,1111}=(\bar l_e\gamma_\mu l_e)(\bar e_R\gamma^\mu e_R)$ probably Bhabha scattering $e^+e^-\to e^+e^-$ provides the most direct probe. Here, $C_{le,1111}$ interferes with the SM photon and Z-exchange amplitudes. A recent scalar-doublet analysis using the LEP differential distributions quotes \cite{erdelyi2025probing}:
\begin{equation}
    \frac{M_{S_4}}{|\lambda_e|}>2.7\text{ TeV}\qquad (95\%\text{ CL})
\end{equation}
At $M_{S_4}=2.7\text{ TeV}$, the unit-coupling prediction is already:
\begin{equation}
    |v^2C_{le}^{2112}|=\frac{v^2}{2M_{S_4}^2}\simeq 0.00416
\end{equation}
This is well inside the Michel parameter interval. However, muon decay still gives complimentary sensitivity to the complex phase. For the coefficient $C_{le,2222}$ a future high-energy muon collider could probe the muonic analogue of Bhabha scattering, $\mu^+\mu\to\mu^+\mu^-$ and significantly extend current sensitivity to muon contact interactions \cite{glioti2025flavor}.\par 
As $S_4$ is an electroweak scalar doublet, it contains charged and neutral components. The gauge interactions are fixed by its representation, independently of $\lambda_e$ or $\lambda_\mu$, and the components could thus be produced through electroweak processes involving $\gamma,Z,W^\pm$. The expected decay would be then:
\begin{align}
    S^0_4&\to e^+e^-\qquad S^0_4\to \mu^+\mu^-\\
    S^\pm_4&\to e^\pm\nu_e\qquad S^\pm_4\to\mu^\pm\nu_\mu
\end{align}
It has been shown that a charged scalar decaying into electrons or muons would provide an especially clean collider signature \cite{alcaide2020lhc}.\par
Falkowski \& Mimouni discuss in \cite{falkowski2016model} the constraints on SMEFT couplings from Michel parameter data as well and remark, like us, that at linear order the best bound comes from $\eta$ and $\beta'/A$ on the real and imaginary part of $C_{le}^{2112}$ ($[c_{le}]_{1221}$ in their notation) respectively. They refrain from giving an explicit bound on $S_4$ (\enquote{inert Higgs}) exactly because LEP data provides the stronger bound. 

\section{Correlated mediator \texorpdfstring{$\mathbf{V_1}$}{V1}}
The heavy mediator $V_1$ produces multiple LEFT coefficients at the scale of muon decay and we use it to obtain a mass bound from a correlated fit. Let us first take a look at its Lagrangian:
\begin{align}
\begin{autobreak}
    {\Delta\mathcal{L}_{V_4}=}

- 2 \mathcal{D}_{{d^\dagger dV_{1}}}^{{rsp}} [(\bar{d}_{{r}})^{{a}}\gamma _{\mu }({d}_{{s}})_{{a}}]{V}_{{1p}}^{\mu } 
 - \mathcal{D}_{{e^\dagger eV_{1}}}^{{rsp}}1[\bar{e}_{{r}}\gamma _{\mu }{e}_{{s}}]{V}_{{1p}}^{\mu } 
 - 2 \mathcal{D}_{{HH^\dagger V_{1}D}}^{{p}} [D_{\mu }{H}_{{i}}]{H}^{{\dagger i}}{V}_{{1p}}^{\mu } 
 + \mathcal{D}_{{ll^\dagger V_{1}}}^{{rsp}} [(\bar{l}_{{s}})^{{i}}\gamma _{\mu }({l}_{{r}})_{{i}}]{V}_{{1p}}^{\mu } 
 + 2 \mathcal{D}_{{QQ^\dagger V_{1}}}^{{rsp}} [(\bar{q}_{{s}})^{{ai}}\gamma _{\mu }({q}_{{r}})_{{ai}}]{V}_{{1p}}^{\mu } 
 - 2 \mathcal{D}_{{u^\dagger uV_{1}}}^{{rsp}} [(\bar{u}_{{r}})^{{a}}\gamma _{\mu }({u}_{{s}})_{{a}}]{V}_{{1p}}^{\mu } 
\end{autobreak}\label{eq:V1Lag}
\end{align}
We choose real unit couplings $\lambda_l=\lambda_e=h=1$ and define:
\begin{align}
    s\equiv \frac{v}{M_V}\qquad v=0.246\text{ TeV}\\
    x_{ab}\equiv v^2 L_{\nu e,ab12}^{V,LL}\qquad y_{ab}\equiv v^2 L_{\nu e,ab12}^{V,LR}
\end{align}
The nonzero coefficients are then:
\begin{equation}
    \begin{aligned}
         (a,b) &\quad x_{ab} & y_{ab}  \\[2mm]
         (2,1) &\quad -s^2   & s^2     \\
         (1,2) &\quad -s^2   & s^2     \\
         (1,1),(2,2),(3,3) &\quad s^2   & -s^2 
    \end{aligned}
\end{equation}
If we substitute this into our full-fermion formulas from chapter 3 we receive:
\begin{align}
    R_G(s)\equiv\frac{G_F^2}{G_{F,\rm SM}}&=1+s^2+\frac{5}{2}s^4\\
    \rho(s)=\delta(s)&=\frac{3}{4}\\
    \xi(s)=\xi'(s)&=1-\frac{5}{2}s^4\\
    \xi''(s)&=1\\
    \eta(s)=-\eta''(s)&=\frac{1}{2}s^2+\frac{3}{4}^4\\
    \frac{\alpha'}{A}=\frac{\beta'}{A}&=0
\end{align}
Define a total chi-squared over our five experimental blocks (TWIST, Danneberg, $\xi$, $\xi'$ and the electroweak input scheme for $G_{F,\rm SM}$):
\begin{equation}
    \chi^2(s)=\chi^2_T(s)+\chi^2_D(s)+\left(\frac{\xi'(s)-1.00}{0.04}\right)^2+\left(\frac{\xi''(s)-0.98}{0.04}\right)^2+\chi_{G_F}^2(s)
\end{equation}
And we obtain:
\begin{equation}
    M_{V_1}>\frac{v}{s_{95}}=\frac{0.246\text{ TeV}}{0.0595}=4.13\text{ TeV}
\end{equation}
This is consistent with our first naïve result in Eq.~\ref{eq:naiveresult} and moreover with the findings of Falkowski \& Mimouni in Sec.~5.2 of \cite{falkowski2016model} where they report $M_V/\lambda_{l,e}>3.3\text{ to }10\text{ TeV}$ depending on the ratio of left- and right-handed couplings.

\subsubsection{Input-parameter effects of $\mathbf{C_{HD}}$}
\label{sec:v1-chd-input}

The neutral vector singlet $V_1\sim(\mathbf 1,\mathbf 1)_0$ may couple to
the Hermitian Higgs current according to:
\begin{equation}
    \Delta\mathcal L_{V_1}
    \supset
    -\lambda_h V_{1,\mu}J_H^\mu
    \qquad
    J_H^\mu=H^\dagger i\overleftrightarrow D^\mu H
\end{equation}
Integrating out $V_1$ generates, among other operators:
\begin{equation}
    C_{HD}=-\frac{2\lambda_h^2}{M_{V_1}^2}
    \label{eq:v1-chd-matching-short}
\end{equation}
This coefficient does not induce a four-lepton contact interaction or
modify the charged-current $Wl\nu$ vertex directly.  Its effect on
muon decay instead arises through the relations between the Lagrangian
parameters and the measured electroweak input observables.  Writing the
Higgs doublet after electroweak symmetry breaking as
\begin{equation}
    H=\frac{1}{\sqrt2}
      \begin{pmatrix}0\\v+h\end{pmatrix},
\end{equation}
one finds, to linear order in $C_{HD}$:
\begin{equation}
    m_W^2=\frac{g_2^2v^2}{4}
    \qquad
    m_Z^2=\frac{(g_1^2+g_2^2)v^2}{4}
    \left(1+\frac12v^2C_{HD}\right)
    \label{eq:chd-input-masses-short}
\end{equation}
Thus $C_{HD}$ changes the values of $g_1$, $g_2$ and $v$ inferred from a
chosen set of electroweak inputs.  In the
$\{\widehat m_W,\widehat m_Z,\widehat\alpha\}$ scheme used here, keeping
the measured inputs fixed gives:
\begin{equation}
    \frac{\delta v}{v}
    =-\frac{c_W^2}{4s_W^2}v^2C_{HD}
    \qquad
    \frac{\delta G_F^{\mathrm{pred}}}{G_F^{\mathrm{pred}}}
    =\frac{c_W^2}{2s_W^2}v^2C_{HD}
    =-\frac{c_W^2}{s_W^2}
      \frac{\lambda_h^2v^2}{M_{V_1}^2}
    \label{eq:chd-input-gf-short}
\end{equation}
Since the decay rate scales as $\Gamma_\mu\propto G_F^2$, this produces
an overall normalization shift:
\begin{equation}
    \left.\frac{\delta\Gamma_\mu}{\Gamma_\mu}\right|_{C_{HD}}
    =-2\frac{c_W^2}{s_W^2}
      \frac{\lambda_h^2v^2}{M_{V_1}^2}.
\end{equation}
It does not, at this order, change the normalized Michel-spectrum shape
parameters.  If $G_F$ is itself chosen as an electroweak input, the same
effect is absorbed into the inferred value of $v$ and instead appears in
the predictions for other electroweak observables.

A consistent treatment would therefore require the simultaneous
redefinition of the electroweak input parameters and the inclusion of
electroweak precision observables together with the remaining operators
generated by $V_1$.  Such a global analysis lies beyond the scope of this
work and is not pursued further.  

\section{Dimension seven results}\label{sec:results-d7}
The constraints from the dimension-seven operators are considerably lower than the ones from the dimension-six SMEFT -- this is because their leading contribution to the Michel observables scale as.
\begin{equation}
\Delta\mathcal{O}\sim |z_7|^2
\sim\left(\frac{v}{\Lambda}\right)^6
\end{equation}
In \cite{fridell2024probing} the authors constrain the operator $\op_{\bar elllH}$ via non-standard muon decay: $L^{S,LL}_{\nu e,2112}$ leads to $\Lambda_\text{NP}\approx 0.4\text{ TeV}$
and $L^{S,LL}_{\nu e,2112}$ to $\Lambda_\text{NP}\approx 0.2\text{ TeV}$ under the assumption of single LEFT operator dominance to $\mu^+\to e^+ +\bar\nu_e+\bar\nu_\mu$. The results agree at the expected order of magnitude -- our somewhat higher coefficient-scale sensitivity follows from a tighter Michel-parameter constraint and from using projected symmetric or mixed flavor directions rather than the single raw coefficient employed in the earlier estimate.

\chapter{Outlook}
We have perturbed muon decay with LEFT operators and constrained two of their possible UV origins. At the many intermediate steps in this investigation we made assumptions and exclusions that could be questioned. We will briefly reflect on some of these which might point toward future analyses:

The heavy mediator $V_1$ also produces bosonic operators which could propagate all the way down to muon decay -- this should be included

Questions:
\begin{itemize}
    \item 
    \item what if new physics does not follow spontaneous symmetry breaking -- how would we perform these calculations in the HEFT -- is there a matching onto the LEFT?
\end{itemize}
\appendix
\chapter{Kinematics}
\section{Michel basis functions}

\section{Kinematics and Phase Space integrals}\label{sec:kinPS}

In this section we closely follow the textbook by Scheck \cite{scheck2012electroweak} and define the momenta as follows:
\begin{equation}
\mu_q\to \overset{\scriptscriptstyle(-)}{\nu}_{k_1}+\overset{\scriptscriptstyle(-)}{\nu}_{k_2}+e_p\, ,
\end{equation}
%
that is the momentum assigned to the muon is $q$, that for the electron is $p$, etc. The labels $k_1$ and $k_2$ are integrated over in deriving the Michel parameters, so the assignment to (massless) neutrinos is arbitrary. When treating tau decays, the tau momentum is also taken to be $q$.

We take $s_0=(0,P_\mu \hat n_0)$ to be the muon spin four-vector, with $P_\mu$ the muon polarization and $\vec n_0$ a unit vector in the direction of the muon spin expectation value. A similar spin vector can be defined for the electron,
%
\begin{equation}
s_1=\left(\frac{1}{m_e}\vec p\cdot \hat n_1, \hat n_1+\frac{\vec p\cdot \hat n_1}{m_e(E+m_e)}\vec p\right)\, ,
\end{equation}
%
which projects onto the electron spin state with spin along $\vec n_1$ \cite{Scheck:1996ur}.

We work in the muon rest frame, $q^\mu=(m_\mu,\vec 0)$, and express the electron four-momentum us $ p^\mu=(E_e,\vec p)$. Working out the relevant products of four-vectors we obtain:
%
\begin{align}
q^2&=m_\mu^2\, ,\\
p^2&=m_e^2\, ,\\
q\cdot p &= m_\mu E_e\, ,\\
q\cdot s_0 &= 0\, ,\\
q\cdot s_1 &= \frac{m_\mu}{m_e}\vec p\cdot \hat n_1\, ,\\
p\cdot s_0 &= -P_\mu\,\vec p\cdot \hat n_0\, ,\\
p\cdot s_1 &= 0\, ,\\
s_0\cdot s_1 &=-P_\mu\left(\hat n_0\cdot\hat n_1+\frac{\vec p\cdot \hat n_1\, \vec p\cdot \hat n_0}{m_e(E_e+m_e)}\right)\, .
\end{align}


The PDG defines muon decay kinematics dependent on the electron spin in terms of the product of $\hat \zeta\cdot \vec P_e$. From a theory perspective we are able to choose to align $\zeta$ with the $\hat n_1$ direction.

The product $p\cdot s_0$ is independent of the electron spin orientation:
%
\begin{equation}
p\cdot s_0=-P_\mu|\vec p|\cos\theta\, .
\end{equation}
%
Taking $\hat e_3$ to be along $\vec p$ we find $q\cdot s_1$ lies strictly along $\hat e_3=\hat p$:
%
\begin{eqnarray}
q\cdot s_1=\frac{m_\mu}{m_e}|\vec p|\zeta_3\, .
\end{eqnarray}
%
The vectors $\hat p$ and $\hat n_0$ define the decay plane. The $\hat e_1$ direction is taken to be the component of the muon spin direction orthogonal to the electron momentum,
%
\begin{equation}
\hat e_1=\frac{\hat n_0-(\hat n_0\cdot \hat p)\hat p}{\sqrt{1-(\hat n_0\cdot \hat p)^2}}=\frac{\hat n_0-\cos\theta\hat p}{\sin\theta}\, ,
\end{equation}
%
and therefore lies in the decay plane orthogonal to $\hat p$.
The product of the electron and muon spins then decomposes over the $\hat e_1$ and $\hat e_3$ directions:
%
\begin{equation}
s_0\cdot s_1=-P_\mu\left(\sin\theta \zeta_1+\frac{E_e}{m_e}\cos\theta\zeta_3\right)\, .
\end{equation}
%
Finally $\hat e_2$ is taken to be normal to the decay plane. In this case only the structure, $\hat n_0\cdot (\hat n_1\times \hat p)$ contributes, which in four vectors corresponds to:
%
\begin{equation}
\epsilon_{\mu\nu\rho\sigma}s_0^\mu s_1^\nu q^\rho p^\sigma=m_\mu P_\mu |\vec p|\sin\theta\zeta_2
\end{equation}

\section{Phase space integrations}
In the LEFT, we need to perform phase space integrations of squared amplitudes dependent on higher orders in momentum than treated in Scheck or for example \cite{Marquez:2022bpg}. Taking $Q\equiv k_1+k_2=q-p$, and defining
%
\begin{equation}
\int d\Phi_2=\int \frac{d^3k_1}{2k_1^0}\frac{d^3k_2}{2k_2^0}\delta^{(4)}(Q-k_1-k_2)\, ,
\end{equation}
%
we obtain:
%
\begin{align}
\int d\Phi_2\, 1&=\frac{\pi}{2}\, ,\\
\int d\Phi_2\, k_1^\mu&=\frac{\pi}{4}Q^\mu\, ,\\
\int d\Phi_2\, k_1^\mu k_1^\nu&=\frac{\pi}{24}\left(4Q^\mu Q^\nu-Q^2\eta^{\mu\nu}\right)\, ,\\
\int d\Phi_2\, k_1^\mu k_1^\nu k_1^\rho&=\frac{\pi}{48}\left[6Q^\mu Q^\nu Q^\rho-Q^2\left(Q^\mu \eta^{\nu\rho}+Q^\nu \eta^{\mu\rho}+Q^\rho\eta^{\mu\nu}\right)\right]\, .
\end{align}
%
It is important to note that all momenta in these basis integrands are the same, i.e. they are all $k_1$. These results also implicitly assume massless neutrinos. When applying these identities to our results we simply substitute $k_2=Q-k_1$ so that our integrals depend only on $k_1$. 

The decay width is for the muon is then given as:
%
\begin{equation}
\frac{d^2\Gamma}{dE_e d\cos\theta}=\frac{1}{4m_\mu}\frac{\sqrt{E_e^2-m_e^2}}{(2\pi)^5}\int d\phi\, d\Phi_2|\, \mathcal M|^2\, ,
\end{equation}
%
with $\mathcal M$ the amplitude for the process. As our definition of the two body phase space excludes factors or $2\pi$ some additional factors appear explicitly in our expression for the differential decay width. For identical neutrinos in the final state an additional factor of one half is necessary.

Defining the maximum electron energy to be,
%
\begin{equation}
W = \frac{1}{2m_\mu}(m_\mu^2+m_e^2)\, ,
\end{equation}
%
we make a change of variables to $x\equiv E_e/W$ giving:
%
\begin{equation}
\frac{d^2\Gamma}{dx d\cos\theta}=\frac{W}{4m_\mu}\frac{\sqrt{W^2x^2-m_e^2}}{(2\pi)^5}\int d\phi\, d\Phi_2|\, \mathcal M|^2\, .
\end{equation}
%
\section{Michel ``basis functions''}\label{sec:michelbasis}

We define the following set of basis functions to aid in extracting the Michel parameters. Starting with electron polarization independent Michel parameters we define:
%
\begin{eqnarray}
\phi_1(x,x_0,\cos\theta)&=&\sqrt{x^2-x_0^2}x(1-x)\, ,\\
\phi_2(x,x_0,\cos\theta)&=&\sqrt{x^2-x_0^2}\frac{2}{9}(4x^2-3x-x_0^2)\, ,\\
\phi_3(x,x_0,\cos\theta)&=&\sqrt{x^2-x_0^2}x_0(1-x)\, ,\\
\phi_4(x,x_0,\cos\theta)&=&\cos\theta(x^2-x_0^2)\frac{1}{3}(1-x)\, ,\\
\phi_5(x,x_0,\cos\theta)&=&\cos\theta(x^2-x_0^2)\frac{2}{9}\left[(4x-3)+\sqrt{1-x_0^2}-1\right]\, .
\end{eqnarray}

These correspond to the kinematic structures multiplying the Michel parameters in the differential decay rate. For example, considering only $\rho$ dependence from the decay rate we have:
%
\begin{equation}
\frac{d^2\Gamma}{dxd\cos\theta}\sim\rho \frac{m_\mu}{4\pi^3}W^4G_F^2\sqrt{x^2-x_0^2}\frac{2}{9}(4x^2-3x-x_0^2)\, .
\end{equation}
%
Defining the Gram matrix,
%
\begin{equation}
G_{ij}=\int_{x_0}^1dx\, \int_{-1}^1d\cos\theta\phi_i\phi_j\, ,
\end{equation}
%
we can solve for the Michel parameters by taking:
%
\begin{equation}
\begin{pmatrix}
A\equiv \frac{G_F^2 m_\mu W^4}{4\pi^3}\\
A\rho\\
A\eta\\
\pm P_\mu A \xi\\
\pm P_\mu A \xi\delta
\end{pmatrix}=G_{ij}^{-1}b_j\, ,\label{eq:solveforMichel}
\end{equation}
%
where we have defined:
%
\begin{equation}
b_j=\left(\int_{x_0}^1dx\, \int_{-1}^{1}d\cos\theta \frac{d^2\Gamma}{dxd\cos\theta}\phi_j\right)\, .
\end{equation}
%
For muon decay $\pm1$ should be understood to be $+1$ for positive muon decay, and $-1$ for negative muon decay.

The same approach is taken for the other Michel parameters. We split these cases by their projection onto our choice of $\zeta_i$ as mentioned in Sec.~\ref{sec:kinPS}.



%
\begin{equation}
\begin{aligned}
G^{-1}b_{\rm unpol}
&=
\begin{pmatrix}
A\\[2pt]
A\rho\\[2pt]
A\eta\\[2pt]
- P_\mu A\xi\\[2pt]
- P_\mu A\xi\delta
\end{pmatrix},
\\[8pt]
G^{-1}b_{\zeta_1}
&=
\begin{pmatrix}
-\dfrac{P_\mu A}{6}
\left[
\xi''+12\left(\rho-\dfrac{3}{4}\right)
\right]
\\[8pt]
-\dfrac{P_\mu A}{4}\eta
\\[8pt]
\dfrac{P_\mu A}{12}\eta''
\end{pmatrix},
\\[8pt]
G^{-1}b_{\zeta_2}
&=
\begin{pmatrix}
P_\mu A\,\dfrac{\alpha'}{A}
\\[8pt]
\dfrac{2}{3}P_\mu A\,\sqrt{1-x_0^2}\,\dfrac{\beta'}{A}
\end{pmatrix},
\\[8pt]
G^{-1}b_{\zeta_3}
&=
\begin{pmatrix}
-\dfrac{A}{6}\xi'
\\[8pt]
-\dfrac{2A}{27}\xi\left(\delta-\dfrac{3}{4}\right)
\\[8pt]
\dfrac{P_\mu A}{6}\xi''
\\[8pt]
\dfrac{2P_\mu A}{3}\left(\rho-\dfrac{3}{4}\right)
\\[8pt]
\dfrac{P_\mu A}{3}\eta''
\end{pmatrix}.
\end{aligned}
\label{eq:solveforMichelAll}
\end{equation}
%



\paragraph*{$\mathbf{\zeta_1}$}
The $\zeta_1$ projection corresponds to $F_{T_1}$ in the PDG notation \cite{particle2024review}. The basis functions we choose are:
%
\begin{eqnarray}
\phi_{\zeta_1,1}(x,x_0,\cos\theta)&=&\sqrt{x^2-x_0^2}\sqrt{1-\cos^2\theta}(1-x)x_0\, ,\\
\phi_{\zeta_1,2}(x,x_0,\cos\theta)&=&\sqrt{x^2-x_0^2}\sqrt{1-\cos^2\theta}(x^2-x_0^2)\, ,\\
\phi_{\zeta_1,3}(x,x_0,\cos\theta)&=&\sqrt{x^2-x_0^2}\sqrt{1-\cos^2\theta}(4x-3x^2-x_0^2)\, .\\
\end{eqnarray}
%
These basis functions allow for the extraction of $\xi''$, $\eta$, and $\eta''$ respectively. The simultaneous extraction of $\eta$ from $\phi_{\zeta_1,2}$ and $\phi_3$ provides a useful consistency check of our results.

\paragraph*{$\mathbf{\zeta_2}$}
The $\zeta_2$ projection corresponds to $F_{T_2}$ in the PDG notation \cite{particle2024review}. The basis functions we choose are:
%
\begin{eqnarray}
\phi_{\zeta_2,1}(x,x_0,\cos\theta)&=&(x^2-x_0^2)\sqrt{1-\cos^2\theta}(1-x)\, ,\\
\phi_{\zeta_2,2}(x,x_0,\cos\theta)&=&(x^2-x_0^2)\sqrt{1-\cos^2\theta}\sqrt{1-x_0^2}\, .
\end{eqnarray}
%
These basis functions allow for the extraction of $\alpha'/A$ and $\beta'/A$ respectively.



\paragraph*{$\mathbf{\zeta_3}$}
The $\zeta_3$ projection corresponds to $F_{IP}$ and $F_{AP}$ in the PDG notation \cite{particle2024review}. The basis functions we choose are:
%
\begin{eqnarray}
\phi_{\zeta_3,1}(x,x_0,\cos\theta)&=&(x^2-x_0^2)(2-2x+\sqrt{1-x_0^2})\, ,\\
\phi_{\zeta_3,2}(x,x_0,\cos\theta)&=&(x^2-x_0^2)(4x-4+\sqrt{1-x_0^2})\,,\\
\phi_{\zeta_3,3}(x,x_0,\cos\theta)&=&\cos\theta\sqrt{x^2-x_0^2}(2x^2-x-x_0^2)\, ,\\
\phi_{\zeta_3,4}(x,x_0,\cos\theta)&=&\cos\theta\sqrt{x^2-x_0^2}(4x^2-3x-x_0^2)\, ,\\
\phi_{\zeta_3,5}(x,x_0,\cos\theta)&=&\cos\theta\sqrt{x^2-x_0^2}(1-x)x_0\, .
\end{eqnarray}
%
These basis functions allow for the extraction of $\xi'$, $\delta$, $\xi''$, $\rho$, and $\eta''$. Again, redundancy with previously extracted Michel parameters provides a useful cross check of our results.







\chapter{Operator bases}\label{sec:base}
\section{LEFT basis}
Here, we give the operator bases of the LEFT at dimension six and the SMEFT at dimensions six and seven. We highlighted the operators that perturb muon decay with a grey background color and the matching operators of the SMEFT as well. 
\begin{table}[htbp]
    \centering
    \renewcommand{\arraystretch}{1.5}
    
    % ================= TOP ROW =================
    \scalebox{0.67}{%
    \begin{tabular}[t]{c|l}
        \multicolumn{2}{c}{\textbf{$\boldsymbol{\nu\nu + \text{h.c.}}$}} \\ \hline
        $\mathcal{O}_{\nu\nu}$ & $(\nu_{Lp}^T C \nu_{Lr})$
    \end{tabular}
    \hspace{0.3em}
    \begin{tabular}[t]{c|c}
        \multicolumn{2}{c}{\textbf{$\boldsymbol{(\nu\nu)X + \text{h.c.}}$}} \\ \hline
        \cellcolor{gray!25}$\mathcal{O}_{\nu\gamma}$ & \cellcolor{gray!25}$(\nu_{Lp}^T C \sigma^{\mu\nu} \nu_{Lr}) F_{\mu\nu}$
    \end{tabular}
    \hspace{0.3em}
    \begin{tabular}[t]{c|c}
        \multicolumn{2}{c}{\textbf{$\boldsymbol{(\bar{L}R)X + \text{h.c.}}$}} \\ \hline
        \cellcolor{gray!25}$\mathcal{O}_{e\gamma}$ & \cellcolor{gray!25}$\bar{e}_{Lp} \sigma^{\mu\nu} e_{Rr} F_{\mu\nu}$ \\
        $\mathcal{O}_{u\gamma}$ & $\bar{u}_{Lp} \sigma^{\mu\nu} u_{Rr} F_{\mu\nu}$ \\
        $\mathcal{O}_{d\gamma}$ & $\bar{d}_{Lp} \sigma^{\mu\nu} d_{Rr} F_{\mu\nu}$ \\
        $\mathcal{O}_{uG}$ & $\bar{u}_{Lp} \sigma^{\mu\nu} T^A u_{Rr} G_{\mu\nu}^A$ \\
        $\mathcal{O}_{dG}$ & $\bar{d}_{Lp} \sigma^{\mu\nu} T^A d_{Rr} G_{\mu\nu}^A$
    \end{tabular}
    \hspace{0.3em}
    \begin{tabular}[t]{c|c}
        \multicolumn{2}{c}{\textbf{$\boldsymbol{X^3}$}} \\ \hline
        $\mathcal{O}_G$ & $f^{ABC} G_\mu^{A\nu} G_\nu^{B\rho} G_\rho^{C\mu}$ \\
        $\mathcal{O}_{\widetilde{G}}$ & $f^{ABC} \widetilde{G}_\mu^{A\nu} G_\nu^{B\rho} G_\rho^{C\mu}$
    \end{tabular}
    } % Close first resizebox
    
    \vspace{2em}
    
    % ================= BOTTOM ROW (Massive Columns) =================
    \scalebox{0.67}{%
    \begin{tabular}[t]{c|c}
        \multicolumn{2}{c}{\textbf{$\boldsymbol{(\bar{L}L)(\bar{L}L)}$}} \\ \hline
        $\mathcal{O}_{\nu\nu}^{V,LL}$ & $(\bar{\nu}_{Lp} \gamma^\mu \nu_{Lr})(\bar{\nu}_{Ls} \gamma_\mu \nu_{Lt})$ \\
        $\mathcal{O}_{ee}^{V,LL}$ & $(\bar{e}_{Lp} \gamma^\mu e_{Lr})(\bar{e}_{Ls} \gamma_\mu e_{Lt})$ \\
        \cellcolor{gray!25}$\mathcal{O}_{\nu e}^{V,LL}$ & \cellcolor{gray!25}$(\bar{\nu}_{Lp} \gamma^\mu \nu_{Lr})(\bar{e}_{Ls} \gamma_\mu e_{Lt})$ \\
        $\mathcal{O}_{\nu u}^{V,LL}$ & $(\bar{\nu}_{Lp} \gamma^\mu \nu_{Lr})(\bar{u}_{Ls} \gamma_\mu u_{Lt})$ \\
        $\mathcal{O}_{\nu d}^{V,LL}$ & $(\bar{\nu}_{Lp} \gamma^\mu \nu_{Lr})(\bar{d}_{Ls} \gamma_\mu d_{Lt})$ \\
        $\mathcal{O}_{eu}^{V,LL}$ & $(\bar{e}_{Lp} \gamma^\mu e_{Lr})(\bar{u}_{Ls} \gamma_\mu u_{Lt})$ \\
        $\mathcal{O}_{ed}^{V,LL}$ & $(\bar{e}_{Lp} \gamma^\mu e_{Lr})(\bar{d}_{Ls} \gamma_\mu d_{Lt})$ \\
        $\mathcal{O}_{\nu edu}^{V,LL}$ & $(\bar{\nu}_{Lp} \gamma^\mu e_{Lr})(\bar{d}_{Ls} \gamma_\mu u_{Lt}) + \text{h.c.}$ \\
        $\mathcal{O}_{uu}^{V,LL}$ & $(\bar{u}_{Lp} \gamma^\mu u_{Lr})(\bar{u}_{Ls} \gamma_\mu u_{Lt})$ \\
        $\mathcal{O}_{dd}^{V,LL}$ & $(\bar{d}_{Lp} \gamma^\mu d_{Lr})(\bar{d}_{Ls} \gamma_\mu d_{Lt})$ \\
        $\mathcal{O}_{ud}^{V1,LL}$ & $(\bar{u}_{Lp} \gamma^\mu u_{Lr})(\bar{d}_{Ls} \gamma_\mu d_{Lt})$ \\
        $\mathcal{O}_{ud}^{V8,LL}$ & $(\bar{u}_{Lp} \gamma^\mu T^A u_{Lr})(\bar{d}_{Ls} \gamma_\mu T^A d_{Lt})$ \\
        \multicolumn{2}{c}{\vspace{-0.5em}} \\ % Small spacer
        \multicolumn{2}{c}{\textbf{$\boldsymbol{(\bar{R}R)(\bar{R}R)}$}} \\ \hline
        $\mathcal{O}_{ee}^{V,RR}$ & $(\bar{e}_{Rp} \gamma^\mu e_{Rr})(\bar{e}_{Rs} \gamma_\mu e_{Rt})$ \\
        $\mathcal{O}_{eu}^{V,RR}$ & $(\bar{e}_{Rp} \gamma^\mu e_{Rr})(\bar{u}_{Rs} \gamma_\mu u_{Rt})$ \\
        $\mathcal{O}_{ed}^{V,RR}$ & $(\bar{e}_{Rp} \gamma^\mu e_{Rr})(\bar{d}_{Rs} \gamma_\mu d_{Rt})$ \\
        $\mathcal{O}_{uu}^{V,RR}$ & $(\bar{u}_{Rp} \gamma^\mu u_{Rr})(\bar{u}_{Rs} \gamma_\mu u_{Rt})$ \\
        $\mathcal{O}_{dd}^{V,RR}$ & $(\bar{d}_{Rp} \gamma^\mu d_{Rr})(\bar{d}_{Rs} \gamma_\mu d_{Rt})$ \\
        $\mathcal{O}_{ud}^{V1,RR}$ & $(\bar{u}_{Rp} \gamma^\mu u_{Rr})(\bar{d}_{Rs} \gamma_\mu d_{Rt})$ \\
        $\mathcal{O}_{ud}^{V8,RR}$ & $(\bar{u}_{Rp} \gamma^\mu T^A u_{Rr})(\bar{d}_{Rs} \gamma_\mu T^A d_{Rt})$
    \end{tabular}
    \hspace{0.3em}
    \begin{tabular}[t]{c|c}
        \multicolumn{2}{c}{\textbf{$\boldsymbol{(\bar{L}L)(\bar{R}R)}$}} \\ \hline
        \cellcolor{gray!25}$\mathcal{O}_{\nu e}^{V,LR}$ & \cellcolor{gray!25}$(\bar{\nu}_{Lp} \gamma^\mu \nu_{Lr})(\bar{e}_{Rs} \gamma_\mu e_{Rt})$ \\
        $\mathcal{O}_{ee}^{V,LR}$ & $(\bar{e}_{Lp} \gamma^\mu e_{Lr})(\bar{e}_{Rs} \gamma_\mu e_{Rt})$ \\
        $\mathcal{O}_{\nu u}^{V,LR}$ & $(\bar{\nu}_{Lp} \gamma^\mu \nu_{Lr})(\bar{u}_{Rs} \gamma_\mu u_{Rt})$ \\
        $\mathcal{O}_{\nu d}^{V,LR}$ & $(\bar{\nu}_{Lp} \gamma^\mu \nu_{Lr})(\bar{d}_{Rs} \gamma_\mu d_{Rt})$ \\
        $\mathcal{O}_{eu}^{V,LR}$ & $(\bar{e}_{Lp} \gamma^\mu e_{Lr})(\bar{u}_{Rs} \gamma_\mu u_{Rt})$ \\
        $\mathcal{O}_{ed}^{V,LR}$ & $(\bar{e}_{Lp} \gamma^\mu e_{Lr})(\bar{d}_{Rs} \gamma_\mu d_{Rt})$ \\
        $\mathcal{O}_{ue}^{V,LR}$ & $(\bar{u}_{Lp} \gamma^\mu u_{Lr})(\bar{e}_{Rs} \gamma_\mu e_{Rt})$ \\
        $\mathcal{O}_{de}^{V,LR}$ & $(\bar{d}_{Lp} \gamma^\mu d_{Lr})(\bar{e}_{Rs} \gamma_\mu e_{Rt})$ \\
        $\mathcal{O}_{\nu edu}^{V,LR}$ & $(\bar{\nu}_{Lp} \gamma^\mu e_{Lr})(\bar{d}_{Rs} \gamma_\mu u_{Rt}) + \text{h.c.}$ \\
        $\mathcal{O}_{uu}^{V1,LR}$ & $(\bar{u}_{Lp} \gamma^\mu u_{Lr})(\bar{u}_{Rs} \gamma_\mu u_{Rt})$ \\
        $\mathcal{O}_{uu}^{V8,LR}$ & $(\bar{u}_{Lp} \gamma^\mu T^A u_{Lr})(\bar{u}_{Rs} \gamma_\mu T^A u_{Rt})$ \\
        $\mathcal{O}_{ud}^{V1,LR}$ & $(\bar{u}_{Lp} \gamma^\mu u_{Lr})(\bar{d}_{Rs} \gamma_\mu d_{Rt})$ \\
        $\mathcal{O}_{ud}^{V8,LR}$ & $(\bar{u}_{Lp} \gamma^\mu T^A u_{Lr})(\bar{d}_{Rs} \gamma_\mu T^A d_{Rt})$ \\
        $\mathcal{O}_{du}^{V1,LR}$ & $(\bar{d}_{Lp} \gamma^\mu d_{Lr})(\bar{u}_{Rs} \gamma_\mu u_{Rt})$ \\
        $\mathcal{O}_{du}^{V8,LR}$ & $(\bar{d}_{Lp} \gamma^\mu T^A d_{Lr})(\bar{u}_{Rs} \gamma_\mu T^A u_{Rt})$ \\
        $\mathcal{O}_{dd}^{V1,LR}$ & $(\bar{d}_{Lp} \gamma^\mu d_{Lr})(\bar{d}_{Rs} \gamma_\mu d_{Rt})$ \\
        $\mathcal{O}_{dd}^{V8,LR}$ & $(\bar{d}_{Lp} \gamma^\mu T^A d_{Lr})(\bar{d}_{Rs} \gamma_\mu T^A d_{Rt})$ \\
        $\mathcal{O}_{uddu}^{V1,LR}$ & $(\bar{u}_{Lp} \gamma^\mu d_{Lr})(\bar{d}_{Rs} \gamma_\mu u_{Rt}) + \text{h.c.}$ \\
        $\mathcal{O}_{uddu}^{V8,LR}$ & $(\bar{u}_{Lp} \gamma^\mu T^A d_{Lr})(\bar{d}_{Rs} \gamma_\mu T^A u_{Rt}) + \text{h.c.}$
    \end{tabular}
    \hspace{0.3em}
    \begin{tabular}[t]{c|c}
        \multicolumn{2}{c}{\textbf{$\boldsymbol{(\bar{L}R)(\bar{L}R) + \text{h.c.}}$}} \\ \hline
        $\mathcal{O}_{ee}^{S,RR}$ & $(\bar{e}_{Lp} e_{Rr})(\bar{e}_{Ls} e_{Rt})$ \\
        $\mathcal{O}_{eu}^{S,RR}$ & $(\bar{e}_{Lp} e_{Rr})(\bar{u}_{Ls} u_{Rt})$ \\
        $\mathcal{O}_{eu}^{T,RR}$ & $(\bar{e}_{Lp} \sigma^{\mu\nu} e_{Rr})(\bar{u}_{Ls} \sigma_{\mu\nu} u_{Rt})$ \\
        $\mathcal{O}_{ed}^{S,RR}$ & $(\bar{e}_{Lp} e_{Rr})(\bar{d}_{Ls} d_{Rt})$ \\
        $\mathcal{O}_{ed}^{T,RR}$ & $(\bar{e}_{Lp} \sigma^{\mu\nu} e_{Rr})(\bar{d}_{Ls} \sigma_{\mu\nu} d_{Rt})$ \\
        $\mathcal{O}_{\nu edu}^{S,RR}$ & $(\bar{\nu}_{Lp} e_{Rr})(\bar{d}_{Ls} u_{Rt})$ \\
        $\mathcal{O}_{\nu edu}^{T,RR}$ & $(\bar{\nu}_{Lp} \sigma^{\mu\nu} e_{Rr})(\bar{d}_{Ls} \sigma_{\mu\nu} u_{Rt})$ \\
        $\mathcal{O}_{uu}^{S1,RR}$ & $(\bar{u}_{Lp} u_{Rr})(\bar{u}_{Ls} u_{Rt})$ \\
        $\mathcal{O}_{uu}^{S8,RR}$ & $(\bar{u}_{Lp} T^A u_{Rr})(\bar{u}_{Ls} T^A u_{Rt})$ \\
        $\mathcal{O}_{ud}^{S1,RR}$ & $(\bar{u}_{Lp} u_{Rr})(\bar{d}_{Ls} d_{Rt})$ \\
        $\mathcal{O}_{ud}^{S8,RR}$ & $(\bar{u}_{Lp} T^A u_{Rr})(\bar{d}_{Ls} T^A d_{Rt})$ \\
        $\mathcal{O}_{dd}^{S1,RR}$ & $(\bar{d}_{Lp} d_{Rr})(\bar{d}_{Ls} d_{Rt})$ \\
        $\mathcal{O}_{dd}^{S8,RR}$ & $(\bar{d}_{Lp} T^A d_{Rr})(\bar{d}_{Ls} T^A d_{Rt})$ \\
        $\mathcal{O}_{uddu}^{S1,RR}$ & $(\bar{u}_{Lp} d_{Rr})(\bar{d}_{Ls} u_{Rt})$ \\
        $\mathcal{O}_{uddu}^{S8,RR}$ & $(\bar{u}_{Lp} T^A d_{Rr})(\bar{d}_{Ls} T^A u_{Rt})$ \\
        \multicolumn{2}{c}{\vspace{-0.5em}} \\ % Small spacer
        \multicolumn{2}{c}{\textbf{$\boldsymbol{(\bar{L}R)(\bar{R}L) + \text{h.c.}}$}} \\ \hline
        $\mathcal{O}_{eu}^{S,RL}$ & $(\bar{e}_{Lp} e_{Rr})(\bar{u}_{Rs} u_{Lt})$ \\
        $\mathcal{O}_{ed}^{S,RL}$ & $(\bar{e}_{Lp} e_{Rr})(\bar{d}_{Rs} d_{Lt})$ \\
        $\mathcal{O}_{\nu edu}^{S,RL}$ & $(\bar{\nu}_{Lp} e_{Rr})(\bar{d}_{Rs} u_{Lt})$
    \end{tabular}
    } % Close second resizebox
    \caption{LEFT operators of dimension three, five, and six that conserve baryon and lepton
    number, and the dimension-three and dimension-five $\Delta L = \pm 2$ operators - all tables from \parencite[p.32]{jenkins2018lowOM}}
    \label{tab:LEFTbasisLNC}
\end{table}
    
\begin{table}[htbp]
    \centering
    \renewcommand{\arraystretch}{1.5}
    
    % Top standalone table
    \scalebox{0.73}{
    \begin{tabular}{c|c}
        \multicolumn{2}{c}{\textbf{$\boldsymbol{\Delta L = 4 + \text{h.c.}}$}} \\ \hline
        $\mathcal{O}_{\nu\nu}^{S,LL}$ & $(\nu_{Lp}^T C \nu_{Lr})(\nu_{Ls}^T C \nu_{Lt})$
    \end{tabular}
    }
    \vspace{1em}
    
    \scalebox{0.73}{%
    % Bottom three columns
    \begin{tabular}[t]{c|l}
        \multicolumn{2}{c}{\textbf{$\boldsymbol{\Delta L = 2 + \text{h.c.}}$}} \\ \hline
        \cellcolor{gray!25}$\mathcal{O}_{\nu e}^{S,LL}$ & \cellcolor{gray!25}$(\nu_{Lp}^T C \nu_{Lr})(\bar{e}_{Rs} e_{Lt})$ \\
        \cellcolor{gray!25}$\mathcal{O}_{\nu e}^{T,LL}$ & \cellcolor{gray!25}$(\nu_{Lp}^T C \sigma^{\mu\nu} \nu_{Lr})(\bar{e}_{Rs} \sigma_{\mu\nu} e_{Lt})$ \\
        \cellcolor{gray!25}$\mathcal{O}_{\nu e}^{S,LR}$ & \cellcolor{gray!25}$(\nu_{Lp}^T C \nu_{Lr})(\bar{e}_{Ls} e_{Rt})$ \\
        $\mathcal{O}_{\nu u}^{S,LL}$ & $(\nu_{Lp}^T C \nu_{Lr})(\bar{u}_{Rs} u_{Lt})$ \\
        $\mathcal{O}_{\nu u}^{T,LL}$ & $(\nu_{Lp}^T C \sigma^{\mu\nu} \nu_{Lr})(\bar{u}_{Rs} \sigma_{\mu\nu} u_{Lt})$ \\
        $\mathcal{O}_{\nu u}^{S,LR}$ & $(\nu_{Lp}^T C \nu_{Lr})(\bar{u}_{Ls} u_{Rt})$ \\
        $\mathcal{O}_{\nu d}^{S,LL}$ & $(\nu_{Lp}^T C \nu_{Lr})(\bar{d}_{Rs} d_{Lt})$ \\
        $\mathcal{O}_{\nu d}^{T,LL}$ & $(\nu_{Lp}^T C \sigma^{\mu\nu} \nu_{Lr})(\bar{d}_{Rs} \sigma_{\mu\nu} d_{Lt})$ \\
        $\mathcal{O}_{\nu d}^{S,LR}$ & $(\nu_{Lp}^T C \nu_{Lr})(\bar{d}_{Ls} d_{Rt})$ \\
        $\mathcal{O}_{\nu edu}^{S,LL}$ & $(\nu_{Lp}^T C e_{Lr})(\bar{d}_{Rs} u_{Lt})$ \\
        $\mathcal{O}_{\nu edu}^{T,LL}$ & $(\nu_{Lp}^T C \sigma^{\mu\nu} e_{Lr})(\bar{d}_{Rs} \sigma_{\mu\nu} u_{Lt})$ \\
        $\mathcal{O}_{\nu edu}^{S,LR}$ & $(\nu_{Lp}^T C e_{Lr})(\bar{d}_{Ls} u_{Rt})$ \\
        $\mathcal{O}_{\nu edu}^{V,RL}$ & $(\nu_{Lp}^T C \gamma^\mu e_{Rr})(\bar{d}_{Ls} \gamma_\mu u_{Lt})$ \\
        $\mathcal{O}_{\nu edu}^{V,RR}$ & $(\nu_{Lp}^T C \gamma^\mu e_{Rr})(\bar{d}_{Rs} \gamma_\mu u_{Rt})$
    \end{tabular}
    \hspace{0.3em}
    \begin{tabular}[t]{c|l}
        \multicolumn{2}{c}{\textbf{$\boldsymbol{\Delta B = \Delta L = 1 + \text{h.c.}}$}} \\ \hline
        $\mathcal{O}_{udd}^{S,LL}$ & $\epsilon_{\alpha\beta\gamma}(u_{Lp}^{\alpha T} C d_{Lr}^\beta)(d_{Ls}^{\gamma T} C \nu_{Lt})$ \\
        $\mathcal{O}_{duu}^{S,LL}$ & $\epsilon_{\alpha\beta\gamma}(d_{Lp}^{\alpha T} C u_{Lr}^\beta)(u_{Ls}^{\gamma T} C e_{Lt})$ \\
        $\mathcal{O}_{uud}^{S,LR}$ & $\epsilon_{\alpha\beta\gamma}(u_{Lp}^{\alpha T} C u_{Lr}^\beta)(d_{Rs}^{\gamma T} C e_{Rt})$ \\
        $\mathcal{O}_{duu}^{S,LR}$ & $\epsilon_{\alpha\beta\gamma}(d_{Lp}^{\alpha T} C u_{Lr}^\beta)(u_{Rs}^{\gamma T} C e_{Rt})$ \\
        $\mathcal{O}_{uud}^{S,RL}$ & $\epsilon_{\alpha\beta\gamma}(u_{Rp}^{\alpha T} C u_{Rr}^\beta)(d_{Ls}^{\gamma T} C e_{Lt})$ \\
        $\mathcal{O}_{duu}^{S,RL}$ & $\epsilon_{\alpha\beta\gamma}(d_{Rp}^{\alpha T} C u_{Rr}^\beta)(u_{Ls}^{\gamma T} C e_{Lt})$ \\
        $\mathcal{O}_{dud}^{S,RL}$ & $\epsilon_{\alpha\beta\gamma}(d_{Rp}^{\alpha T} C u_{Rr}^\beta)(d_{Ls}^{\gamma T} C \nu_{Lt})$ \\
        $\mathcal{O}_{ddu}^{S,RL}$ & $\epsilon_{\alpha\beta\gamma}(d_{Rp}^{\alpha T} C d_{Rr}^\beta)(u_{Ls}^{\gamma T} C \nu_{Lt})$ \\
        $\mathcal{O}_{duu}^{S,RR}$ & $\epsilon_{\alpha\beta\gamma}(d_{Rp}^{\alpha T} C u_{Rr}^\beta)(u_{Rs}^{\gamma T} C e_{Rt})$
    \end{tabular}
    \hspace{0.3em}
    \begin{tabular}[t]{c|l}
        \multicolumn{2}{c}{\textbf{$\boldsymbol{\Delta B = -\Delta L = 1 + \text{h.c.}}$}} \\ \hline
        $\mathcal{O}_{ddd}^{S,LL}$ & $\epsilon_{\alpha\beta\gamma}(d_{Lp}^{\alpha T} C d_{Lr}^\beta)(\bar{e}_{Rs} d_{Lt}^\gamma)$ \\
        $\mathcal{O}_{udd}^{S,LR}$ & $\epsilon_{\alpha\beta\gamma}(u_{Lp}^{\alpha T} C d_{Lr}^\beta)(\bar{\nu}_{Ls} d_{Rt}^\gamma)$ \\
        $\mathcal{O}_{ddu}^{S,LR}$ & $\epsilon_{\alpha\beta\gamma}(d_{Lp}^{\alpha T} C d_{Lr}^\beta)(\bar{\nu}_{Ls} u_{Rt}^\gamma)$ \\
        $\mathcal{O}_{ddd}^{S,LR}$ & $\epsilon_{\alpha\beta\gamma}(d_{Lp}^{\alpha T} C d_{Lr}^\beta)(\bar{e}_{Ls} d_{Rt}^\gamma)$ \\
        $\mathcal{O}_{ddd}^{S,RL}$ & $\epsilon_{\alpha\beta\gamma}(d_{Rp}^{\alpha T} C d_{Rr}^\beta)(\bar{e}_{Rs} d_{Lt}^\gamma)$ \\
        $\mathcal{O}_{udd}^{S,RR}$ & $\epsilon_{\alpha\beta\gamma}(u_{Rp}^{\alpha T} C d_{Rr}^\beta)(\bar{\nu}_{Ls} d_{Rt}^\gamma)$ \\
        $\mathcal{O}_{ddd}^{S,RR}$ & $\epsilon_{\alpha\beta\gamma}(d_{Rp}^{\alpha T} C d_{Rr}^\beta)(\bar{e}_{Ls} d_{Rt}^\gamma)$
    \end{tabular}
    } % Close resizebox
    \caption{LEFT dimension-six four-fermion operators that violate baryon and/or lepton number - all tables from \parencite[p.33]{jenkins2018lowOM}}
    \label{tab:LEFTbasisLNV}
\end{table}


\section{SMEFT basis}
See table \ref{tab:dim6SMEFTbasis}
\begin{table}[htbp]
    \centering
    \renewcommand{\arraystretch}{1.5}
    \scalebox{0.75}{%
    
    % ================= ROW 1 =================
    \begin{tabular}[t]{c|c}
        \multicolumn{2}{c}{\textbf{1 : $\boldsymbol{X^3}$}} \\ \hline
        $Q_G$ & $f^{ABC} G_\mu^{A\nu} G_\nu^{B\rho} G_\rho^{C\mu}$ \\
        $Q_{\widetilde{G}}$ & $f^{ABC} \widetilde{G}_\mu^{A\nu} G_\nu^{B\rho} G_\rho^{C\mu}$ \\
        $Q_W$ & $\epsilon^{IJK} W_\mu^{I\nu} W_\nu^{J\rho} W_\rho^{K\mu}$ \\
        $Q_{\widetilde{W}}$ & $\epsilon^{IJK} \widetilde{W}_\mu^{I\nu} W_\nu^{J\rho} W_\rho^{K\mu}$
    \end{tabular}
    \hspace{0.3em}
    \begin{tabular}[t]{c|c}
        \multicolumn{2}{c}{\textbf{2 : $\boldsymbol{H^6}$}} \\ \hline
        $Q_H$ & $(H^\dagger H)^3$ \\
    \end{tabular}
    \hspace{0.3em}
    \begin{tabular}[t]{c|c}
        \multicolumn{2}{c}{\textbf{3 : $\boldsymbol{H^4 D^2}$}} \\ \hline
        $Q_{H\square}$ & $(H^\dagger H)\square(H^\dagger H)$ \\
        $Q_{HD}$ & $\left(H^\dagger D^\mu H\right)^* \left(H^\dagger D_\mu H\right)$ \\
    \end{tabular}
    \hspace{0.3em}
    \begin{tabular}[t]{c|c}
        \multicolumn{2}{c}{\textbf{5 : $\boldsymbol{\psi^2 H^3 + \text{h.c.}}$}} \\ \hline
        $Q_{eH}$ & $(H^\dagger H)(\bar{l}_p e_r H)$ \\
        $Q_{uH}$ & $(H^\dagger H)(\bar{q}_p u_r \widetilde{H})$ \\
        $Q_{dH}$ & $(H^\dagger H)(\bar{q}_p d_r H)$ \\
    \end{tabular}
    } % Close first resizebox to allow line break
    
    \vspace{1em}
    \scalebox{0.75}{%
    % ================= ROW 2 =================
    \begin{tabular}[t]{c|c}
        \multicolumn{2}{c}{\textbf{4 : $\boldsymbol{X^2 H^2}$}} \\ \hline
        $Q_{HG}$ & $H^\dagger H\, G_{\mu\nu}^A G^{A\mu\nu}$ \\
        $Q_{H\widetilde{G}}$ & $H^\dagger H\, \widetilde{G}_{\mu\nu}^A G^{A\mu\nu}$ \\
        $Q_{HW}$ & $H^\dagger H\, W_{\mu\nu}^I W^{I\mu\nu}$ \\
        $Q_{H\widetilde{W}}$ & $H^\dagger H\, \widetilde{W}_{\mu\nu}^I W^{I\mu\nu}$ \\
        $Q_{HB}$ & $H^\dagger H\, B_{\mu\nu} B^{\mu\nu}$ \\
        $Q_{H\widetilde{B}}$ & $H^\dagger H\, \widetilde{B}_{\mu\nu} B^{\mu\nu}$ \\
        $Q_{HWB}$ & $H^\dagger \tau^I H\, W_{\mu\nu}^I B^{\mu\nu}$ \\
        $Q_{H\widetilde{W}B}$ & $H^\dagger \tau^I H\, \widetilde{W}_{\mu\nu}^I B^{\mu\nu}$
    \end{tabular}
    \hspace{0.3em}
    \begin{tabular}[t]{c|c}
        \multicolumn{2}{c}{\textbf{6 : $\boldsymbol{\psi^2 XH + \text{h.c.}}$}} \\ \hline
        $Q_{eW}$ & $(\bar{l}_p \sigma^{\mu\nu} e_r) \tau^I H W_{\mu\nu}^I$ \\
        $Q_{eB}$ & $(\bar{l}_p \sigma^{\mu\nu} e_r) H B_{\mu\nu}$ \\
        $Q_{uG}$ & $(\bar{q}_p \sigma^{\mu\nu} T^A u_r) \widetilde{H} G_{\mu\nu}^A$ \\
        $Q_{uW}$ & $(\bar{q}_p \sigma^{\mu\nu} \tau^I u_r) \widetilde{H} W_{\mu\nu}^I$ \\
        $Q_{uB}$ & $(\bar{q}_p \sigma^{\mu\nu} u_r) \widetilde{H} B_{\mu\nu}$ \\
        $Q_{dG}$ & $(\bar{q}_p \sigma^{\mu\nu} T^A d_r) H G_{\mu\nu}^A$ \\
        $Q_{dW}$ & $(\bar{q}_p \sigma^{\mu\nu} \tau^I d_r) H W_{\mu\nu}^I$ \\
        $Q_{dB}$ & $(\bar{q}_p \sigma^{\mu\nu} d_r) H B_{\mu\nu}$
    \end{tabular}
    \hspace{0.3em}
    \begin{tabular}[t]{c|c}
        \multicolumn{2}{c}{\textbf{7 : $\boldsymbol{\psi^2 H^2 D}$}} \\ \hline
        \cellcolor{gray!25}$Q_{Hl}^{(1)}$ & \cellcolor{gray!25}$(H^\dagger i\overleftrightarrow{D}_\mu H)(\bar{l}_p \gamma^\mu l_r)$ \\
        \cellcolor{gray!25}$Q_{Hl}^{(3)}$ & \cellcolor{gray!25}$(H^\dagger i\overleftrightarrow{D}^I_\mu H)(\bar{l}_p \tau^I \gamma^\mu l_r)$ \\
        \cellcolor{gray!25}$Q_{He}$ & \cellcolor{gray!25}$(H^\dagger i\overleftrightarrow{D}_\mu H)(\bar{e}_p \gamma^\mu e_r)$ \\
        $Q_{Hq}^{(1)}$ & $(H^\dagger i\overleftrightarrow{D}_\mu H)(\bar{q}_p \gamma^\mu q_r)$ \\
        $Q_{Hq}^{(3)}$ & $(H^\dagger i\overleftrightarrow{D}^I_\mu H)(\bar{q}_p \tau^I \gamma^\mu q_r)$ \\
        $Q_{Hu}$ & $(H^\dagger i\overleftrightarrow{D}_\mu H)(\bar{u}_p \gamma^\mu u_r)$ \\
        $Q_{Hd}$ & $(H^\dagger i\overleftrightarrow{D}_\mu H)(\bar{d}_p \gamma^\mu d_r)$ \\
        $Q_{Hud} + \text{h.c.}$ & $i(\widetilde{H}^\dagger D_\mu H)(\bar{u}_p \gamma^\mu d_r)$
    \end{tabular}
    } % Close second resizebox
    
    \vspace{1em}
    \scalebox{0.75}{%
    % ================= ROW 3 =================
    \begin{tabular}[t]{c|c}
        \multicolumn{2}{c}{\textbf{8 : $\boldsymbol{(\bar{L}L)(\bar{L}L)}$}} \\ \hline
        \cellcolor{gray!25}$Q_{ll}$ & \cellcolor{gray!25}$(\bar{l}_p \gamma^\mu l_r)(\bar{l}_s \gamma_\mu l_t)$ \\
        $Q_{qq}^{(1)}$ & $(\bar{q}_p \gamma^\mu q_r)(\bar{q}_s \gamma_\mu q_t)$ \\
        $Q_{qq}^{(3)}$ & $(\bar{q}_p \gamma^\mu \tau^I q_r)(\bar{q}_s \gamma_\mu \tau^I q_t)$ \\
        $Q_{lq}^{(1)}$ & $(\bar{l}_p \gamma^\mu l_r)(\bar{q}_s \gamma_\mu q_t)$ \\
        $Q_{lq}^{(3)}$ & $(\bar{l}_p \gamma^\mu \tau^I l_r)(\bar{q}_s \gamma_\mu \tau^I q_t)$
    \end{tabular}
    \hspace{0.3em}
    \begin{tabular}[t]{c|c}
        \multicolumn{2}{c}{\textbf{8 : $\boldsymbol{(\bar{R}R)(\bar{R}R)}$}} \\ \hline
        $Q_{ee}$ & $(\bar{e}_p \gamma^\mu e_r)(\bar{e}_s \gamma_\mu e_t)$ \\
        $Q_{uu}$ & $(\bar{u}_p \gamma^\mu u_r)(\bar{u}_s \gamma_\mu u_t)$ \\
        $Q_{dd}$ & $(\bar{d}_p \gamma^\mu d_r)(\bar{d}_s \gamma_\mu d_t)$ \\
        $Q_{eu}$ & $(\bar{e}_p \gamma^\mu e_r)(\bar{u}_s \gamma_\mu u_t)$ \\
        $Q_{ed}$ & $(\bar{e}_p \gamma^\mu e_r)(\bar{d}_s \gamma_\mu d_t)$ \\
        $Q_{ud}^{(1)}$ & $(\bar{u}_p \gamma^\mu u_r)(\bar{d}_s \gamma_\mu d_t)$ \\
        $Q_{ud}^{(8)}$ & $(\bar{u}_p \gamma^\mu T^A u_r)(\bar{d}_s \gamma_\mu T^A d_t)$
    \end{tabular}
    \hspace{0.3em}
    \begin{tabular}[t]{c|c}
        \multicolumn{2}{c}{\textbf{8 : $\boldsymbol{(\bar{L}L)(\bar{R}R)}$}} \\ \hline
        \cellcolor{gray!25}$Q_{le}$ & \cellcolor{gray!25}$(\bar{l}_p \gamma^\mu l_r)(\bar{e}_s \gamma_\mu e_t)$ \\
        $Q_{lu}$ & $(\bar{l}_p \gamma^\mu l_r)(\bar{u}_s \gamma_\mu u_t)$ \\
        $Q_{ld}$ & $(\bar{l}_p \gamma^\mu l_r)(\bar{d}_s \gamma_\mu d_t)$ \\
        $Q_{qe}$ & $(\bar{q}_p \gamma^\mu q_r)(\bar{e}_s \gamma_\mu e_t)$ \\
        $Q_{qu}^{(1)}$ & $(\bar{q}_p \gamma^\mu q_r)(\bar{u}_s \gamma_\mu u_t)$ \\
        $Q_{qu}^{(8)}$ & $(\bar{q}_p \gamma^\mu T^A q_r)(\bar{u}_s \gamma_\mu T^A u_t)$ \\
        $Q_{qd}^{(1)}$ & $(\bar{q}_p \gamma^\mu q_r)(\bar{d}_s \gamma_\mu d_t)$ \\
        $Q_{qd}^{(8)}$ & $(\bar{q}_p \gamma^\mu T^A q_r)(\bar{d}_s \gamma_\mu T^A d_t)$
    \end{tabular}
    } % Close third resizebox
    
    \vspace{1em}
    \scalebox{0.75}{% Scaling down slightly less for the smaller last row
    % ================= ROW 4 =================
    \begin{tabular}[t]{c|c}
        \multicolumn{2}{c}{\textbf{8 : $\boldsymbol{(\bar{L}R)(\bar{R}L) + \text{h.c.}}$}} \\ \hline
        $Q_{ledq}$ & $(\bar{l}_p^j e_r)(\bar{d}_s q_{tj})$
    \end{tabular}
    \hspace{0.3em}
    \begin{tabular}[t]{c|c}
        \multicolumn{2}{c}{\textbf{8 : $\boldsymbol{(\bar{L}R)(\bar{L}R) + \text{h.c.}}$}} \\ \hline
        $Q_{quqd}^{(1)}$ & $(\bar{q}_p^j u_r) \epsilon_{jk} (\bar{q}_s^k d_t)$ \\
        $Q_{quqd}^{(8)}$ & $(\bar{q}_p^j T^A u_r) \epsilon_{jk} (\bar{q}_s^k T^A d_t)$ \\
        $Q_{lequ}^{(1)}$ & $(\bar{l}_p^j e_r) \epsilon_{jk} (\bar{q}_s^k u_t)$ \\
        $Q_{lequ}^{(3)}$ & $(\bar{l}_p^j \sigma_{\mu\nu} e_r) \epsilon_{jk} (\bar{q}_s^k \sigma^{\mu\nu} u_t)$
    \end{tabular}
    } % Close fourth resizebox
    \caption{76 operators from the Warsaw basis \parencite{grzadkowski2010dimension}, that conserve baryon and lepton number $\to$ table from \parencite{jenkins2018lowOM}.}
    \label{tab:dim6SMEFTbasis}
\end{table}
\subsection{SMEFT dim-7 operator basis}
See table \ref{tab:dim7SMEFTbl02} and \ref{tab:dim7SMEFTbl1-1}.
\begin{table}[htbp]
    \centering
    \renewcommand{\arraystretch}{1.5}
    
    % ================= ROW 1 =================
    \scalebox{0.75}{%
    \begin{tabular}[t]{c|c}
        \multicolumn{2}{c}{\textbf{$\boldsymbol{\psi^2 H^4}$}} \\ \hline
        $\mathcal{O}_{LH}^{pr}$ & $\epsilon_{ij}\epsilon_{mn} (\overline{L_p^{iC}} L_r^m) H^j H^n (H^\dagger H)$
    \end{tabular}
    \hspace{0.3em}
    \begin{tabular}[t]{c|c}
        \multicolumn{2}{c}{\textbf{$\boldsymbol{\psi^2 H^3 D}$}} \\ \hline
        \cellcolor{gray!25}$\mathcal{O}_{LeHD}^{pr}$ & \cellcolor{gray!25}$\epsilon_{ij}\epsilon_{mn} (\overline{L_p^{iC}} \gamma_\mu e_r) H^j H^m iD^\mu H^n$
    \end{tabular}
    } % Close row 1
    
    \vspace{1em}
    \scalebox{0.75}{
    \begin{tabular}[t]{c|c}
        \multicolumn{2}{c}{\textbf{$\boldsymbol{\psi^2 H^2 D^2}$}} \\ \hline
        $\mathcal{O}_{LHD1}^{pr}$ & $\epsilon_{ij}\epsilon_{mn} (\overline{L_p^{iC}} D^\mu L_r^j) H^m D_\mu H^n$ \\
        $\mathcal{O}_{LHD2}^{pr}$ & $\epsilon_{im}\epsilon_{jn} (\overline{L_p^{iC}} D^\mu L_r^j) H^m D_\mu H^n$
    \end{tabular}
    \hspace{0.3em}
    \begin{tabular}[t]{c|c}
        \multicolumn{2}{c}{\textbf{$\boldsymbol{\psi^2 H^2 X}$}} \\ \hline
        $\mathcal{O}_{LHB}^{pr}$ & $\epsilon_{ij}\epsilon_{mn} (\overline{L_p^{iC}} \sigma_{\mu\nu} L_r^m) H^j H^n B^{\mu\nu}$ \\
        $\mathcal{O}_{LHW}^{pr}$ & $\epsilon_{ij}(\epsilon\tau^I)_{mn} (\overline{L_p^{iC}} \sigma_{\mu\nu} L_r^m) H^j H^n W^{I\mu\nu}$
    \end{tabular}
    }
    
    % ================= ROW 2 =================
    \vspace{1em}
    \scalebox{0.75}{%
    \begin{tabular}[t]{c|c}
        \multicolumn{2}{c}{\textbf{$\boldsymbol{\psi^4 D}$}} \\ \hline
        $\mathcal{O}_{\bar{d}uLLD}^{prst}$ & $\epsilon_{ij} (\overline{d_p} \gamma_\mu u_r)(\overline{L_s^{iC}} iD^\mu L_t^j)$ \\
        $\mathcal{O}_{\bar{e}dddD}^{prst}$ & $\epsilon_{\alpha\beta\gamma} (\overline{e_p} \gamma_\mu d_r^\alpha) (\overline{d_s^{\beta C}} iD^\mu d_t^\gamma)$ \\
        $\mathcal{O}_{\bar{L}QddD}^{prst}$ & $\epsilon_{\alpha\beta\gamma} (\overline{L_p} \gamma_\mu Q_r^\alpha) (\overline{d_s^{\beta C}} iD^\mu d_t^\gamma)$
    \end{tabular}
    } % Close row 2
    \caption{Dim-7 SMEFT operators with $(\Delta B, \Delta L) = (0, 2)$ - all tables from \parencite[p.5]{liao2025rge}}
    \label{tab:dim7SMEFTbl02}
    \end{table}
    
\begin{table}[htbp]
    \centering
    \scalebox{0.75}{%
    \begin{tabular}[t]{c|l}
        \multicolumn{2}{c}{\textbf{$\boldsymbol{\psi^4 H \quad (\Delta L = 2)}$}} \\ \hline
        \cellcolor{gray!25}$\mathcal{O}_{\bar{e}LLLH}^{prst}$ & \cellcolor{gray!25}$\epsilon_{ij}\epsilon_{mn} (\overline{e_p} L_r^i)(\overline{L_s^{jC}} L_t^m) H^n$ \\
        $\mathcal{O}_{\bar{d}LQLH1}^{prst}$ & $\epsilon_{ij}\epsilon_{mn} (\overline{d_p} L_r^i)(\overline{Q_s^{jC}} L_t^m) H^n$ \\
        $\mathcal{O}_{\bar{d}LQLH2}^{prst}$ & $\epsilon_{im}\epsilon_{jn} (\overline{d_p} L_r^i)(\overline{Q_s^{jC}} L_t^m) H^n$ \\
        $\mathcal{O}_{\bar{d}LueH}^{prst}$ & $\epsilon_{ij} (\overline{d_p} L_r^i)(u_s^C e_t) H^j$ \\
        $\mathcal{O}_{\bar{Q}uLLH}^{prst}$ & $\epsilon_{ij} (\overline{Q_p} u_r)(\overline{L_s^C} L_t^i) H^j$
    \end{tabular}
    \hspace{0.3em}
    \begin{tabular}[t]{c|l}
        \multicolumn{2}{c}{\textbf{$\boldsymbol{\psi^4 H \quad (\Delta B = 1, \Delta L = -1)}$}} \\ \hline
        $\mathcal{O}_{\bar{L}dud\tilde{H}}^{prst}$ & $\epsilon_{\alpha\beta\gamma} (\overline{L_p} d_r^\alpha) (u_s^{\beta C} d_t^\gamma) \widetilde{H}$ \\
        $\mathcal{O}_{\bar{L}dddH}^{prst}$ & $\epsilon_{\alpha\beta\gamma} (\overline{L_p} d_r^\alpha) (\overline{d_s^{\beta C}} d_t^\gamma) H$ \\
        $\mathcal{O}_{\bar{e}Qdd\tilde{H}}^{prst}$ & $\epsilon_{ij}\epsilon_{\alpha\beta\gamma} (\overline{e_p} Q_r^{i\alpha}) (\overline{d_s^{\beta C}} d_t^\gamma) \widetilde{H}^j$ \\
        $\mathcal{O}_{\bar{L}dQQ\tilde{H}}^{prst}$ & $\epsilon_{ij}\epsilon_{\alpha\beta\gamma} (\overline{L_p} d_r^\alpha) (\overline{Q_s^{\beta C}} Q_t^{i\gamma}) \widetilde{H}^j$
    \end{tabular}
    } % Close row 3
    \caption{Dim-7 SMEFT operators with $(\Delta B, \Delta L) = (1, -1)$ - all tables from \parencite[p.5]{liao2025rge}}
    \label{tab:dim7SMEFTbl1-1}
\end{table}
\clearpage
\section{Renormalization Group Equations}\label{sec:apprge}
Here we cite the full one-loop RGE for the relevant SMEFT operators. The overall structure of the dimension-six SMEFT renormalization group equations was presented in \parencite{jenkins2013rge1}. The Yukawa loop contributions were given in \parencite{jenkins2014rge2} and the Gauge loop contributions in \parencite{alonso2014rge3} - from the latter two papers we collect the results and add them here together. The notation in the papers is the following:
\begin{equation}
    \dot{C} \equiv 16\pi^2 \mu \frac{d C}{d\mu}
\end{equation}
\begin{align}
%%%%%
\dot L_{\substack{ ll \\ prst }} &= -\frac{1}{2} [Y_e^\dagger Y_e]_{pr} L_{\substack{ Hl \\ st }}^{(1)} -\frac{1}{2} [Y_e^\dagger Y_e]_{st} L_{\substack{ Hl \\ pr }}^{(1)}  +\frac{1}{2} [Y_e^\dagger Y_e]_{pr} L_{\substack{ Hl \\ st }}^{(3)} +\frac{1}{2} [Y_e^\dagger Y_e]_{st} L_{\substack{ Hl \\ pr }}^{(3)} \nn
&-  [Y_e^\dagger Y_e]_{sr} L_{\substack{ Hl \\ pt }}^{(3)} - [Y_e^\dagger Y_e]_{pt} L_{\substack{ Hl \\ sr }}^{(3)}  - \frac{1}{2} \, [Y_e^\dagger]_{sv} \, [Y_e]_{wt} \, L_{\substack{le  \\ prvw}} - \frac{1}{2} \, [Y_e^\dagger]_{pv} \, [Y_e]_{wr} \, L_{\substack{le  \\ stvw}} \nn
&+ \gamfY lpv L_{\substack{ll \\ vrst}} + \gamfY lsv L_{\substack{ll \\ prvt}}
+  L_{\substack{ll \\ pvst}} \gamfY lvr + L_{\substack{ll \\ prsv}}  \gamfY lvt\\
&+ \frac23 g_1^2 \hyp_h \hyp_l C^{(1)}_{ \substack {Hl \\ s t } } \delta_{p r} 
-       \frac16 g_2^2 C^{(3)}_{ \substack {Hl \\ s t } } \delta_{p r}
+ \frac13 g_2^2 C^{(3)}_{ \substack {Hl \\ s r } } \delta_{p t} 
+    \frac13   g_2^2 C^{(3)}_{ \substack {Hl \\ p t } } \delta_{r s} 
\nn
&
+ \frac23 g_1^2 \hyp_h \hyp_l C^{(1)}_{ \substack {Hl \\ p r } }   \delta_{s t} 
- \frac16 g_2^2 C^{(3)}_{ \substack {Hl \\ p r } } \delta_{s t}  
+     \frac43 g_1^2 \hyp_l^2 L_{ \substack {ll \\ p r w w } } \delta_{s t}
+        \frac43 g_1^2 \hyp_l^2 L_{ \substack {ll \\ s t w w } } \delta_{p r} 
+        \frac43 g_1^2 \hyp_l^2 L_{ \substack {ll \\ w w s t } } \delta_{p r} 
\nn
&
+        \frac43 g_1^2 \hyp_l^2 L_{ \substack {ll \\ w w p r } } \delta_{s t} 
+        \frac23 g_1^2 \hyp_l^2 L_{ \substack {ll \\ p w w r } } \delta_{s t} 
+        \frac23 g_1^2 \hyp_l^2 L_{ \substack {ll \\ s w w t } } \delta_{p r} 
+        \frac23 g_1^2 \hyp_l^2 L_{ \substack {ll \\ w r p w } } \delta_{s t}
+        \frac23 g_1^2 \hyp_l^2 L_{ \substack {ll \\ w t s w } } \delta_{p r} 
\nn
&
-        \frac16 g_2^2 L_{ \substack {ll \\ p w w r } } \delta_{s t} 
-        \frac16 g_2^2 L_{ \substack {ll \\ s w w t } } \delta_{p r}
-        \frac16 g_2^2 L_{ \substack {ll \\ w r p w } } \delta_{s t} 
-        \frac16 g_2^2 L_{ \substack {ll \\ w t s w } } \delta_{p r} 
+        \frac13 g_2^2 L_{ \substack {ll \\ s w w r } } \delta_{p t}
\nn
&
+        \frac13 g_2^2 L_{ \substack {ll \\ p w w t } } \delta_{r s} 
+         \frac13    g_2^2 L_{ \substack {ll \\ w r s w } } \delta_{p t} 
+        \frac13 g_2^2 L_{ \substack {ll \\ w t p w } } \delta_{r s} 
+        \frac43 g_1^2 N_c \hyp_l \hyp_q C^{(1)}_{ \substack {lq \\ p r w w } } \delta_{s t}
+        \frac43 g_1^2 N_c \hyp_l \hyp_q C^{(1)}_{ \substack {lq \\ s t w w } } \delta_{p r} 
\nn
&
-        \frac13 g_2^2 N_c C^{(3)}_{ \substack {lq \\ p r w w } } \delta_{s t}
-        \frac13 g_2^2 N_c C^{(3)}_{ \substack {lq \\ s t w w } } \delta_{p r} 
+        \frac23 g_2^2 N_c C^{(3)}_{ \substack {lq \\ s r w w } } \delta_{p t} 
+        \frac23 g_2^2 N_c C^{(3)}_{ \substack {lq \\ p t w w } } \delta_{r s} 
\nn
&
+        \frac23 g_1^2 N_c \hyp_l \hyp_u L_{ \substack {lu \\ p r w w } } \delta_{s t}
+        \frac23 g_1^2 N_c \hyp_l \hyp_u L_{ \substack {lu \\ s t w w } } \delta_{p r} 
+        \frac23 g_1^2 N_c \hyp_d \hyp_l L_{ \substack {ld \\ p r w w } } \delta_{s t} 
+        \frac23 g_1^2 N_c \hyp_d \hyp_l L_{ \substack {ld \\ s t w w } } \delta_{p r} 
\nn
&
+        \frac23 g_1^2 \hyp_e \hyp_l L_{ \substack {le \\ p r w w } } \delta_{s t} 
+        \frac23 g_1^2 \hyp_e \hyp_l L_{ \substack {le \\ s t w w } } \delta_{p r} 
+ 6 g_2^2 \, L_{\substack{ll \\ ptsr}} - 3\left( g_2^2 - 4 \hyp_l^2 g_1^2 \right) L_{\substack{ll \\ prst}} 
\end{align}
\begin{align}
%
\dot L_{\substack{Hl  \\ pr}}^{(3)} &= -\frac12 [Y_e^\dagger Y_e]_{pr}   L_{H\Box}  +\frac12 [Y_e^\dagger Y_e]_{pt} \bigl( 3 L_{\substack{Hl \\ tr}}^{(1)} +  L_{\substack{Hl \\ tr}}^{(3)} \bigr) + \frac12 \bigl( 3 L_{\substack{Hl \\ pt}}^{(1)} +  L_{\substack{Hl \\ pt}}^{(3)} \bigr)  [Y_e^\dagger Y_e]_{tr} \nn
&
 - \bigl( L_{\substack{ll \\ ptsr}}+ L_{\substack{ll \\ srpt}} \bigr) [{Y_e}^\dagger Y_e  ]_{ts} -2N_c  C^{(3)}_{\substack{lq \\ prst }}   [Y_d^\dagger Y_d+Y_u^\dagger Y_u]_{ts}+ 2 \gamHY   L_{\substack{Hl  \\ pr}}^{(3)}+ \gamfY lpt L_{\substack{Hl\\ tr}} ^{(3)} + L_{\substack{Hl \\ pt}}^{(3)}\gamfY ltr\\
 &+\frac16 g_2^2 L_{H \Box} \delta_{pr}+ \frac23 g_2^2  C^{(3)}_{\substack{H l \\ tt}}\delta_{pr} +  \frac23 g_2^2 N_c  C^{(3)}_{\substack{H q \\ tt}}\delta_{pr}
+\frac13 g_2^2  L_{\substack{Hl \\ pr}}^{(3)} +\frac{1}{3} g_2^2 L_{ \substack {ll \\ p w w r } } +\frac{1}{3} g_2^2 L_{ \substack {ll \\ w r p w } } 
+\frac{2}{3} g_2^2 N_c C^{(3)}_{ \substack {lq \\ p r w w } }\nn
& - 6 g_2^2  L_{\substack{Hl \\ pr}}^{(3)} 
\end{align}
\begin{align}
%%%%%%%%%%%%%%%%%%%%%%%%%%%%%%%%%%%%%%%%%%%%%%%%%%%%%%%%%%%%
\dot L_{\substack{ le \\ prst }} &=  [Y_e]_{sr} \xi_{\substack{e \\ pt}} +[Y_e^\dagger]_{pt} \xi^*_{\substack{e \\ rs}}- [Y_e^\dagger Y_e]_{pr} L_{\substack{ He \\ st }} +2  [Y_e Y_e^\dagger]_{st} L_{\substack{ Hl \\ pr }}^{(1)}  - [Y_e^\dagger]_{pv} \, [Y_e]_{wr} \, L_{\substack{ee  \\ vtsw}} -  [Y_e^\dagger]_{pw} \, [Y_e]_{vr} \, L_{\substack{ee  \\ wtsv}} \nn
& - 2 [Y_e^\dagger]_{pv} \, [Y_e]_{wr} \, L_{\substack{ee  \\ vwst}}
 + [Y_e^\dagger]_{pw} \, [Y_e]_{sv}  \, L_{\substack{le  \\ vrwt}}    - [Y_e^\dagger]_{wt} \, [Y_e]_{sv}  \, L_{\substack{ll  \\ pwvr}} -  [Y_e^\dagger]_{vt} \, [Y_e]_{sw} \, L_{\substack{ll \\ pvwr}} \nn
& - 4 [Y_e^\dagger]_{wt} \, [Y_e]_{sv} \, L_{\substack{ll \\ prvw}} + [Y_e^\dagger]_{vt} \, [Y_e]_{wr}  \, L_{\substack{le \\ pvsw}}   + \gamfY lpv L_{\substack{le \\ vrst}}  + \gamfY esv  L_{\substack{le \\ prvt}}
+  L_{\substack{le \\ pvst}}  \gamfY lvr  + L_{\substack{le \\ prsv}}  \gamfY evt\\
 &+    \frac43 g_1^2 \hyp_h \hyp_l L_{ \substack {He \\ s t } } \delta_{p r}  
 +        \frac43 g_1^2 \hyp_e \hyp_h C^{(1)}_{ \substack {Hl \\ p r } } \delta_{s t}    
+        \frac83 g_1^2 \hyp_e \hyp_l L_{ \substack {ll \\ p r w w } } \delta_{s t} 
+        \frac83 g_1^2 \hyp_e \hyp_l L_{ \substack {ll \\ w w p r } } \delta_{s t} 
+        \frac43 g_1^2 \hyp_e \hyp_l L_{ \substack {ll \\ p w w r } } \delta_{s t}
\nn
&
+        \frac43 g_1^2 \hyp_e \hyp_l L_{ \substack {ll \\ w r p w } } \delta_{s t} 
+        \frac83 g_1^2 N_c \hyp_e \hyp_q C^{(1)}_{ \substack {lq \\ p r w w } } \delta_{s t} 
+        \frac83 g_1^2 N_c \hyp_l \hyp_q L_{ \substack {qe \\ w w s t } } \delta_{p r}
+        \frac43 g_1^2 \hyp_e^2 L_{ \substack {le \\ p r w w } } \delta_{s t} 
+        \frac83 g_1^2 \hyp_l^2 L_{ \substack {le \\ w w s t } } \delta_{p r} 
\nn
&
+        \frac43 g_1^2 N_c \hyp_e \hyp_u L_{ \substack {lu \\ p r w w } } \delta_{s t}
+        \frac43 g_1^2 N_c \hyp_d \hyp_e L_{ \substack {ld \\ p r w w } } \delta_{s t} 
+        \frac43 g_1^2 N_c \hyp_l \hyp_u L_{ \substack {eu \\ s t w w } } \delta_{p r} 
+        \frac43 g_1^2 N_c \hyp_d \hyp_l L_{ \substack {ed \\ s t w w } } \delta_{p r} 
\nn
&
+        \frac43 g_1^2 \hyp_e \hyp_l L_{ \substack {ee \\ s t w w } } \delta_{p r} 
+        \frac43 g_1^2 \hyp_e \hyp_l L_{ \substack {ee \\ s w w t } } \delta_{p r} 
+        \frac43 g_1^2 \hyp_e \hyp_l L_{ \substack {ee \\ w t s w } } \delta_{p r}
+        \frac43 g_1^2 \hyp_e \hyp_l L_{ \substack {ee \\ w w s t } } \delta_{p r} 
-12 \hyp_l \hyp_e g_1^2 L_{\substack{ le \\ prst }} 
\end{align}
For the dimension-seven SMEFT operators, we cite here the results from \parencite{liao2019renormalization} - the authors present their results directly in terms of the counterterms that each operator requires at one loop, so we collect the respective contributions from all SMEFT operators.
\begin{align}
    \begin{autobreak}
    \dot L_{\bar{e}LLLH}=
    % eLLLH
    +\frac{1}{8}\Big\{\big(9g_1^2+7g_2^2-4W_H\big)C^{prst}_{\bar{e}LLLH}
    -8g_2^2\big(C^{prts}_{\bar{e}LLLH}-C^{pstr}_{\bar{e}LLLH}+2C^{ptsr}_{\bar{e}LLLH}\big)
    -4(Y^\dagger_e)_{pr}(Y_d)_{vw}\big(2C^{wvst}_{\bar{e}LLLH}+C^{wsvt}_{\bar{e}LLLH}
    +C^{wstv}_{\bar{e}LLLH}-C^{wtsv}_{\bar{e}LLLH}\big)
    -2\Big[6(Y^\dagger_e Y_e)_{pv}C^{vrst}_{\bar{e}LLLH}
    +(Y_e Y_e^\dagger)_{vr}\big(C^{pvst}_{\bar{e}LLLH}+4C^{pvts}_{\bar{e}LLLH}\big)+5(Y_e Y_e^\dagger)_{vs}C^{prvt}_{\bar{e}LLLH}
    -(Y_e Y_e^\dagger)_{vt}\big(4C^{prvs}_{\bar{e}LLLH}+3C^{prsv}_{\bar{e}LLLH}\big)
    \Big]\Big\}
    % LHD1
    -\frac{1}{4}\Big\{3g_1^2(Y^\dagger_e)_{pt}\big(C^{rs}_{LHD1}+C^{sr}_{LHD1}\big)
    -g_2^2\Big[(Y^\dagger_e)_{pt}\big(C^{rs}_{LHD1}-C^{sr}_{LHD1} \big)-2(Y^\dagger_e)_{ps}\big(C^{rt}_{LHD1}-C^{tr}_{LHD1} \big)\Big]
    -4(Y^\dagger_e)_{pr}W^1_{ts}\Big\}
    % LHD2
    -\frac{1}{4}\Big\{3\big(g_1^2-g_2^2\big)\Big[(Y^\dagger_e)_{pt}C^{rs}_{LHD2}
    -(Y^\dagger_e)_{ps}\big(C^{rt}_{LHD2}-C^{tr}_{LHD2}\big)\Big]
    -4(Y^\dagger_e)_{pr}W^2_{ts}\Big\}
    % LHB
    -3g^2_1\Big[(Y_e^\dagger)_{pr}C^{st}_{LHB}
    +(Y_e^\dagger)_{pt}C^{rs}_{LHB}-4(Y_e^\dagger)_{ps}C^{rt}_{LHB}\Big]
    % LHW
    +\frac{1}{2}g^2_2
    \Big\{(Y^\dagger_e)_{pt}\big(5C^{rs}_{LHW}+C^{sr}_{LHW}\big)
    -4(Y^\dagger_e)_{ps}\big(C^{rt}_{LHW}-C^{tr}_{LHW}\big)\Big\}
    % dLQLH1
    -\frac{3}{2}(Y^\dagger_e)_{pr}(Y_d)_{vw}C^{wsvt}_{\bar{d}LQLH1}
    % dLQLH2
    -\frac{3}{2}(Y^\dagger_e)_{pr}(Y_d)_{vw}\big(C^{wsvt}_{\bar{d}LQLH2}-C^{wtvs}_{\bar{d}LQLH2}\big)L_{\bar{d}LQLH2}^{prst}
    % QuLLH
    +3(Y^\dagger_e)_{pr}(Y^\dagger_u)_{vw}C^{wvst}_{\bar{Q}uLLH}
    \end{autobreak}
\end{align}
\begin{align}
    \begin{autobreak}
    \dot L_{LeHD}^{pr}=
    % LeHD
    \frac{1}{4}\Big[\big(3g_1^2-12\lambda-6W_H\big)C^{pr}_{LeHD}
    -(Y_eY^\dagger_e)_{vp}C^{vr}_{LeHD}-2(Y_e)_{vr}(Y^\dagger_e)_{wp}C^{vw}_{LeHD}
    -8(Y^\dagger_eY_e)_{vr}C^{pv}_{LeHD}\Big]
    % LHD1
    -\frac{1}{4}\Big\{(Y_e)_{vr}\Big[\big(3g_1^2
    -4g_2^2\big)C^{pv}_{LHD1} +\big(3g_1^2+2g_2^2\big)C^{vp}_{LHD1} \Big]
    -\Big[(Y_eY^\dagger_e)_{vp}(Y_e)_{wr}\big(2C^{vw}_{LHD1}-C^{wv}_{LHD1}\big)
    +4(Y_eY^\dagger_eY_e)_{vr}C^{pv}_{LHD1}\Big]\Big\}
    % LHD2
    -\frac{1}{16}\Big[\big( 13g_1^2-17g_2^2-8\lambda\big)(Y_e)_{vr}C^{pv}_{LHD2}
    -4\Big(5(Y_eY^\dagger_eY_e)_{vr}C^{pv}_{LHD2}
    +(Y_eY^\dagger_e)_{vp}(Y_e)_{wr}C^{vw}_{LHD2}\Big)
    \Big]
    % LeHD
    +3(Y^\dagger_uY_d)_{vw}C^{wpvr}_{LeHD}
    \end{autobreak}
\end{align}


\chapter{UV Lagrangian}
\section{Irreps}
UV scalars, fermions and vectors that contribute to SMEFT operators up to dimension 7 at tree-level in tab. \ref{tab:BSMsfv} \cite{li2024complete}.
\begin{table}[]
    \centering
    %\renewcommand{\arraystretch}{1.5}
    \small
    \begin{tabular}{lcccccccc}
        \toprule
        Notation \cite{li2024complete}& \cellcolor{gray!25} $S_1$ & $S_2$ & $S_3$ & \cellcolor{gray!25}$S_4$ & $S_5$ & \cellcolor{gray!25}$S_6$ & $S_7$ & $S_8$ \\
        Notation \cite{de2018effective} & \cellcolor{gray!25}$\mathcal{S}$ & $\mathcal{S}_1$ & $\mathcal{S}_2$ & \cellcolor{gray!25}$\varphi$ & $\Xi$ & \cellcolor{gray!25}$\Xi_1$ & $\Theta_1$ & $\Theta_3$ \\
        Irrep & \cellcolor{gray!25}$({\bf 1,1})_0$ & $({\bf 1,1})_1$ & $({\bf 1,1})_2$ & \cellcolor{gray!25}$({\bf 1,2})_{\frac{1}{2}}$ & $({\bf 1,3})_0$ & \cellcolor{gray!25}$({\bf 1,3})_1$ & $({\bf 1,4})_{\frac{1}{2}}$ & $({\bf 1,4})_{\frac{3}{2}}$ \\
        \midrule
        Notation \cite{li2024complete}& $S_9$ & $S_{10}$ & $S_{11}$ & $S_{12}$ & $S_{13}$ & $S_{14}$ & & \\
        Notation \cite{de2018effective} & $\omega_4$ & $\omega_1$ & $\omega_2$ & $\Pi_1$ & $\Pi_7$ & $\zeta$ & & \\
        Irrep & $({\bf 3,1})_{-\frac{4}{3}}$ & $({\bf 3,1})_{-\frac{1}{3}}$ & $({\bf 3,1})_{\frac{2}{3}}$ & $({\bf 3,2})_{\frac{1}{6}}$ & $({\bf 3,2})_{\frac{7}{6}}$ & $({\bf 3,3})_{-\frac{1}{3}}$ & & \\
        \midrule
        Notation \cite{li2024complete}& $S_{15}$ & $S_{16}$ & $S_{17}$ & $S_{18}$ & $S_{19}$ & & & \\
        Notation \cite{de2018effective} & $\Omega_2$ & $\Omega_1$ & $\Omega_4$ & $\Upsilon_1$ & $\Phi$ & & & \\
        Irrep & $({\bf 6,1})_{-\frac{2}{3}}$ & $({\bf 6,1})_{\frac{1}{3}}$ & $({\bf 6,1})_{\frac{4}{3}}$ & $({\bf 6,3})_{\frac{1}{3}}$ & $({\bf 8,2})_{\frac{1}{2}}$ & & & \\
        \midrule
        Notation \cite{li2024complete}& \cellcolor{gray!25}$F_1$ & \cellcolor{gray!25}$F_2$ & \cellcolor{gray!25}$F_3$ & $F_4$ & \cellcolor{gray!25}$F_5$ & \cellcolor{gray!25}$F_6$ & $F_7$ \\
        Notation \cite{de2018effective} & \cellcolor{gray!25}$N$ & \cellcolor{gray!25}$E^c$ & \cellcolor{gray!25}$\Delta_1^c$ & $\Delta_3^c$ & \cellcolor{gray!25}$\Sigma$ & \cellcolor{gray!25}$\Sigma_1^c$ & \\
        Irrep & \cellcolor{gray!25}$({\bf 1,1})_0$ & \cellcolor{gray!25}$({\bf 1,1})_{1}$ & \cellcolor{gray!25}$({\bf 1,2})_{\frac{1}{2}}$ & $({\bf 1,2})_{\frac{3}{2}}$ & \cellcolor{gray!25}$({\bf 1,3})_0$ & \cellcolor{gray!25}$({\bf 1,3})_{1}$ & $({\bf 1,4})_{\frac{1}{2}}$ \\
        \midrule 
        Notation \cite{li2024complete} & $F_8$ & $F_{9}$ & $F_{10}$ & $F_{11}$ & $F_{12}$ & $F_{13}$ & $F_{14}$ \\
        Notation \cite{de2018effective} & $D$ & $U$ & $Q_5$ & $Q_1$ & $Q_7$ & $T_1$ & $T_2$ \\
        Irrep & $({\bf 3,1})_{-\frac{1}{3}}$ & $({\bf 3,1})_{\frac{2}{3}}$ & $({\bf 3,2})_{-\frac{5}{6}}$ & $({\bf 3,2})_{\frac{1}{6}}$ & $({\bf 3,2})_{\frac{7}{6}}$ & $({\bf 3,3})_{-\frac{1}{3}}$ & $({\bf 3,3})_{\frac{2}{3}}$ \\
        \midrule
        Notation \cite{li2024complete}& \cellcolor{gray!25}$V_1$ & \cellcolor{gray!25}$V_2$ & \cellcolor{gray!25}$V_3$ & \cellcolor{gray!25}$V_4$ & $V_5$ & $V_6$ & $V_7$ \\
        Notation \cite{de2018effective} & \cellcolor{gray!25}$\mathcal{B}$ & \cellcolor{gray!25}$\mathcal{B}_1$ & \cellcolor{gray!25}$\mathcal{L}_3^\dagger$ & \cellcolor{gray!25}$\mathcal{W}$ & $\mathcal{U}_2$ & $\mathcal{U}_5$ & $\mathcal{Q}_5$ \\
        Irrep & \cellcolor{gray!25}$({\bf 1,1})_0$ & \cellcolor{gray!25}$({\bf 1,1})_1$ & \cellcolor{gray!25}$({\bf 1,2})_{\frac{3}{2}}$ & \cellcolor{gray!25}$({\bf 1,3})_0$ & $({\bf 3,1})_{\frac{2}{3}}$ & $({\bf 3,1})_{\frac{5}{3}}$ &  $({\bf 3,2})_{-\frac{5}{6}}$ \\
        \midrule
        Notation \cite{li2024complete}& $V_8$ & $V_{9}$ & $V_{10}$ & $V_{11}$ & $V_{12}$ & $V_{13}$ & $V_{14}$ \\
        Notation \cite{de2018effective}& $\mathcal{Q}_1$ & $\mathcal{X}$ & $\mathcal{Y}_1^\dagger$ & $\mathcal{Y}_5^\dagger$ & $\mathcal{G}$ & $\mathcal{G}_1$ & $\mathcal{H}$ \\
        Irrep & $({\bf 3,2})_{\frac{1}{6}}$ &$({\bf 3,3})_{\frac{2}{3}}$ & $({\bf {6}, 2})_{-\frac{1}{6}}$ & $({\bf {6}, 2})_{\frac{5}{6}}$ & $({\bf 8,1})_0$ & $({\bf 8,1})_1$ & $({\bf 8,3})_0$ \\
        \bottomrule
    \end{tabular}
    \caption{BSM scalars, fermions and vectors that can contribute to operators up to dimension 7 in SMEFT at tree level. Table(s) from \cite{li2024complete}, name in the second row corresponding to \cite{de2018effective}. Fermion superscript ${}^c$ if they are conjugated with each other, Majorana fermions $F_1$ and $F_5$ have parity left, i.e. $F = F_L = P_L F$. Vector superscript ${}^\dagger$ if they are conjugated with each other. Fields that produce the operators $\mathcal{O}_{ll}, \mathcal{O}_{le}, \mathcal{O}^{(1)}_{Hl}, \mathcal{O}^{(3)}_{Hl}, \mathcal{O}_{LeHD}, \mathcal{O}_{\bar eLLLH}$ relevant for muon decay are highlighted in gray.}
    \label{tab:BSMsfv}
\end{table}
\section{UV Scalars}
\begin{align}
    \begin{autobreak}
        \Delta{\cal L}_{\rm S_2,S_4,S_6}^{(d = 3)} =  
 + \frac{\mathcal{C}_{{HHS_{6}^\dagger }}^{{p}}}{\sqrt{2}}\epsilon ^{{kj}} (\tau ^{{I}})_{{k}}^{{i}}{H}_{{i}}{H}_{{j}}({S}_{{6p}}^{\dagger })^{{I}} 
 + \mathcal{C}_{{HS_{2}^\dagger S_{4}}}^{{rp}}\epsilon ^{{ji}}{H}_{{j}}({S}_{{4r}})_{{i}}{S}_{{2p}}^{\dagger } 
 + \frac{\mathcal{C}_{{HS_{4}S_{6}^\dagger }}^{{pr}}}{\sqrt{2}}\epsilon ^{{ki}} (\tau ^{{I}})_{{k}}^{{j}}{H}_{{j}}({S}_{{4p}})_{{i}}({S}_{{6r}}^{\dagger })^{{I}} 
    \end{autobreak}
\end{align}
\begin{align}
\begin{autobreak}
\Delta{\cal L}_{\rm S_2,S_4,S_6}^{(d = 4)} = 
        - \mathcal{D}_{{LLS_{2}}}^{{rsp}}\epsilon ^{{ij}}[({l}_{{r}})_{{i}}{C}({l}_{{s}})_{{j}}]{S}_{{2p}} 
 - \mathcal{D}_{{e^\dagger LS_{4}^\dagger }}^{{rsp}} [\bar{e}_{{r}}({l}_{{s}})_{{i}}]({S}_{{4p}}^{\dagger })^{{i}} 
 - 2 \mathcal{D}_{{QS_{4}u^\dagger }}^{{rps}}\epsilon ^{{ji}} [(\bar{u}_{{s}})^{{a}}({q}_{{r}})_{{aj}}]({S}_{{4p}})_{{i}} 
 - \frac{\mathcal{D}_{{LLS_{6}}}^{{rsp}}}{\sqrt{2}}\epsilon ^{{jk}} (\tau ^{{I}})_{{k}}^{{i}}[({l}_{{r}})_{{i}}{C}({l}_{{s}})_{{j}}]({S}_{{6p}})^{{I}} 
 - \mathcal{D}_{{HH^\dagger H^\dagger S_{4}(1)}}^{{p}}\delta _{{k}}^{{i}} \delta _{{l}}^{{j}}{H}_{{j}}{H}^{{\dagger k}}{H}^{{\dagger l}}({S}_{{4p}})_{{i}} 
 + \frac{\mathcal{D}_{{HH^\dagger S_{2}S_{6}^\dagger }}^{{pr}}}{\sqrt{2}}(\tau ^{{I}})_{{j}}^{{i}}{H}_{{i}}{H}^{{\dagger j}}({S}_{{6r}}^{\dagger })^{{I}}{S}_{{2p}} 
 + \frac{\mathcal{D}_{{HH^\dagger S_{6}S_{6}^\dagger (1)}}^{{pr}}}{2}(\tau ^{{I}})_{{k}}^{{i}} (\tau ^{{J}})_{{j}}^{{k}}{H}_{{i}}{H}^{{\dagger j}}({S}_{{6p}})^{{I}}({S}_{{6r}}^{\dagger })^{{J}} 
 + \mathcal{D}_{{HH^\dagger S_{6}S_{6}^\dagger (2)}}^{{pr}}\delta ^{{IJ}} \delta _{{j}}^{{i}}{H}_{{i}}{H}^{{\dagger j}}({S}_{{6p}})^{{I}}({S}_{{6r}}^{\dagger })^{{J}} 
    \end{autobreak}
\end{align}
\subsection{New UV fermions}

\subsection{New UV vectors}
\subsection{New UV mixing terms}
\subsection{SMEFT to UV matching conditions}
\begin{align}
\begin{autobreak}
 {C_{ll}} =
 
-\frac{\mathcal{D}_{{LLS_{2}}}^{{f_1f_2p_1}} \mathcal{D}_{{LLS_{2}}}^{{f_4f_3p_1*}}}{2 {M_{S_{2}}^2}}
+\frac{\mathcal{D}_{{LLS_{6}}}^{{f_1f_2p_1}} \mathcal{D}_{{LLS_{6}}}^{{f_4f_3p_1*}}}{4 {M_{S_{6}}^2}}
-\frac{\mathcal{D}_{{LL^\dagger V_{1}}}^{{f_1f_3p_1}} \mathcal{D}_{{LL^\dagger V_{1}}}^{{f_2f_4p_1}}}{2 {M_{V_{1}}^2}}
+\frac{\mathcal{D}_{{LL^\dagger V_{4}}}^{{f_1f_3p_1}} \mathcal{D}_{{LL^\dagger V_{4}}}^{{f_2f_4p_1}}}{4 {M_{V_{4}}^2}}
-\frac{\mathcal{D}_{{LL^\dagger V_{4}}}^{{f_1f_4p_1}} \mathcal{D}_{{LL^\dagger V_{4}}}^{{f_2f_3p_1}}}{2 {M_{V_{4}}^2}}
+\frac{\mathcal{D}_{{LLS_{2}}}^{{f_2f_1p_1}} \mathcal{D}_{{LLS_{2}}}^{{f_4f_3p_1*}}}{2 {M_{S_{2}}^2}}
+\frac{\mathcal{D}_{{LLS_{6}}}^{{f_2f_1p_1}} \mathcal{D}_{{LLS_{6}}}^{{f_4f_3p_1*}}}{4 {M_{S_{6}}^2}} 
\end{autobreak}\label{eq:ll-UV}
\end{align}
\begin{align}
\begin{autobreak}
 {C_{le}} =
 \frac{\mathcal{D}_{{e^\dagger eV_{1}}}^{{f_1f_3p_1}} \mathcal{D}_{{LL^\dagger V_{1}}}^{{f_2f_4p_1}}}{{M_{V_{1}}^2}}
-\frac{\mathcal{D}_{{e^\dagger LS_{4}^\dagger }}^{{f_1f_2p_1}} \mathcal{D}_{{e^\dagger LS_{4}^\dagger }}^{{f_3f_4p_1*}}}{2 {M_{S_{4}}^2}}
+\frac{\mathcal{D}_{{e^\dagger L^\dagger V_{3}^\dagger }}^{{f_1f_4p_1}} \mathcal{D}_{{e^\dagger L^\dagger V_{3}^\dagger }}^{{f_3f_2p_1*}}}{{M_{V_{3}}^2}} 
\end{autobreak}
\end{align}
\begin{align}
\begin{autobreak}
 {C_{Hl}^{(1)}} =
 
-\frac{i \mathcal{D}_{{LL^\dagger V_{1}}}^{{f_1f_4p_1}} \mathcal{D}_{{HH^\dagger V_{1}D}}^{{p_1}}}{{M_{V_{1}}^2}}
+\frac{i \mathcal{D}_{{LL^\dagger V_{1}}}^{{f_1f_4p_1}} \mathcal{D}_{{HH^\dagger V_{1}D}}^{{p_1*}}}{{M_{V_{1}}^2}}
+\frac{\mathcal{D}_{{F_{1}HL}}^{{p_1f_1}} \mathcal{D}_{{F_{1}HL}}^{{p_1f_4*}}}{4 {M_{F_{1}}^2}}
-\frac{\mathcal{D}_{{F_{2L}H^\dagger L}}^{{p_1f_1}} \mathcal{D}_{{F_{2L}H^\dagger L}}^{{p_1f_4*}}}{4 {M_{F_{2}}^2}}
+\frac{3 \mathcal{D}_{{F_{5}HL}}^{{p_1f_1}} \mathcal{D}_{{F_{5}HL}}^{{p_1f_4*}}}{8 {M_{F_{5}}^2}}
-\frac{3 \mathcal{D}_{{F_{6L}H^\dagger L}}^{{p_1f_1}} \mathcal{D}_{{F_{6L}H^\dagger L}}^{{p_1f_4*}}}{8 {M_{F_{6}}^2}} 
\end{autobreak}
\end{align}
\begin{align}
\begin{autobreak}
 {C_{Hl}^{(3)}} =
 
-\frac{i \mathcal{D}_{{LL^\dagger V_{4}}}^{{f_1f_4p_1}} \mathcal{D}_{{HH^\dagger V_{4}D}}^{{p_1}}}{2 {M_{V_{4}}^2}}
-\frac{i \mathcal{D}_{{LL^\dagger V_{4}}}^{{f_1f_4p_1}} \mathcal{D}_{{HH^\dagger V_{4}D}}^{{p_1*}}}{2 {M_{V_{4}}^2}}
-\frac{\mathcal{D}_{{F_{1}HL}}^{{p_1f_1}} \mathcal{D}_{{F_{1}HL}}^{{p_1f_4*}}}{4 {M_{F_{1}}^2}}
-\frac{\mathcal{D}_{{F_{2L}H^\dagger L}}^{{p_1f_1}} \mathcal{D}_{{F_{2L}H^\dagger L}}^{{p_1f_4*}}}{4 {M_{F_{2}}^2}}
+\frac{\mathcal{D}_{{F_{5}HL}}^{{p_1f_1}} \mathcal{D}_{{F_{5}HL}}^{{p_1f_4*}}}{8 {M_{F_{5}}^2}}
+\frac{\mathcal{D}_{{F_{6L}H^\dagger L}}^{{p_1f_1}} \mathcal{D}_{{F_{6L}H^\dagger L}}^{{p_1f_4*}}}{8 {M_{F_{6}}^2}} 
\end{autobreak}
\end{align}
\begin{align}
\begin{autobreak}
 {C_{He}} =
 \frac{i \mathcal{D}_{{e^\dagger eV_{1}}}^{{f_1f_4p_1}} \mathcal{D}_{{HH^\dagger V_{1}D}}^{{p_1}}}{{M_{V_{1}}^2}}
-\frac{i \mathcal{D}_{{e^\dagger eV_{1}}}^{{f_1f_4p_1}} \mathcal{D}_{{HH^\dagger V_{1}D}}^{{p_1*}}}{{M_{V_{1}}^2}}
+\frac{\mathcal{D}_{{e^\dagger F_{3R}^\dagger H^\dagger }}^{{f_1p_1}} \mathcal{D}_{{e^\dagger F_{3R}^\dagger H^\dagger }}^{{f_4p_1*}}}{2 {M_{F_{3}}^2}}
-\frac{\mathcal{D}_{{e^\dagger F_{4R}^\dagger H}}^{{f_1p_1}} \mathcal{D}_{{e^\dagger F_{4R}^\dagger H}}^{{f_4p_1*}}}{2 {M_{F_{4}}^2}} 
\end{autobreak}
\end{align}
\begin{align}
\begin{autobreak}
 {C_{LeHD}} =
-\frac{{y}_{{e}}^{{p_1f_5}} \mathcal{D}_{{F_{1}HL}}^{{p_2f_1}} \mathcal{D}_{{F_{1}HL}}^{{p_2p_1}}}{2 {M_{F_{1}}^3}}
-\frac{{y}_{{e}}^{{p_1f_5}} \mathcal{D}_{{LLS_{6}}}^{{f_1p_1p_2}} \mathcal{C}_{{HHS_{6}^\dagger }}^{{p_2}}}{{M_{S_{6}}^4}}
-\frac{{y}_{{e}}^{{p_1f_5}} \mathcal{D}_{{F_{5}HL}}^{{p_2f_1}} \mathcal{D}_{{F_{5}HL}}^{{p_2p_1}}}{4 {M_{F_{5}}^3}}
+\frac{{y}_{{e}}^{{p_1f_5}} \mathcal{C}_{{HHS_{6}^\dagger }}^{{p_2}} \mathcal{D}_{{LLS_{6}}}^{{p_1f_1p_2}}}{{M_{S_{6}}^4}}
+\frac{2 i \mathcal{D}_{{eF_{1}V_{2}}}^{{f_5p_1p_2}} \mathcal{D}_{{F_{1}HL}}^{{p_1f_1}} \mathcal{D}_{{HHV_{2}^\dagger D}}^{{p_2}}}{{M_{F_{1}}} {M_{V_{2}}^2}}
-\frac{\mathcal{D}_{{e^\dagger F_{3R}^\dagger H^\dagger }}^{{f_5p_1*}} \mathcal{D}_{{F_{1}F_{3R}^\dagger H}}^{{p_2p_1}} \mathcal{D}_{{F_{1}HL}}^{{p_2f_1}}}{{M_{F_{1}}} {M_{F_{3}}^2}}
-\frac{\mathcal{D}_{{e^\dagger F_{3R}^\dagger H^\dagger }}^{{f_5p_1*}} \mathcal{D}_{{F_{3L}F_{5}H^\dagger }}^{{p_1p_2*}} \mathcal{D}_{{F_{5}HL}}^{{p_2f_1}}}{{M_{F_{3}}} {M_{F_{5}}^2}}
+\frac{2 i \mathcal{D}_{{e^\dagger F_{3R}^\dagger H^\dagger }}^{{f_5p_1*}} \mathcal{D}_{{F_{3L}L^\dagger V_{2}^\dagger }}^{{p_1f_1p_2*}} \mathcal{D}_{{HHV_{2}^\dagger D}}^{{p_2}}}{{M_{F_{3}}} {M_{V_{2}}^2}}
-\frac{\mathcal{D}_{{e^\dagger F_{3R}^\dagger H^\dagger }}^{{f_5p_1*}} \mathcal{D}_{{F_{3R}^\dagger F_{5}H}}^{{p_1p_2}} \mathcal{D}_{{F_{5}HL}}^{{p_2f_1}}}{2 {M_{F_{3}}^2} {M_{F_{5}}}}
-\frac{\mathcal{D}_{{e^\dagger F_{3R}^\dagger H^\dagger }}^{{f_5p_1*}} \mathcal{D}_{{F_{3R}^\dagger LS_{6}}}^{{p_1f_1p_2}} \mathcal{C}_{{HHS_{6}^\dagger }}^{{p_2}}}{{M_{F_{3}}^2} {M_{S_{6}}^2}}
+\frac{\mathcal{D}_{{e^\dagger F_{5}S_{6}^\dagger }}^{{f_5p_1p_2*}} \mathcal{D}_{{F_{5}HL}}^{{p_1f_1}} \mathcal{C}_{{HHS_{6}^\dagger }}^{{p_2}}}{{M_{F_{5}}^2} {M_{S_{6}}^2}}
-\frac{2 i \mathcal{D}_{{e^\dagger L^\dagger V_{3}^\dagger }}^{{f_5f_1p_1*}} \mathcal{C}_{{HHS_{6}^\dagger }}^{{p_2}} \mathcal{D}_{{HS_{6}V_{3}^\dagger D}}^{{p_2p_1}}}{{M_{S_{6}}^2} {M_{V_{3}}^2}}
-\frac{4 i \mathcal{D}_{{e^\dagger L^\dagger V_{3}^\dagger }}^{{f_5f_1p_1*}} \mathcal{D}_{{HHV_{2}^\dagger D}}^{{p_2}} \mathcal{C}_{{HV_{2}V_{3}^\dagger }}^{{p_2p_1}}}{{M_{V_{2}}^2} {M_{V_{3}}^2}} 
\end{autobreak}\label{eq:LeHD-UV}
\end{align}
\begin{align}
\begin{autobreak}
 {C_{eLLLH}} =
-\frac{{y}_{{e}}^{{f_1f_2*}} \mathcal{D}_{{F_{1}HL}}^{{p_1f_3}} \mathcal{D}_{{F_{1}HL}}^{{p_1f_4}}}{2 {M_{F_{1}}^3}}
-\frac{3 \mathcal{D}_{{e^\dagger F_{1}S_{2}^\dagger }}^{{f_1p_1p_2}} \mathcal{D}_{{LLS_{2}}}^{{f_2f_3p_2}} \mathcal{D}_{{F_{1}HL}}^{{p_1f_4}}}{2 {M_{F_{1}}} {M_{S_{2}}^2}}
+\frac{3 \mathcal{D}_{{e^\dagger F_{1}S_{2}^\dagger }}^{{f_1p_1p_2}} \mathcal{D}_{{LLS_{2}}}^{{f_3f_2p_2}} \mathcal{D}_{{F_{1}HL}}^{{p_1f_4}}}{2 {M_{F_{1}}} {M_{S_{2}}^2}}
-\frac{\mathcal{D}_{{e^\dagger LS_{4}^\dagger }}^{{f_1f_2p_1}} \mathcal{D}_{{F_{1}HL}}^{{p_2f_4}} \mathcal{D}_{{F_{1L}S_{4}}}^{{p_2f_3p_1}}}{{M_{F_{1}}} {M_{S_{4}}^2}}
-\frac{{y}_{{e}}^{{f_1f_2*}} \mathcal{D}_{{F_{5}HL}}^{{p_1f_3}} \mathcal{D}_{{F_{5}HL}}^{{p_1f_4}}}{4 {M_{F_{5}}^3}}
-\frac{\mathcal{D}_{{e^\dagger LS_{4}^\dagger }}^{{f_1f_4p_1}} \mathcal{D}_{{F_{5}HL}}^{{p_2f_3}} \mathcal{D}_{{F_{5L}S_{4}}}^{{p_2f_2p_1}}}{{M_{F_{5}}} {M_{S_{4}}^2}}
+\frac{\mathcal{D}_{{e^\dagger LS_{4}^\dagger }}^{{f_1f_3p_1}} \mathcal{D}_{{F_{5}HL}}^{{p_2f_4}} \mathcal{D}_{{F_{5L}S_{4}}}^{{p_2f_2p_1}}}{2 {M_{F_{5}}} {M_{S_{4}}^2}}
+\frac{\mathcal{D}_{{e^\dagger LS_{4}^\dagger }}^{{f_1f_4p_1}} \mathcal{D}_{{F_{5}HL}}^{{p_2f_2}} \mathcal{D}_{{F_{5L}S_{4}}}^{{p_2f_3p_1}}}{2 {M_{F_{5}}} {M_{S_{4}}^2}}
+\frac{\mathcal{D}_{{e^\dagger LS_{4}^\dagger }}^{{f_1f_3p_1}} \mathcal{D}_{{F_{5}HL}}^{{p_2f_2}} \mathcal{D}_{{F_{5L}S_{4}}}^{{p_2f_4p_1}}}{2 {M_{F_{5}}} {M_{S_{4}}^2}}
+\frac{\mathcal{D}_{{e^\dagger F_{4R}^\dagger H}}^{{f_1p_1}} \mathcal{D}_{{F_{4L}LS_{2}^\dagger }}^{{p_1f_4p_2}} \mathcal{D}_{{LLS_{2}}}^{{f_2f_3p_2}}}{{M_{F_{4}}} {M_{S_{2}}^2}}
+\frac{2 \mathcal{D}_{{e^\dagger LS_{4}^\dagger }}^{{f_1f_4p_1}} \mathcal{C}_{{HS_{2}^\dagger S_{4}}}^{{p_2p_1}} \mathcal{D}_{{LLS_{2}}}^{{f_2f_3p_2}}}{{M_{S_{2}}^2} {M_{S_{4}}^2}}
-\frac{\mathcal{D}_{{e^\dagger F_{4R}^\dagger H}}^{{f_1p_1}} \mathcal{D}_{{F_{4L}LS_{2}^\dagger }}^{{p_1f_4p_2}} \mathcal{D}_{{LLS_{2}}}^{{f_3f_2p_2}}}{{M_{F_{4}}} {M_{S_{2}}^2}}
-\frac{2 \mathcal{D}_{{e^\dagger LS_{4}^\dagger }}^{{f_1f_4p_1}} \mathcal{C}_{{HS_{2}^\dagger S_{4}}}^{{p_2p_1}} \mathcal{D}_{{LLS_{2}}}^{{f_3f_2p_2}}}{{M_{S_{2}}^2} {M_{S_{4}}^2}}
+\frac{\mathcal{D}_{{e^\dagger F_{1}S_{2}^\dagger }}^{{f_1p_1p_2}} \mathcal{D}_{{F_{1}HL}}^{{p_1f_2}} \mathcal{D}_{{LLS_{2}}}^{{f_3f_4p_2}}}{2 {M_{F_{1}}} {M_{S_{2}}^2}}
-\frac{\mathcal{D}_{{e^\dagger LS_{4}^\dagger }}^{{f_1f_2p_1}} \mathcal{C}_{{HS_{2}^\dagger S_{4}}}^{{p_2p_1}} \mathcal{D}_{{LLS_{2}}}^{{f_3f_4p_2}}}{{M_{S_{2}}^2} {M_{S_{4}}^2}}
-\frac{\mathcal{D}_{{e^\dagger F_{1}S_{2}^\dagger }}^{{f_1p_1p_2}} \mathcal{D}_{{F_{1}HL}}^{{p_1f_2}} \mathcal{D}_{{LLS_{2}}}^{{f_4f_3p_2}}}{2 {M_{F_{1}}} {M_{S_{2}}^2}}
+\frac{\mathcal{D}_{{e^\dagger LS_{4}^\dagger }}^{{f_1f_2p_1}} \mathcal{C}_{{HS_{2}^\dagger S_{4}}}^{{p_2p_1}} \mathcal{D}_{{LLS_{2}}}^{{f_4f_3p_2}}}{{M_{S_{2}}^2} {M_{S_{4}}^2}}
-\frac{\mathcal{D}_{{e^\dagger F_{4R}^\dagger H}}^{{f_1p_1}} \mathcal{D}_{{F_{4L}LS_{6}^\dagger }}^{{p_1f_4p_2}} \mathcal{D}_{{LLS_{6}}}^{{f_2f_3p_2}}}{2 {M_{F_{4}}} {M_{S_{6}}^2}}
-\frac{3 \mathcal{D}_{{e^\dagger F_{5}S_{6}^\dagger }}^{{f_1p_1p_2}} \mathcal{D}_{{F_{5}HL}}^{{p_1f_4}} \mathcal{D}_{{LLS_{6}}}^{{f_2f_3p_2}}}{4 {M_{F_{5}}} {M_{S_{6}}^2}}
-\frac{\mathcal{D}_{{e^\dagger LS_{4}^\dagger }}^{{f_1f_4p_1}} \mathcal{C}_{{HS_{4}S_{6}^\dagger }}^{{p_1p_2}} \mathcal{D}_{{LLS_{6}}}^{{f_2f_3p_2}}}{2 {M_{S_{4}}^2} {M_{S_{6}}^2}}
+\frac{\mathcal{D}_{{e^\dagger F_{4R}^\dagger H}}^{{f_1p_1}} \mathcal{D}_{{F_{4L}LS_{6}^\dagger }}^{{p_1f_4p_2}} \mathcal{D}_{{LLS_{6}}}^{{f_3f_2p_2}}}{2 {M_{F_{4}}} {M_{S_{6}}^2}}
+\frac{3 \mathcal{D}_{{e^\dagger F_{5}S_{6}^\dagger }}^{{f_1p_1p_2}} \mathcal{D}_{{F_{5}HL}}^{{p_1f_4}} \mathcal{D}_{{LLS_{6}}}^{{f_3f_2p_2}}}{4 {M_{F_{5}}} {M_{S_{6}}^2}}
+\frac{\mathcal{D}_{{e^\dagger LS_{4}^\dagger }}^{{f_1f_4p_1}} \mathcal{C}_{{HS_{4}S_{6}^\dagger }}^{{p_1p_2}} \mathcal{D}_{{LLS_{6}}}^{{f_3f_2p_2}}}{2 {M_{S_{4}}^2} {M_{S_{6}}^2}}
+\frac{\mathcal{D}_{{e^\dagger F_{5}S_{6}^\dagger }}^{{f_1p_1p_2}} \mathcal{D}_{{F_{5}HL}}^{{p_1f_2}} \mathcal{D}_{{LLS_{6}}}^{{f_3f_4p_2}}}{4 {M_{F_{5}}} {M_{S_{6}}^2}}
-\frac{\mathcal{D}_{{e^\dagger F_{4R}^\dagger H}}^{{f_1p_1}} \mathcal{D}_{{F_{4L}LS_{6}^\dagger }}^{{p_1f_2p_2}} \mathcal{D}_{{LLS_{6}}}^{{f_4f_3p_2}}}{{M_{F_{4}}} {M_{S_{6}}^2}}
-\frac{5 \mathcal{D}_{{e^\dagger F_{5}S_{6}^\dagger }}^{{f_1p_1p_2}} \mathcal{D}_{{F_{5}HL}}^{{p_1f_2}} \mathcal{D}_{{LLS_{6}}}^{{f_4f_3p_2}}}{4 {M_{F_{5}}} {M_{S_{6}}^2}}
-\frac{\mathcal{D}_{{e^\dagger LS_{4}^\dagger }}^{{f_1f_2p_1}} \mathcal{C}_{{HS_{4}S_{6}^\dagger }}^{{p_1p_2}} \mathcal{D}_{{LLS_{6}}}^{{f_4f_3p_2}}}{{M_{S_{4}}^2} {M_{S_{6}}^2}} 
\end{autobreak}\label{eq:eLLLH-UV}
\end{align}
\printbibliography
\end{fmffile}
\end{document}
